\chapter{Intrinsic geometric and automorphic comparison}
\label{v4-arith:ic:chapter}

A local Hecke polynomial determines two eigenvalues. It does not
determine an automorphic representation, the extensions between
representations, or the category on which the Hecke operators act.
For example, the functor
\[
 D(k[t])\longrightarrow D(k[t])\times D(k[t]),
 \qquad M\longmapsto(M,0),
\]
is fully faithful and respects a faithful scalar action on the target.
Its image misses every nonzero object in the second factor.
This elementary example separates the scalar comparison from
categorical reconstruction.

The arithmetic constructions of the preceding chapters provide
operators, coefficient algebras, marked interfaces and their
composition maps. An intrinsic realization must identify these data
with independently defined geometric or automorphic objects.
This chapter constructs the formal and local comparisons that enter
that problem. It retains the deformation groupoids, central actions,
extension classes and tensor transitions needed to pass between them.

The coefficients matter throughout. Smooth representations with
characteristic-zero coefficients away from the residue characteristic
have Bernstein blocks and Weil--Deligne parameters. Smooth
\(p\)-power-torsion representations have a different spectral recipient,
built from étale \((\varphi,\Gamma)\)-modules. Global determinant
deformations use one fixed coefficient ring at every place.
Each comparison below specifies the regime in which its maps are
defined.

\section{Marked objects and categorical recognition}
\label{v4-arith:ic:recognition}

Let \(\mathscr D\) be an \(\infty\)-category of arithmetic data.
An object may contain a coefficient algebra, selected modules,
Hecke or Galois actions, intertwiners, a distinguished state and
a scalar evaluation on a specified domain. Morphisms preserve
these markings. Let \(\mathscr R\) be a category of proposed geometric
realizations and let \(U:\mathscr R\to\mathscr D\) retain the data.
Write \(\mathscr D^\simeq\) for its maximal \(\infty\)-groupoid,
and similarly for \(\mathscr R\).

\begin{definition}\label{v4-arith:ic:marked-realizations}
For \(d\in\mathscr D\), the space of marked realizations is
\[
 \operatorname{Real}_U(d)=
 \operatorname{hofib}_d
 \bigl(\mathscr R^\simeq\xrightarrow{U^\simeq}
       \mathscr D^\simeq\bigr).
\]
A point consists of \(r\in\mathscr R\) and an equivalence
\(U(r)\simeq d\), with their higher identifications.
\end{definition}

Nonemptiness expresses existence. Contractibility expresses uniqueness
with coherent identifications. A functorial realization requires a
functor \(S:\mathscr D\to\mathscr R\) and an equivalence
\(US\simeq\operatorname{id}_{\mathscr D}\).
Objectwise choices do not define \(S\) on noninvertible morphisms.
When the desired comparison is a map, its realizations instead
retain that map in the appropriate arrow category.

This definition applies to several different arithmetic problems.
The relative singular models retain a chosen geometric background.
The marked fusion models also retain their associator, duality and
distinguished arithmetic junction. A global automorphic realization
must in addition preserve its ramified and archimedean actions,
coefficient topology and evaluation functional. A scalar identity
obtained after evaluation supplies one consequence of such a map.

\subsection{Coefficients and retained operations}

A coefficient change must be a map in the stated coefficient category.
For example,
\[
 (\mathbb Z/\ell^r)\otimes_{\mathbb Z}\mathbb Q=0,
 \qquad
 \left(\varprojlim_r\mathbb Z/\ell^r\right)
       \otimes_{\mathbb Z}\mathbb Q=\mathbb Q_\ell.
\]
Thus rationalization before completion destroys this coefficient
tower. Likewise, over \(\mathbb F_p\), the big Witt vector with only
\(a_p=1\) is nonzero although all its ghosts
\(\sum_{d\mid n}d\,a_d^{n/d}\) vanish.
The integral operations require their original coordinates.

Even an equivalence of underlying categories need not preserve a
chosen central action. Give \(D(k)\) the two \(k[t]\)-actions
\(t=0\) and \(t=1\). Its identity equivalence cannot intertwine them:
the proposed equality on the object \(k\) would be \(0=1\).
The action and its specified compatibility are therefore part of
the realization datum.

\subsection{Comparison on generators}

\begin{theorem}[Comparison on compact generators]
\label{v4-arith:ic:w5-e:local:generator-criterion}
Let $\mathcal C,\mathcal D$ be presentable stable
$\infty$-categories, and let $F:\mathcal C\to\mathcal D$ be exact
and colimit-preserving. Let $\{P_i\}_{i\in I}$ be a set of compact
generators of $\mathcal C$.
Here generation means that their smallest full stable subcategory
closed under all small colimits is $\mathcal C$.
Assume each $F(P_i)$ is compact and that the maps of mapping spectra
\[
 \operatorname{Map}_{\mathcal C}(P_i,P_j)
 \longrightarrow
 \operatorname{Map}_{\mathcal D}(F(P_i),F(P_j))
\]
are equivalences for all $i,j$. Then $F$ is fully faithful.
It is an equivalence if the objects $F(P_i)$ also generate
$\mathcal D$.
\end{theorem}

\begin{proof}
Fix $i$. Compactness and exactness imply that both mapping-spectrum
functors in the comparison, with second variable $X$, preserve
coproducts and finite colimits. They therefore preserve all colimits
between stable categories. The objects $X$ where the comparison is an
equivalence form a localizing subcategory of $\mathcal C$.
It contains every $P_j$, so it is all of $\mathcal C$.

Fix $X$ and vary the first argument. Mapping spectra convert colimits
in this argument into limits. Hence the objects $Y$ for which
\[
 \operatorname{Map}_{\mathcal C}(Y,X)
 \longrightarrow\operatorname{Map}_{\mathcal D}(F(Y),F(X))
\]
is an equivalence form a localizing subcategory containing all $P_i$.
This proves full faithfulness. Its essential image is closed under
colimits and finite cones, since diagrams in that image lift through
a fully faithful functor. If the image contains generators of
$\mathcal D$, it is all of $\mathcal D$.
\end{proof}

For an $E$-linear differential graded enhancement, the hypothesis is a
quasi-isomorphism on each derived Hom complex induced by the chosen
functor. An equality of dimensions or a match in degree one is weaker.
Use the augmentation modules $k=k[t]/(t)$ and
$k=(k[\epsilon]/(\epsilon^2))/(\epsilon)$. Then $k$ has
$\operatorname{Ext}^1(k,k)=k$ over both $k[t]$ and
$k[\epsilon]/(\epsilon^2)$. Over $k[t]$, its length-one free resolution
makes all higher Ext groups zero. Over the dual numbers, the periodic
resolution with differential multiplication by $\epsilon$ gives
$\operatorname{Ext}^n(k,k)=k$ for every $n\ge0$.

The local categorical Langlands question requires the actual source
subcategory, spectral support condition, comparison functor, and
objects $P_i$. Frobenius parameters and invariant vectors supply inputs
for this question. A categorical
equivalence follows only after the derived Hom and generation
hypotheses have been established on those categories.

\section{Determinants and local parameter spaces}
\label{v4-arith:ic:determinants}

A representation gives characteristic polynomials after every scalar
extension. A determinant polynomial law retains these compatible
polynomials even when no free representation has been chosen.
Its definition works in characteristic two, where a trace alone
does not contain every characteristic coefficient.


\begin{definition}[Chenevier determinant and associated trace]
\label{v4-arith:ic:w5-h:def:chenevier-det-trace}
Let $R$ be a commutative ring, $\Gamma$ a group and $d\geq1$.
An $R$-valued determinant of dimension $d$ is a multiplicative,
homogeneous polynomial law of degree $d$ from $R[\Gamma]$ to $R$.
Thus it specifies maps $D_A:A[\Gamma]\to A$ for every commutative
$R$-algebra $A$, natural in $A$, with
$D_A(1)=1$, $D_A(xy)=D_A(x)D_A(y)$ and
$D_A(ax)=a^dD_A(x)$. Write
\[
 D_{R[t]}(t-x)=\sum_{j=0}^d(-1)^j\Lambda_j(x)t^{d-j}.
\]

The associated trace is the $R$-linear map $T=\Lambda_1$. The
notation $(T,D)$ retains this dependence. The characteristic
polynomial gives the Cayley--Hamilton relations in the quotient
by their coefficients after polynomial base change. These relations
need not hold in $R[\Gamma]$ itself.
For a profinite $\Gamma$ and a topological coefficient ring,
continuity means that all characteristic-polynomial coefficient
functions on $\Gamma$ are continuous.
\end{definition}

\begin{proposition}[determinantal deformation rings]
\label{v4-arith:ic:w5-h:prop:chenevier-moduli-via-det-trace}

Let $k$ be a finite field of characteristic $p_0$, and put $W=W(k)$.
Fix a continuous residual determinant $\bar D$ of dimension $d$
of a profinite group $\Gamma$. On finite local $W$-algebras $A$
with residue field $k$, let $\mathcal D_\Gamma(A)$ be the set of
continuous determinants lifting $\bar D$. This covariant functor
is prorepresented by a profinite local $W$-algebra $R_\Gamma^{\mathrm{Ch}}$:
\[
 \mathcal D_\Gamma(A)=
 \operatorname{Hom}_{W,\mathrm{cont}}(R_\Gamma^{\mathrm{Ch}},A).
\]
Assume that every open subgroup has only finitely many continuous
homomorphisms to $\mathbb Z/p_0\mathbb Z$.

This is Mazur's finiteness condition, and it makes the ring complete
Noetherian. A closed deformation condition
represented by a quotient is imposed by that quotient.
The representing formal space is $\operatorname{Spf}R_\Gamma^{\mathrm{Ch}}$.
For $\Gamma=G_{\mathbb Q,S}$ with finite $S$, and for local Galois
groups in characteristic zero, the finiteness condition holds.
If the residual representation is absolutely irreducible and defined
over $k$, determinant deformations agree with its representation
deformations up to strict equivalence.
\end{proposition}

\begin{proof}[Attribution]
Chenevier, arXiv:0809.0415, \S1.1 and \S1.10 define polynomial
laws and their characteristic coefficients. Proposition~3.3,
p.~42, proves prorepresentability. Example~3.6 and
Proposition~3.7, p.~43, give the finiteness assertions.
Example~3.4 and Theorem~2.22 give the absolutely irreducible
comparison. Bella\"iche--Chenevier, \emph{Families of Galois
representations and Selmer groups}, Chapter~1, provides the related
deformation-theoretic setting. A chosen local condition must be
shown to define the quotient used in any particular application.
\end{proof}

\subsection{The Frobenius parameter stack}

Let $E$ be a field of characteristic zero and let $d\in E^\times$.
For an $E$-algebra $A$, consider a rank-two projective $A$-module $V$
with an automorphism $F$ satisfying
$\bigwedge^2F=d\,\operatorname{id}_{\bigwedge^2V}$.
Isomorphisms intertwine $F$. No trivialization of $\bigwedge^2V$ is
part of the data. This defines the quotient stack
\[
 \mathcal U_d=[Y_d/GL_2],\qquad
 Y_d=\operatorname{Spec}
 E[x_{11},x_{12},x_{21},x_{22}]/(x_{11}x_{22}-x_{12}x_{21}-d),
\]
where $GL_2$ acts by conjugation. It is the algebraic parameter stack
for an unramified rank-two Weil representation with chosen generator
acting by $F$.

Put $t=\operatorname{tr}F$. The regular semisimple open substack
$\mathcal U_d^{\mathrm{rs}}$ is defined by $t^2-4d\ne0$.
The ordered eigenvalue scheme and diagonal torus are
\[
 T_d=\operatorname{Spec}E[a,a^{-1},(a^2-d)^{-1}],\qquad
 T=\mathbb G_m^2.
\]
The involution $w$ sends $a$ to $d/a$ and exchanges the factors of $T$.

\begin{proposition}[Eigenvalues and eigenline automorphisms]
\label{v4-arith:ic:w5-e:local:frobenius-stack}
There is an equivalence
\[
 \mathcal U_d^{\mathrm{rs}}\simeq[T_d/(T\rtimes\langle w\rangle)],
\]
where $T$ acts trivially on $T_d$. The scheme of unordered eigenvalues is
\[
 U_d=\operatorname{Spec}E[t,(t^2-4d)^{-1}].
\]
It has dimension one. Over an algebraically closed field, each
geometric object of $\mathcal U_d^{\mathrm{rs}}$ has automorphism group
$\mathbb G_m^2$.
\end{proposition}

\begin{proof}
After a finite \textrm{\'etale} base change, the polynomial
$X^2-tX+d$ has distinct roots $a,b$, with $b=d/a$.
The endomorphisms
\[
 e_a=\frac{F-b}{a-b},\qquad e_b=\frac{F-a}{b-a}
\]
are complementary idempotents. Their images are rank-one projective
modules, and $F$ acts on them by $a$ and $b$.
Conversely, two line bundles with these scalar actions give $(V,F)$.

Changing their frames gives the $T$-action. Exchanging their order
gives the involution. Descent gives the stated quotient stack.

The map $a\mapsto t=a+d/a$ identifies the unordered roots.
Its coordinate extension satisfies $a^2-ta+d=0$.
Inverting $t^2-4d$ makes this a finite \textrm{\'etale} double cover
with the stated involution. The invariant ring is
$E[t,(t^2-4d)^{-1}]$. Each automorphism preserves the two eigenspaces
and scales them independently.
\end{proof}

The atlas $Y_d$ has dimension three, since the determinant equation
has nonzero differential on invertible matrices.
Thus $\mathcal U_d$ has dimension $3-4=-1$ as an algebraic stack.
Its eigenvalue scheme has a different dimension and omits stabilizers.
If a trivialization of $\bigwedge^2V$ is required, the quotient uses
$SL_2$ instead. Its regular semisimple stabilizer is $\mathbb G_m$.
These are distinct moduli conventions.

\subsection{An unramified deformation ring}

Let $\mathcal O$ be a complete discrete valuation ring of mixed
characteristic, with finite residue field $k$ of characteristic $p$.
Fix a nonarchimedean local field $L$. Write $G_L$ for its absolute
Galois group and $I_L$ for inertia, so $G_L/I_L\simeq\widehat{\mathbb Z}$.
Fix $\overline F\in GL_2(k)$ and $d\in\mathcal O^\times$ lifting
$\det\overline F$.

A test algebra is a local Artinian
$\mathcal O$-algebra $A$ with residue field $k$.
An unramified framed deformation sends arithmetic Frobenius to
$F_A\in GL_2(A)$, with $F_A\bmod\mathfrak m_A=\overline F$
and $\det F_A=d$.
Here arithmetic Frobenius acts on a residue field by its cardinality
power. Inertia acts trivially.

\begin{proposition}[Framed deformations and distinct residual roots]
\label{v4-arith:ic:w5-e:local:framed-unramified}
Choose a matrix lift $F_0\in M_2(\mathcal O)$ of $\overline F$.
The framed deformation functor is represented by
\[
 R_d^\square=
 \mathcal O[[X_{11},X_{12},X_{21},X_{22}]]/
       (\det(F_0+X)-d).
\]
This ring is formally smooth of relative dimension three.
Suppose additionally that
$\overline F=\operatorname{diag}(\bar a,\bar b)$, with
$\bar a\ne\bar b$.

Choose $a_0\in\mathcal O^\times$ lifting $\bar a$
and put $b_0=d/a_0$. Then
\[
 R_d^\square\simeq\mathcal O[[x,y,z]],\qquad
 F=\begin{pmatrix}
 a_0+x&y\\ z&(d+yz)/(a_0+x)
 \end{pmatrix}.
\]
Strict conjugacy classes of these deformations are represented by
$\mathcal O[[u]]$. The corresponding ordered eigenvalues are
$a=a_0+u$ and $b=d/a$.
\end{proposition}

\begin{proof}
Every test algebra is finite. A matrix in its finite group $GL_2(A)$
defines a unique continuous homomorphism from $\widehat{\mathbb Z}$.
Thus specifying $F_A$ is equivalent to specifying the unramified
representation. The matrix entries and determinant equation give the
first ring.

At least one cofactor of $\overline F$ is nonzero. The derivative in
the corresponding matrix coordinate is a unit. Solving for that
coordinate proves formal smoothness and leaves three parameters.
In the displayed diagonal case, $a_0+x$ is a unit and the determinant
equation solves directly for the lower-right entry.

For each $F_A$, Hensel lifting gives the unique characteristic root
$a$ reducing to $\bar a$. The other root is $b=d/a$.
Since $a-b$ is a unit, the projectors above split $A^2$ into two
free rank-one eigenspaces. Their bases can be chosen to reduce to the
standard residual bases.

The resulting change of basis is the identity
modulo $\mathfrak m_A$, so it is a strict conjugation.
Two diagonal lifts are strictly conjugate exactly when their ordered
roots agree. The choice of $a-a_0\in\mathfrak m_A$ therefore represents
the set of strict conjugacy classes.
\end{proof}

The last ring represents classes, not the deformation groupoid.
For a diagonal lift, strict automorphisms contain
$(1+\mathfrak m_A)^2$. A quotient stack retains these automorphisms.
The construction concerns unramified lifts only. A crystalline lifting
condition or a larger Emerton--Gee component requires its own map and
coefficient conventions.


\section{Ordinary quotients and characteristic-zero types}
\label{v4-arith:ic:ordinary}

Fix $p\ge5$. Let $E/\mathbb Q_p$ be finite, with integer ring
$\mathcal O$, uniformizer $\varpi$, and finite residue field $k$.
Put $G=G_{\mathbb Q_p}$ and let $I\subset G$ be inertia.
Fix a continuous representation
$\bar\rho:G\to\operatorname{GL}_2(k)$, a basis of $k^2$, and a
continuous character $\delta:G\to\mathcal O^\times$ lifting
$\det\bar\rho$.

Give an Artinian local $\mathcal O$-algebra $A$ with residue field $k$
the discrete topology. A framed lift is a continuous homomorphism
$\rho_A:G\to\operatorname{GL}_2(A)$ reducing to the specified matrix
representation $\bar\rho$. Framed lifts are not divided by conjugation.
The framed deformation ring $R^\square$ is a complete Noetherian
local $\mathcal O$-algebra with universal representation
$\rho^\square$.

Its existence is the framed representability statement
in (Kisin, \emph{The structure of potentially semi-stable deformation rings}, \S1.2).
Set
\[
 R=R^{\square,\delta}
 =R^\square/\overline{(\det\rho^\square(g)-\delta(g):g\in G)}.
\]
The bar denotes closure in the maximal-ideal topology.
This quotient represents framed lifts of determinant $\delta$.
Write its universal matrices as
\[
 \rho_R(g)=\begin{pmatrix}a_g&b_g\\c_g&d_g\end{pmatrix}.
\]

Assume a chosen residual quotient
$\bar q:k^2\twoheadrightarrow k(\bar\eta)$ is unramified.
Here $k(\bar\eta)$ is the one-dimensional representation with
character $\bar\eta$, and unramified means $\bar\eta|_I=1$.
Choose the basis so that $\bar q=(0,1)$.
Then
\[
 \bar\rho(g)=
 \begin{pmatrix}\bar\alpha(g)&*\\0&\bar\eta(g)\end{pmatrix},
 \qquad \bar\alpha=\bar\delta\bar\eta^{-1}.
\]

An ordinary lift relative to $\bar q$ is a framed lift together with
a free rank-one quotient lifting $\bar q$, on which inertia acts
trivially. Quotients are identified up to an equivariant target
isomorphism. Each such quotient has a unique normalized row
$q=(t,1)$, where $t\in\mathfrak m_A$.
Its character is allowed to vary with $A$.

\begin{proposition}[explicit ordinary deformation equations]
\label{v4-arith:ic:ordinary-local:equations}
Suppose $\bar\alpha\ne\bar\eta$. Choose $g_0\in G$ with
$\bar\alpha(g_0)\ne\bar\eta(g_0)$.
There is a unique $t_0\in\mathfrak m_R$ satisfying
\begin{equation}
\label{v4-arith:ic:ordinary-local:hensel}
 b_{g_0}t_0^2+(d_{g_0}-a_{g_0})t_0-c_{g_0}=0.
\end{equation}
Let $J_{\mathrm{ord}}\subset R$ be the closed ideal generated by
\[
 \begin{split}
 b_gt_0^2+(d_g-a_g)t_0-c_g&\quad(g\in G),\\
 t_0b_h+d_h-1&\quad(h\in I).
 \end{split}
\]

Then $R^{\mathrm{ord}}=R/J_{\mathrm{ord}}$ represents ordinary
lifts relative to $\bar q$. The quotient line is uniquely determined
by the framed representation. Thus forgetting the quotient gives
a closed immersion
$\operatorname{Spf}R^{\mathrm{ord}}\hookrightarrow\operatorname{Spf}R$.
\end{proposition}

\begin{proof}
The polynomial in \eqref{v4-arith:ic:ordinary-local:hensel} vanishes at
zero modulo $\mathfrak m_R$. Its derivative there is
$\bar\eta(g_0)-\bar\alpha(g_0)\ne0$.
Completeness and Hensel's lemma give the unique root $t_0$.
For completeness, successive correction modulo powers of
$\mathfrak m_R$ solves a linear equation with this invertible
coefficient.

These corrections converge in $R$.
If two roots lie in $\mathfrak m_R$, their difference times
$b_{g_0}(t+t')+d_{g_0}-a_{g_0}$ is zero. The second factor is a
unit, so the roots agree.

For a row $q=(t,1)$, stability of its kernel under $\rho_A(g)$
is equivalent to
\[
 q\rho_A(g)=\eta_A(g)q,\qquad
 \eta_A(g)=t b_g+d_g,
\]
where the matrix entries are specialized to $A$.
The first component is exactly
$b_gt^2+(d_g-a_g)t-c_g=0$.
Since $\eta_A(g)$ reduces to $\bar\eta(g)\ne0$, it is a unit.
Surjectivity of $q$ and multiplication of the matrices show that
$\eta_A$ is a character. It is continuous because $\rho_A$ is
continuous. The equations indexed by $I$ make it unramified.

Every ordinary quotient satisfies the equation for $g_0$.
The uniqueness argument over the Artinian local ring $A$ forces
$t$ to be the image of $t_0$.
The remaining equations are precisely the images of
$J_{\mathrm{ord}}$. Conversely, a map $R^{\mathrm{ord}}\to A$
gives the quotient $(t_0,1)$ with these properties.

Both constructions commute with local $\mathcal O$-algebra maps
of Artinian coefficient rings. The construction also applies to
complete Noetherian local coefficient rings by passage to their Artinian
quotients. The normalized rows are compatible by uniqueness.
\end{proof}

The kernel of this quotient is freely generated by $(1,-t_0)$.
Consequently every ordinary lift has an exact sequence
\begin{equation}
\label{v4-arith:ic:ordinary-local:sequence}
 0\longrightarrow A(\delta\eta_A^{-1})
 \longrightarrow A^2_{\rho_A}
 \longrightarrow A(\eta_A)\longrightarrow0,
 \qquad \eta_A|_I=1.
\end{equation}
Indeed the action on the kernel is multiplication by
$a_g-b_gt_0$, and the determinant is the product of the two
characters. The sequence need not split.
To prescribe an unramified lift $\eta:G\to\mathcal O^\times$,
add the equations $t_0b_g+d_g=\eta(g)$ for $g\in G$.

\begin{remark}[coincident residual characters]
Without $\bar\alpha\ne\bar\eta$, the quotient must be retained
as part of the deformation problem. That problem is still represented
by
\[
 R[[t]]\big/
 \overline{\bigl(
 b_gt^2+(d_g-a_g)t-c_g\ (g\in G),\quad
 tb_h+d_h-1\ (h\in I)
 \bigr)}.
\]
Here $t$ maps into the maximal ideal of each coefficient ring.
The same row calculation proves the assertion.

For the trivial residual representation and trivial determinant,
take the constant trivial lift to $k[u]/(u^2)$.
The rows $(a u,1)$ for $a\in k$ give distinct ordinary quotient
lines lifting $(0,1)$. Thus the forgetful map need not be injective
on deformation functors when the residual characters coincide.
\end{remark}

\subsection{Hodge types on characteristic-zero lifts}

The ordinary equations make sense over Artinian coefficient rings.
Hodge and inertial types require characteristic-zero period modules.
For a continuous representation $V$ on a finite-dimensional vector
space over a finite extension $E'/E$, put
\[
 D_{\mathrm{dR}}(V)
   =(B_{\mathrm{dR}}\otimes_{\mathbb Q_p}V)^G.
\]
Here $B_{\mathrm{dR}}$ is Fontaine's filtered de Rham period field.

The representation is de Rham when
$\dim_{E'}D_{\mathrm{dR}}(V)=\dim_{E'}V$.
We use filtration degrees as Hodge--Tate weights, so the cyclotomic
character $\varepsilon$ has weight $-1$.
A type $\mathbf h=\{h_1,h_2\}$ prescribes the dimensions of
$\operatorname{gr}^jD_{\mathrm{dR}}(V)$.
Equivalently, Kisin's cocharacter is the conjugacy class of
$z\mapsto\operatorname{diag}(z^{-h_1},z^{-h_2})$.

Let $B_{\mathrm{cris}}\subset B_{\mathrm{st}}$ be Fontaine's
crystalline and semistable period rings. For a finite extension
$K'/\mathbb Q_p$ with maximal unramified subfield $K'_0$, set
\[
 D_{\mathrm{st},K'}(V)=(B_{\mathrm{st}}\otimes_{\mathbb Q_p}V)^{G_{K'}},
 \qquad
 D_{\mathrm{cris},K'}(V)=(B_{\mathrm{cris}}\otimes_{\mathbb Q_p}V)^{G_{K'}}.
\]
The representation is potentially semistable, respectively potentially
crystalline, when for some $K'$ the corresponding module is free of
rank $\dim_{E'}V$ over $K'_0\otimes_{\mathbb Q_p}E'$.
Its potentially semistable period module is
\[
 D_{\mathrm{pst}}(V)=\varinjlim_{K'}D_{\mathrm{st},K'}(V).
\]
Inertia acts linearly on this module over
$\overline{\mathbb Q}_p^{\mathrm{nr}}\otimes_{\mathbb Q_p}E'$.

An inertial type is a continuous representation
$\tau:I\to\operatorname{GL}_2(E)$ with open kernel which extends to
the Weil group. The representation $V$ has this
type when its inertia representation on $D_{\mathrm{pst}}(V)$
is isomorphic to $\tau$ after scalar extension. Here
$\overline{\mathbb Q}_p^{\mathrm{nr}}$ is the maximal unramified
extension of $\mathbb Q_p$ inside its chosen algebraic closure.
This is the period-module action, not the inertia action on $V$.

Choose $E$ large enough to define $\tau$ and the Hodge type.
Kisin constructs quotients $R^{\square,\mathbf h,\tau}_{\mathrm{pst}}$
and $R^{\square,\mathbf h,\tau}_{\mathrm{pcris}}$ of $R^\square$.
Their characteristic-zero quotients impose the potentially semistable
and potentially crystalline type conditions on finite local
characteristic-zero coefficient algebras. In particular, their field-valued points specify lifts of these types.
These are Kisin's Theorems 1.2.1--1.2.2, pp.~6--7.

For $\star\in\{\mathrm{pst},\mathrm{pcris}\}$ define the actual
fixed-determinant intersection ring and its generic fiber by
\[
 S_\star=R^{\square,\mathbf h,\tau}_\star
             \otimes_{R^\square}R,
 \qquad X_\star=\operatorname{Spec}S_\star[1/p].
\]
Both factors are quotient rings, so this tensor product is an
ordinary quotient of $R^\square$. It is already complete and
Noetherian. A continuous map $S_\star\to\mathcal O_{E'}$
gives a framed lattice of determinant $\delta$ whose rational
representation has the indicated type.

Conversely, choose a frame of such a lattice reducing to the
given residual frame. Kisin's factorization criterion applies over
$E'$. The defining ideal then vanishes in the torsion-free ring
$\mathcal O_{E'}$, giving the required integral map. Here coefficient extensions
carry a fixed residue-field embedding, and all integral maps are
continuous local $\mathcal O$-algebra maps.

The scheme $X_\star$ also retains the infinitesimal structure of
the intersection. No flatness or reducedness of $S_\star$ is needed
for this construction. Incompatible determinant and type conditions
can give an empty generic fiber.

The schematic closure of the generic fiber in $\operatorname{Spec}S_\star$
is defined by the $p$-power torsion ideal. Its coordinate ring is
$S_\star/S_\star[p^\infty]$, where
$S_\star[p^\infty]=\{x:p^n x=0\text{ for some }n\}$.
This quotient is $p$-torsion-free, hence flat over the discrete
valuation ring $\mathcal O$. It has no effect after inverting $p$,
but can change the special fiber. The integral intersection and its generic fiber
must therefore remain distinct.

Under the distinguished quotient hypothesis, define
\[
 X_\star^{\mathrm{ord}}
 =\operatorname{Spec}\bigl((S_\star/J_{\mathrm{ord}}S_\star)[1/p]\bigr),
 \qquad
 X_\star^{\mathrm{rest}}=X_\star\setminus X_\star^{\mathrm{ord}}.
\]
The first is a closed subscheme and the second its open complement.
This gives a two-piece decomposition within the specified lifting
space. The ordinary condition is relative to the chosen residual
quotient $\bar q$.

\begin{corollary}[ordinary lifts of a prescribed type]
\label{v4-arith:ic:ordinary-local:typed-ordinary}
A continuous point $S_\star\to\mathcal O_{E'}$ lies in the ordinary
piece precisely when its framed lattice has an unramified quotient
lifting $\bar q$. On this piece its rational representation has
an exact sequence
\[
 0\longrightarrow E'(\delta\eta^{-1})\longrightarrow V
 \longrightarrow E'(\eta)\longrightarrow0,
 \qquad \eta|_I=1,
\]
and has the prescribed Hodge and inertial types.
\end{corollary}

\begin{proof}
Factoring through $R^{\mathrm{ord}}$ is equivalent to the ordinary
quotient property by the proposition. Apply
\eqref{v4-arith:ic:ordinary-local:sequence} and invert $p$.
The quotient $R^{\square,\mathbf h,\tau}_\star$ imposes the type.
If the equations vanish only after inverting $p$, they already vanish
in the torsion-free ring $\mathcal O_{E'}$.
\end{proof}

For the potentially Barsotti--Tate problem in the Tate-module
convention, take $\star=\mathrm{pcris}$ and
$\mathbf h=\{-1,0\}$, with a specified inertial type.
The dual convention has weights $\{0,1\}$.
Other Hodge types require other quotient rings.
These spaces do not give an exhaustive decomposition of all residual
points.

They describe lifts above a fixed residual representation.
For example, with $a=p-1$, the representations
$1\oplus\varepsilon$ and
$\varepsilon^a\oplus\varepsilon^{1-a}$ have the same reduction and
determinant. Their Hodge--Tate multisets are $\{0,-1\}$ and
$\{1-p,p-2\}$, respectively.


\section{A local categorical realization}
\label{v4-arith:ic:deg}

The formal lifting spaces above have natural maps into a spectral
stack at the coefficient prime. The determinant identification
belongs to this stack. It is more structure than the determinant
value used for the Frobenius quotient stack.



To compare these deformation spaces with categorical local Langlands,
fix a continuous central character
$\zeta:\mathbb Q_p^\times\to\mathcal O^\times$.
Use local reciprocity sending a uniformizer to geometric Frobenius,
and require the already fixed determinant to satisfy
$\delta=\zeta\varepsilon^{-1}$ as a Galois character.
These are the conventions of (Dotto--Emerton--Gee, arXiv:2603.26887, \S1.3, p.~8).

Let $\mathcal A_\zeta$ be smooth $\operatorname{GL}_2(\mathbb Q_p)$
representations on $p$-power torsion $\mathcal O$-modules with
central character $\zeta$. Let $\mathcal X_\delta$ classify rank-two projective \'etale
$(\varphi,\Gamma)$-modules with a determinant identification.
This identification is an isomorphism with the rank-one module
attached to $\delta$. Morphisms preserve this isomorphism.

This is the fiber-product definition in
(Dotto--Emerton--Gee, arXiv:2603.26887, Definition C.1.1, p.~265).
The Dotto--Emerton--Gee functor is
\[
 F:D(\mathcal A_\zeta)\longrightarrow
          \operatorname{Ind}D^b_{\mathrm{coh}}(\mathcal X_\delta).
\]
It is exact, colimit-preserving and fully faithful
(Dotto--Emerton--Gee, arXiv:2603.26887, Theorem 1.1.1 and its following paragraph, p.~3).

The functor has an actual kernel. Write
\(\mathcal O[[G]]_\zeta\) for the completed group algebra with
central character \(\zeta\), where \(G=GL_2(\mathbb Q_p)\).
Dotto--Emerton--Gee construct a pro-coherent right
\(\mathcal O[[G]]_\zeta\)-module \(\mathcal L_\infty\) on
\(\mathcal X_\delta\), and on the bounded finitely presented source
their functor is
\[
 F(\pi)=\mathcal L_\infty
              \otimes_{\mathcal O[[G]]_\zeta}^{\mathbb L}\pi.
\]
The right action has a specified variance.

The initial Colmez
construction has left central character
\(\zeta\varepsilon^{-2}\). Twisting by
\((\zeta^{-1}\varepsilon)\circ\det\) changes this to
\(\zeta^{-1}\). Precomposition with \(g\mapsto g^{-1}\)
then gives the right \(\zeta\)-action.

On a smooth formal chart \(\operatorname{Spf}A\), the kernel comes
from the compatible systems attached to the étale modules over
\(A/I^n\). Their coefficient base change and descent construct the
global pro-coherent object. Bounded coherence of the displayed tensor and full faithfulness
are separate theorems of that construction.
Their locators are Definitions 5.1.31--5.1.33,
Proposition 5.3.23 and Theorem 5.5.1.
A symbol for the tensor product alone would prove neither assertion.


Let $\bar\theta$ be the residual pseudorepresentation of $\bar\rho$.
Complete $\mathcal X_\delta$ along the fiber over $\bar\theta$.
This classifies Galois representations deforming that
pseudorepresentation, with the specified determinant identification
(Dotto--Emerton--Gee, \S1.1.6, p.~4).
The universal representation defines a morphism
\[
 \operatorname{Spf}R\longrightarrow
       \widehat{\mathcal X}_{\delta,\bar\theta}
       \longrightarrow\mathcal X_\delta
\]
as follows.

Reduce $\rho_R$ modulo each power of its maximal ideal.
Over each finite coefficient ring, apply the equivalence between
continuous Galois representations and \'etale $(\varphi,\Gamma)$-modules.
Retain the determinant identification sending the wedge of the
two frame vectors to $1$ in the character module of $\delta$.
The construction preserves tensor products and base change.
See Dotto--Emerton--Gee, Appendix C.2, especially (C.2.1)--(C.2.2),
p.~268, and Remark~C.2.17, p.~272.

The compatible objects define the displayed formal morphism.
Their residual pseudorepresentation is $\bar\theta$, so the map
factors through the indicated completion. A change of basis preserving
the determinant identification gives an isomorphism in that stack.
Scalar automorphisms $\pm1$ remain.

The ordinary and type quotients give
morphisms into this diagram. For the type quotients these are the maps from
$\operatorname{Spf}S_\star$. The ordinary intersection uses
$\operatorname{Spf}(S_\star/J_{\mathrm{ord}}S_\star)$.
These are morphisms of formal moduli problems, whereas $X_\star$
denotes the separately defined scheme after inverting $p$. They do not identify the source formal schemes with the
stack, which retains automorphisms.

An ordinary or potentially Barsotti--Tate comparison must specify
full stable source and target subcategories. Essential surjectivity
first requires a proof that $F$ restricts between them. The deformation maps alone do not specify these
categories or their generators. The functor on the entire target
is not an equivalence (Dotto--Emerton--Gee, arXiv:2603.26887, Remark 1.1.2, p.~3).
This statement does not decide equivalence on a specified smaller target.

\begin{proposition}[generation criterion for a restricted functor]
\label{v4-arith:ic:ordinary-local:generation}
Let $\mathcal C,\mathcal D$ be presentable stable
$\infty$-categories. Suppose an exact colimit-preserving fully
faithful functor $F_0:\mathcal C\to\mathcal D$ is supplied.
If objects $P_i\in\mathcal C$ have images generating $\mathcal D$
under colimits, shifts and cofibers, then $F_0$ is an equivalence.
\end{proposition}

\begin{proof}
Full faithfulness makes the essential image a full subcategory
equivalent to $\mathcal C$. Every diagram in that image can
therefore be lifted to $\mathcal C$.
Preservation of colimits makes its colimit in $\mathcal D$ another
image object. Exactness gives closure under shifts and cofibers.
The essential image contains the specified generators and hence
all of $\mathcal D$.
\end{proof}

The local arithmetic construction needed for this criterion is the
identification of $F(P_i)$ with actual spectral objects, together
with a proof of target generation. If lifting conditions overlap,
their union must generate all required extension objects.
A direct-sum assembly additionally requires orthogonality.
Neither modularity of a lift nor a ring of scalar actions establishes
these categorical properties.

\paragraph{Line and quotient conventions.}
The ordinary lifting problem just constructed retains an unramified
\emph{quotient}. Its kernel has character $\delta\eta^{-1}$ and need
not be unramified. An ordinary-line condition instead asks for
an unramified subrepresentation.

Dualizing a quotient
of $V$ gives an unramified subrepresentation of $V^\vee$, with determinant
$\delta^{-1}$. It does not identify the two conditions on the same $V$.
Any use of the line convention in the comparison below must specify
this dualization or another proved flag comparison, including the
central-character and Hodge-sign changes.


\subsection{Smooth objects with the required central character}

Suppose the local residual rank-two representation is unramified
at \(p\), and use the determinant convention
\(\det\bar\rho=\chi\bar\varepsilon^{-1}\).
Under geometric reciprocity,
\(\bar\varepsilon(u)=\bar u\) for \(u\in\mathbb Z_p^\times\)
and \(\bar\varepsilon(p)=1\).
Consequently \(\chi(u)=\bar u\).
An induction from two unramified characters has a different
central character on units. The following induction gives actual
objects in the required torsion category.

\begin{proposition}[central-character-compatible smooth source objects]
\label{v4-arith:ic:w5-e:local:ramified-source}
Fix a finite extension of $\mathbb Q_p$ with integer ring $\mathcal O$
and residue field $k$. Choose a continuous lift
$\zeta:\mathbb Q_p^\times\to\mathcal O^\times$ of $\chi$.
Under geometric reciprocity, set $\delta=\zeta\varepsilon^{-1}$,
where $\varepsilon$ is the $p$-adic cyclotomic character. Let $A$ be a finite local Artinian $\mathcal O$-algebra
with residue field $k$. Write $\zeta_A$ for the induced character,
and choose an unramified character $\eta_A:\mathbb Q_p^\times\to A^\times$.

Put $G=GL_2(\mathbb Q_p)$ and let $B$ be its upper triangular subgroup.
The unnormalized smooth induction
\[
 I_A=\operatorname{Ind}_B^G(\zeta_A\eta_A^{-1}\otimes\eta_A)
\]
consists of locally constant $f:G\to A$ satisfying
\[
 f\left(\begin{pmatrix}a&b\\0&d\end{pmatrix}g\right)
 =\zeta_A(a)\eta_A(a)^{-1}\eta_A(d)f(g),
 \qquad (g_0f)(g)=f(gg_0).
\]
It is a nonzero smooth $p$-power-torsion $\mathcal O$-representation
with central character $\zeta$. It has no $GL_2(\mathbb Z_p)$-fixed
vectors. In particular, no square-root modulus factor is used over $A$.
\end{proposition}
\begin{proof}
The finite coefficient characters have open kernels. Choose a principal
congruence subgroup $K_m\subset GL_2(\mathbb Z_p)$ so that the inducing
character $\theta_A$ of $B$ is trivial on $B\cap K_m$.
Define $f(bk)=\theta_A(b)$ on $BK_m$ and zero outside. Two
factorizations differ by $B\cap K_m$, so this is well-defined.
The image of $K_m$ in the Hausdorff space $B\backslash G$ is compact
and open, hence closed as well.

The resulting zero extension is
locally constant and $f(1)=1$.
Restrict a locally constant induced function to the compact group
$GL_2(\mathbb Z_p)$. A finite clopen partition refines its local
constancy neighborhoods. It is therefore fixed on the right by
some sufficiently small open subgroup. This proves
smoothness.

The ring $A$ is killed by a power of $p$, so the same is
true of $I_A$. A central scalar $x$ acts by
$\zeta_A(x)\eta_A(x)^{-1}\eta_A(x)=\zeta_A(x)$.
Choose a unit $u$ whose reduction is not $1$. The element
$\zeta_A(u)-1$ is a unit of $A$. If a vector were fixed by
$GL_2(\mathbb Z_p)$, it would be annihilated by this unit, hence vanish.
\end{proof}

A determinant twist cannot turn an unramified principal series into
this residual central-character family when $p$ is odd. On the cyclic
group $\mathbb F_p^\times$, such a twist changes the central exponent
by $2a$. The required exponent is $1$. The congruence
$2a=1\pmod{p-1}$ has no solution. This remains an obstruction after
coefficient extension and, by reduction, for the finite rings $A$.
The characteristic-zero unramified invariant-vector calculation has
a different central-character carrier.


Let \(\mathcal C_{\mathrm{ind}}\) be the stable colimit closure
of these induced objects in \(D(\mathcal A_\zeta)\).
It is a well-defined source category. To identify it with an
ordinary comparison category one must identify its spectral image
and generators, including the chosen ordinary flag.
An invariant line and an invariant quotient are distinct conditions.

\section{Quadratic induction and ordinary lifting spaces}

A quadratic global induction can become reducible at a nonsplit local
place. The local character determines this behavior. Ordinary lifting
conditions require further data: a residual quotient line, a determinant,
and a characteristic-zero lifting space. These data specify the local
objects whose images under a categorical Langlands functor must be
compared.

\subsection{Restriction of a quadratic induction}

Let $K/\mathbb Q$ be quadratic. Fix a place $v$ of $\mathbb Q$ and
compatible decomposition groups. Let $L$ be a coefficient field, and
let $\psi:G_K\to L^\times$ be continuous. Here $G_F$ denotes the absolute Galois group of a field $F$,
and $L$ is equipped with its specified field topology. All representation
spaces carry the corresponding finite-dimensional topology. Induction from an open subgroup
of finite index means ordinary algebraic induction, with its induced
continuous action.

If $v$ splits in $K$, the double-coset decomposition gives
\[
 \left.\operatorname{Ind}_{G_K}^{G_{\mathbb Q}}\psi\right|_{G_{\mathbb Q_v}}
 \simeq \psi_w\oplus\psi_{w'},
\]
where $w,w'$ are the two places above $v$. The two local characters
can coincide. If $v$ is nonsplit, there is one place $w$ above $v$,
and the same decomposition gives
\[
 \left.\operatorname{Ind}_{G_K}^{G_{\mathbb Q}}\psi\right|_{G_{\mathbb Q_v}}
 \simeq\operatorname{Ind}_{G_{K_w}}^{G_{\mathbb Q_v}}\psi_w.
\]
These formulas follow by separating functions according to their support
on the double cosets $G_K\backslash G_{\mathbb Q}/G_{\mathbb Q_v}$.
There are two cosets in the split case and one in the nonsplit case.

To determine the second representation, let $H$ be an open normal
subgroup of index two in a profinite group $G$. Choose $s\in G\setminus H$.
For a continuous character $\chi:H\to L^\times$, put
\[
 \chi^s(h)=\chi(s^{-1}hs),\qquad c=\chi(s^2).
\]
On the induced basis $e_1=1\otimes1$, $e_2=s\otimes1$, the action is
\begin{equation}
\label{v4-arith:ic:quadratic-local:matrices}
 \rho(h)=\begin{pmatrix}\chi(h)&0\\0&\chi^s(h)\end{pmatrix},
 \qquad
 \rho(s)=\begin{pmatrix}0&c\\1&0\end{pmatrix}.
\end{equation}

\begin{proposition}[local irreducibility criterion]
\label{v4-arith:ic:quadratic-local:criterion}
If $\chi\ne\chi^s$, then $\operatorname{Ind}_H^G\chi$ is absolutely
irreducible. If $\chi=\chi^s$, its reducibility over $L$ is equivalent
to the existence of a root of $T^2-c$ in $L$.

Suppose $\chi=\chi^s$ and $b^2=c$ with $b\in L$. There is a continuous
extension $\widetilde\chi:G\to L^\times$ with
$\widetilde\chi|_H=\chi$ and $\widetilde\chi(s)=b$.
If $\operatorname{char}L\ne2$, then
\[
 \operatorname{Ind}_H^G\chi
 \simeq\widetilde\chi\oplus(\widetilde\chi\,\omega_{G/H}),
\]
where $\omega_{G/H}$ is trivial on $H$ and sends $s$ to $-1$.
In characteristic two the induced representation is a nonsplit
self-extension of $\widetilde\chi$.
\end{proposition}

\begin{proof}
When the characters differ, choose $h\in H$ with
$\chi(h)\ne\chi^s(h)$. Any invariant line over an algebraic closure
must be one of the two eigenspaces of $\rho(h)$. The matrix $\rho(s)$
interchanges them, so neither is $G$-stable.

When the characters agree, $H$ acts by scalars. A $G$-stable line
therefore exists precisely when $\rho(s)$ has an eigenvalue in $L$.
Its characteristic and minimal polynomials are both $T^2-c$.
The prescribed values of $\widetilde\chi$ respect conjugation by $s$
and the relation $s^2\in H$, hence define a character.

It is continuous because $H$ is open and $\chi$ is continuous.
In characteristic different from two the two roots $b,-b$ are distinct,
which gives the displayed direct sum. In characteristic two the
minimal polynomial is $(T-b)^2$. The matrix is not scalar, so its
unique eigenline has no invariant complement.
\end{proof}

Thus over an algebraically closed field, local irreducibility in the
nonsplit case is equivalent to $\psi_w\ne\psi_w^s$. Inertness and
ramification alone do not decide this equality. Over a smaller
coefficient field, an invariant character can instead give an
irreducible representation which becomes reducible after scalar extension.
Over coefficient rings the matrices in
\eqref{v4-arith:ic:quadratic-local:matrices} still define induction.
The field criterion applies after specialization and does not assert
an integral direct-sum decomposition.

For a ramified quadratic extension of $\mathbb Q_p$ with $p\ge5$,
the ramification index is two, so the extension is tame.
A characteristic-zero inducing character can still be nontrivial on
wild inertia. For a continuous character valued in
$\overline{\mathbb F}_p^\times$ with the discrete topology, its image
is finite of order prime to $p$. Its restriction to wild inertia is
therefore trivial. Since the quadratic extension is tame, wild inertia
lies in the inducing subgroup and acts trivially on the induced
residual representation.

\begin{example}[an inert place with reducible restriction]
Let $L_0$ be the splitting field of $X^3-2$. Eisenstein's criterion at
$2$ and the discriminant $-108$ give
$\operatorname{Gal}(L_0/\mathbb Q)=S_3$. The quadratic subfield is
$K=\mathbb Q(\sqrt{-3})$. Over $\overline{\mathbb F}_5$, induction
of a nontrivial character of $A_3$ gives an absolutely irreducible
representation of $S_3$. Complex conjugation is a transposition,
so this representation is odd.

The prime $5$ is unramified in $L_0$, and
\[
 X^3-2=(X-3)(X^2+3X+4)\pmod5.
\]
The quadratic discriminant is $3$, which is nonsquare modulo $5$.
Frobenius is consequently a transposition. Also $-3=2\pmod5$ is
nonsquare, so $5$ is inert in $K$. The local representation factors
through the order-two decomposition group and splits into its
$+1$ and $-1$ Frobenius eigenspaces.
\end{example}


\begin{lemma}[restriction of a dihedral residual representation]
\label{v4-arith:ic:w6-e:lem:CP-induced}
For a quadratic induction, the split restriction is the sum of the two
local characters, which can coincide. At a nonsplit place the restriction is induced from the local
quadratic subgroup. Over an algebraically closed coefficient field,
it is irreducible exactly when its two conjugate characters differ. Ramification of the quadratic
extension alone specifies no Hodge type. At $p\ge5$ that quadratic
extension is tame, and a discrete mod-$p$ character is trivial on wild inertia.
\end{lemma}
\begin{proof}
The double-coset decomposition, matrix computation and characteristic-two
case in Proposition~\ref{v4-arith:ic:quadratic-local:criterion} prove these
statements. The inert prime $5$ example shows why inertness alone cannot
imply irreducibility. A comparison with the Colmez Montreal functor,
including its actual block and coefficient category, remains additional
arithmetic data. It is not a consequence of the restriction formula.
\end{proof}


\section{Residual parameters and deformation objects}
\label{v4-arith:ic:w10-d:X2}

The complex parameters of residual automorphic representations and the mod-$p$
parameters of a deformation stack have different coefficient fields.
Their comparison requires more than the same two-dimensional notation.
Even over $\mathbb C$, the two $SL_2$ factors in a local Arthur parameter
have different roles.

\begin{definition}[The non-tempered parameter]
\label{v4-arith:ic:w10-d:def:arthur-param-s1}
At a finite place, use the domain
$W_F\times SL_2(\mathbb C)_{\mathrm D}\times SL_2(\mathbb C)_{\mathrm A}$.
The parameter takes values in $GL_2(\mathbb C)$.
Here $W_F$ is the Weil group, the factor marked $\mathrm D$ is the
Deligne factor, and the factor marked $\mathrm A$ is the additional
Arthur factor. The trivial-character rank-two non-tempered parameter is
\begin{equation}\label{v4-arith:ic:w10-d:eq:arthur-param-s1}
 \psi_1(w,g_{\mathrm D},g_{\mathrm A})=g_{\mathrm A}.
\end{equation}
The associated Langlands parameter is obtained by substituting
$\operatorname{diag}(|w|^{1/2},|w|^{-1/2})$ in the Arthur factor.
This convention records its non-tempered exponents.
\end{definition}

\begin{definition}[The monodromy parameter]
\label{v4-arith:ic:w10-d:def:arthur-param-s0}
On the same domain, the standard rank-two Deligne parameter, regarded
as an Arthur parameter trivial on its additional factor, is
\begin{equation}\label{v4-arith:ic:w10-d:eq:arthur-param-s0}
 \psi_0(w,g_{\mathrm D},g_{\mathrm A})=g_{\mathrm D}.
\end{equation}
The distinction from \eqref{v4-arith:ic:w10-d:eq:arthur-param-s1} lies in
which factor acts. Restricting both displays to an unnamed $SL_2$
would erase that distinction. A determinant twist changes the Weil
part and does not exchange the two factors.
\end{definition}

\begin{remark}[Local and global classification]
\label{v4-arith:ic:w10-d:rem:arthur-gln}
Arthur's parameter formalism separates the Deligne and Arthur factors.
His endoscopic classification concerns orthogonal and symplectic groups,
with $GL_n$ classification as input. The residual classification for
$GL_n$ is the theorem of Mo{\oe}glin--Waldspurger.
For global data one can specify a formal cuspidal datum
$\mu\boxtimes\nu_2$, where $\nu_2$ is the standard two-dimensional
representation.

This notation does not construct a global Langlands
group or a map to a residual-characteristic stack. Packet assertions
must specify twists, central character and local-global realization.
In particular, \eqref{v4-arith:ic:w10-d:eq:arthur-param-s0} does not
construct an automorphic representation Steinberg at every finite place.
\end{remark}

The direct sum of the trivial and mod-\(p\) cyclotomic characters
admits several independent cyclotomic extension directions. Their dimensions, and the automorphisms
of their deformations, depend on the prime. These data determine what a
comparison with a distinguished extension class must preserve.
The constructions below give the continuous local deformation groupoids,
explicit Kummer classes, and a formal family integrating those classes.

\subsection{First-order deformations and their automorphisms}

Fix a separable closure \(K^{\mathrm{sep}}\) of \(K=\mathbb Q_p\), and
put \(\Gamma=\operatorname{Gal}(K^{\mathrm{sep}}/K)\), with its profinite topology.
Let \(k/\mathbb F_p\) be a finite field with the discrete topology.
Choose a primitive \(p\)th root \(\zeta_p\), with \(\zeta_2=-1\), and define
\[
 g(\zeta_p)=\zeta_p^{\omega(g)},\qquad
 \omega:\Gamma\longrightarrow\mathbb F_p^\times\longrightarrow k^\times.
\]
Write \(k(r)\) for the one-dimensional \(k\)-representation with character
\(\omega^r\).

In particular, \(\omega=1\) when \(p=2\).
Fix the ordered basis \(e_1,e_2\) of \(V=k^2\), and set
\[
 \rho_0(g)=\begin{pmatrix}1&0\\0&\omega(g)\end{pmatrix},\qquad
 \mathfrak a=\operatorname{End}_k(V),\qquad
 \mathfrak a_0=\ker(\operatorname{tr}:\mathfrak a\to k).
\]
Write \(E_{ij}\) for the matrix sending \(e_j\) to \(e_i\) and the other basis vector to zero.
The action on \(\mathfrak a\) is \(gX=\rho_0(g)X\rho_0(g)^{-1}\).
All cochains and representations are continuous.

For a finite discrete \(\Gamma\)-module \(N\), use the inhomogeneous cochains
\(C^n(\Gamma,N)=\operatorname{Map}_{\mathrm{cts}}(\Gamma^n,N)\), with \(C^0(\Gamma,N)=N\).
Their differential is
\begin{align*}
 (df)(g_1,\ldots,g_{n+1})
 &=g_1f(g_2,\ldots,g_{n+1})\\
 &\quad+\sum_{i=1}^n(-1)^i
 f(g_1,\ldots,g_ig_{i+1},\ldots,g_{n+1})\\
 &\quad+(-1)^{n+1}f(g_1,\ldots,g_n).
\end{align*}
Thus \(dX(g)=gX-X\) and
\(dc(g,h)=gc(h)-c(gh)+c(g)\).
Write \(Z^1(\Gamma,N)=\ker(d:C^1\to C^2)\),
\(B^1(\Gamma,N)=dN\), and \(H^i(\Gamma,N)=H^i(C^\bullet(\Gamma,N))\).

Set \(A=k[\epsilon]/(\epsilon^2)\), with the discrete topology.
A framed deformation is a continuous representation
\(\rho:\Gamma\to\operatorname{GL}_2(A)\) reducing to \(\rho_0\) in the specified basis.
Morphisms preserve this entire \(A\)-basis.
A residual-marked unframed deformation is a free rank-two \(A\)-module
with a continuous \(\Gamma\)-action and a specified identification of its
reduction with \(V\). Its morphisms induce the identity on \(V\).

\begin{proposition}[Continuous first-order groupoids]
\label{v4-arith:ic:local-cyclotomic:groupoids}
Framed deformations form the discrete groupoid on \(Z^1(\Gamma,\mathfrak a)\).
Residual-marked unframed deformations form the action groupoid
\[
 [Z^1(\Gamma,\mathfrak a)/\mathfrak a],\qquad
 X:c\longmapsto c-dX.
\]
Its isomorphism classes form \(H^1(\Gamma,\mathfrak a)\).
The automorphism group of each object is the additive group
\(H^0(\Gamma,\mathfrak a)\).
\end{proposition}

\begin{proof}
Every framed lift has a unique expression
\[
 \rho_c(g)=(1+\epsilon c(g))\rho_0(g).
\]
Multiplication gives
\[
 \rho_c(g)\rho_c(h)
 =\bigl(1+\epsilon(c(g)+gc(h))\bigr)\rho_0(gh).
\]
The homomorphism condition is therefore \(dc=0\).
A basis lifting \(e_1,e_2\) identifies any residual-marked deformation
with such a matrix representation.

Two such choices differ by
\(1+\epsilon X\), and
\[
 (1+\epsilon X)\rho_c(g)(1-\epsilon X)=\rho_{c-dX}(g).
\]
The quotient and automorphism statements follow from this equality.
An automorphism fixes \(c\) exactly when \(dX=0\).
\end{proof}

If the residual marking is forgotten, the residual centralizer also acts.
Write
\[
 C=\{B\in\operatorname{GL}_2(k):B\rho_0(g)=\rho_0(g)B\text{ for all }g\}.
\]
Every permissible change of basis is \((1+\epsilon X)B\), and its action is
\[
 c\longmapsto BcB^{-1}-dX.
\]
Thus the acting group is \(\mathfrak a\rtimes C\).

For odd \(p\), the character \(\omega\) is nontrivial and \(C\) is the group of invertible diagonal matrices.
For \(p=2\), it is \(\operatorname{GL}_2(k)\).
The stabilizer of \(c\) consists of the pairs satisfying
\(BcB^{-1}-dX=c\). At \(c=0\), this is
\(H^0(\Gamma,\mathfrak a)\rtimes C\).

A formal scheme, viewed as a groupoid-valued functor, has only
identity automorphisms at its points. It cannot represent these
unframed groupoids, which retain nonzero infinitesimal automorphisms.
This assertion concerns the groupoids, rather than representability of
their sets of isomorphism classes.

\subsection{Determinant character and determinant identification}

Fixing a determinant character restricts the objects of a deformation
problem. Choosing an identification of determinant representations also
restricts its morphisms. To separate these conditions, fix
\[
 \delta_\eta(g)=\omega(g)(1+\epsilon\eta(g)),\qquad
 \eta\in Z^1(\Gamma,k).
\]
Here the action on \(k\) is trivial, so \(\eta\) is a continuous additive
homomorphism.

For a character \(\delta:\Gamma\to A^\times\), write \(A(\delta)\) for the free rank-one module with that action.
The constant determinant lift has \(\eta=0\).
The identity
\[
 \det\rho_c(g)=\omega(g)(1+\epsilon\operatorname{tr}c(g))
\]
holds in every characteristic.

\begin{proposition}[The two fixed-determinant groupoids]
\label{v4-arith:ic:local-cyclotomic:determinant}
For the constant lift \(\delta_0=\omega\), the residual-marked deformation
groupoid with fixed determinant character is
\[
 \mathcal D_{\mathrm{char}}(A)
 =[Z^1(\Gamma,\mathfrak a_0)/\mathfrak a].
\]
Fix the residual orientation \(e_1\wedge e_2\).
If objects include an identification
\(\lambda:\det(V_A)\xrightarrow{\sim}A(\omega)\) lifting that orientation,
and arrows preserve \(\lambda\), the groupoid is
\[
 \mathcal D_{\mathrm{det}}(A)
 =[Z^1(\Gamma,\mathfrak a_0)/\mathfrak a_0].
\]
Both groupoids have isomorphism classes \(H^1(\Gamma,\mathfrak a_0)\).
Their automorphism groups are respectively \(H^0(\Gamma,\mathfrak a)\)
and \(H^0(\Gamma,\mathfrak a_0)\).
\end{proposition}

\begin{proof}
The fixed-character equation is \(\operatorname{tr}c=0\).
Every strict conjugation preserves the determinant, and
\(\operatorname{tr}(dX)=0\). This proves the first groupoid formula.
For the second, choose a lifted basis with
\(\lambda(e_1\wedge e_2)=1\).
Such a basis exists because
\[
 \det:\,1+\epsilon\mathfrak a\longrightarrow1+\epsilon k
\]
is surjective: \(\operatorname{diag}(1+\epsilon t,1)\) has determinant
\(1+\epsilon t\). Changes between these bases satisfy
\(\operatorname{tr}X=0\), which proves the second formula.

The trace admits the \(\Gamma\)-equivariant splitting
\[
 k\longrightarrow\mathfrak a,\qquad t\longmapsto tE_{11}.
\]
Indeed, \(E_{11}\) is invariant and has trace one. Hence
\(\mathfrak a=\mathfrak a_0\oplus kE_{11}\) and
\(d\mathfrak a=d\mathfrak a_0\), including when \(p=2\).
The two orbit sets are therefore \(H^1(\Gamma,\mathfrak a_0)\).
The equations \(dX=0\), with the respective trace restrictions,
give the stated automorphism groups.
\end{proof}

For a general \(\eta\), the cocycle \(c_\eta(g)=\eta(g)E_{11}\) has trace
\(\eta\). The object spaces are therefore the affine translates
\(c_\eta+Z^1(\Gamma,\mathfrak a_0)\), with the same gauge formulas.
A determinant-compatible framing requires
\(\lambda(e_1\wedge e_2)=1\), and gives the discrete groupoid on
\(Z^1(\Gamma,\mathfrak a_0)\) when \(\eta=0\).
If a full framing and an independent determinant identification are both
retained, the discrete object set is instead
\(Z^1(\Gamma,\mathfrak a_0)\times k\): the second coordinate records
\(\lambda(e_1\wedge e_2)=1+\epsilon t\).
For a chosen determinant identification, forgetting the residual marking
allows residual changes in \(C\cap\operatorname{SL}_2(k)\).

The coefficient decomposition is
\begin{align*}
 \mathfrak a&=kE_{11}\oplus kE_{22}\oplus k(-1)E_{12}\oplus k(1)E_{21},\\
 \mathfrak a_0&=k(E_{11}-E_{22})\oplus k(-1)E_{12}\oplus k(1)E_{21}.
\end{align*}
At \(p=2\), the diagonal trace-zero generator is \(I\), and every
summand has trivial action. At \(p=3\), the two off-diagonal characters
coincide, but their matrix subspaces remain distinct.

The first-order tangent groupoid is presented by the two-term complex
\[
 [\mathfrak a\xrightarrow{-d}Z^1(\Gamma,\mathfrak a)]
 \quad\text{in degrees }[-1,0].
\]
Its degree \(-1\) cohomology records automorphisms and its degree \(0\)
cohomology records deformation classes. This is the truncation through
degree zero of \(C^\bullet(\Gamma,\mathfrak a)[1]\), with the cohomological
shift convention \((C[1])^i=C^{i+1}\) and shifted differential \(-d\).
Thus \(H^2(\Gamma,\mathfrak a)\) occurs in degree one of the full shifted complex.

For a determinant identification, replace \(\mathfrak a\) by \(\mathfrak a_0\).
For the fixed-character problem alone, the two-term complex is
\([\mathfrak a\xrightarrow{-d}Z^1(\Gamma,\mathfrak a_0)]\).
Identifying any full cochain complex with the tangent complex of a
specified derived stack requires an equivalence of the underlying
moduli problems.

\subsection{Kummer classes at every prime}

The lower-left summand has an explicit arithmetic description.
For \(a\in K^\times\), choose \(\beta_a\in K^{\mathrm{sep}}\) with
\(\beta_a^p=a\), and define
\[
 \frac{g(\beta_a)}{\beta_a}=\zeta_p^{\kappa_a(g)},
 \qquad \kappa_a(g)\in\mathbb F_p\subset k.
\]
The action on roots of unity gives
\[
 \kappa_a(gh)=\kappa_a(g)+\omega(g)\kappa_a(h).
\]
This cocycle is continuous, since \(\beta_a\) and \(\zeta_p\) lie in a
finite Galois extension. Replacing \(\beta_a\) by \(\zeta_p^r\beta_a\)
changes it by the coboundary \(dr(g)=\omega(g)r-r\).

\begin{proposition}[Explicit cyclotomic cohomology]
\label{v4-arith:ic:local-cyclotomic:kummer}
The Kummer map gives an isomorphism
\[
 (\mathbb Q_p^\times/\mathbb Q_p^{\times p})\otimes_{\mathbb F_p}k
 \xrightarrow{\sim} H^1(\Gamma,k(1)),\qquad [a]\otimes1\longmapsto[\kappa_a].
\]
For odd \(p\), a basis is \([\kappa_p],[\kappa_{1+p}]\).
For \(p=2\), a basis is \([\kappa_2],[\kappa_{-1}],[\kappa_5]\).
The dimensions are respectively two and three.
\end{proposition}

\begin{proof}
Multiplication of roots makes the Kummer map additive modulo \(p\)th
powers. If \(\kappa_a=dr\), replacing \(\beta_a\) by
\(\zeta_p^{-r}\beta_a\) makes it invariant. Hence \(a\) is a \(p\)th
power in \(K\). Conversely, Hilbert 90 implies that every continuous
\(\mu_p\)-valued cocycle is \(g(\beta)/\beta\) for some
\(\beta\in(K^{\mathrm{sep}})^\times\). Then \(\beta^p\in K^\times\).
This proves the isomorphism over \(\mathbb F_p\).

The required Hilbert 90 statement reduces to a finite Galois extension
\(L/K\) through which the cocycle factors. If \(z_\sigma\) is a
multiplicative cocycle for \(\operatorname{Gal}(L/K)\), choose \(a\in L\)
so that
\[
 b=\sum_\sigma z_\sigma\sigma(a)\ne0.
\]
Such an \(a\) exists by linear independence of distinct field automorphisms.
The cocycle identity gives \(\tau(b)=z_\tau^{-1}b\).

Consequently \(z_\tau=\tau(b^{-1})/b^{-1}\), as required.
This is the proof of Milne, \emph{Class Field Theory}, Proposition~II.1.22, p.~67.
A basis of the finite extension \(k/\mathbb F_p\) decomposes each continuous
cochain into finitely many coordinate cochains. Tensoring with \(k\)
therefore commutes with cohomology.

For odd \(p\), the multiplicative decomposition is
\[
 \mathbb Q_p^\times=p^{\mathbb Z}\times\mu_{p-1}\times(1+p\mathbb Z_p).
\]
The valuation supplies the first factor. Hensel lifting of the distinct
nonzero residues supplies the roots of unity in the second factor.
The \(p\)th-power map on \(\mu_{p-1}\) is bijective.

The logarithm and exponential give inverse homomorphisms
\[
 \log:1+p\mathbb Z_p\xrightarrow{\sim}p\mathbb Z_p.
\]
For \(x,y\in p\mathbb Z_p\), the estimates
\[
 v_p(x^n/n)\ge n-v_p(n),\qquad
 v_p(y^n/n!)\ge n-v_p(n!)
\]
prove convergence of the logarithm and exponential, using
\(v_p(n!)\le n/(p-1)\). Their formal inverse
and product identities pass to these convergent series.
Logarithm carries \(p\)th powers to \(p^2\mathbb Z_p\), and
\(\log(1+p)\equiv p\pmod{p^2}\).
Thus valuation and the principal-unit quotient give the basis
\([p],[1+p]\).

At \(p=2\), each odd unit has a unique sign making it congruent to
one modulo four. Hence
\[
 \mathbb Q_2^\times=2^{\mathbb Z}\times\{1,-1\}\times(1+4\mathbb Z_2).
\]
The preceding convergence estimates, with input valuation at least two,
give \(\log(1+4\mathbb Z_2)=4\mathbb Z_2\).
Squaring becomes multiplication by two, and
\(\log(5)\equiv4\pmod8\). The valuation, sign, and principal-unit
factors give the independent classes \([2],[-1],[5]\).
\end{proof}

In particular, \(p=3\) gives the basis \([\kappa_3],[\kappa_4]\).
The choice of \(\zeta_p\) identifies \(\mu_p\) with \(\mathbb F_p(1)\).
Changing it rescales all numerical cocycles by one nonzero scalar.
The intrinsic classes take values in \(\mu_p\otimes_{\mathbb F_p}k\).

The unit classes form an arithmetic subspace, and valuation gives an
exact sequence
\[
 0\longrightarrow
 (\mathbb Z_p^\times/(\mathbb Z_p^\times)^p)\otimes_{\mathbb F_p}k
 \longrightarrow H^1(\Gamma,k(1))
 \xrightarrow{v_p}k\longrightarrow0.
\]
The unit subspace has dimension one for odd \(p\) and two for \(p=2\).
The specified uniformizer \(p\) supplies the splitting
\(1\mapsto[\kappa_p]\). The quotient itself does not specify a splitting.

There is also an unramified condition. If \(I\subset \Gamma\) is inertia,
define its subspace as
\[
 H^1_{\mathrm{ur}}(\Gamma,k(1))
 =\operatorname{im}\bigl(H^1(\Gamma/I,k(1)^I)\to H^1(\Gamma,k(1))\bigr).
\]
For odd \(p\), this is zero: \(\mathbb Q_p(\zeta_p)/\mathbb Q_p\) is
totally ramified of degree \(p-1\), since \(((1+T)^p-1)/T\) is Eisenstein at \(p\). Thus \(k(1)^I=0\).

At \(p=2\), the unramified subspace is the line \(k[\kappa_5]\).
Indeed, \((1+\sqrt5)/2\) satisfies \(x^2-x-1=0\), whose reduction
is separable and irreducible over \(\mathbb F_2\).
It therefore generates the unramified quadratic extension.
Since \(\Gamma/I\cong\widehat{\mathbb Z}\), its continuous additive
characters into \(k\) form a one-dimensional \(k\)-space.

\subsection{An explicit formal family}

Each cyclotomic cocycle defines both an extension of residual
representations and a tangent direction at their split sum.
For \(b\in Z^1(\Gamma,k(1))\), set
\[
 r_b(g)=\begin{pmatrix}1&0\\b(g)&\omega(g)\end{pmatrix}.
\]
The line \(ke_2\) is invariant, and the quotient is trivial. Thus
\(r_b\) gives an exact sequence
\[
 0\longrightarrow k(1)\longrightarrow E_b\longrightarrow k\longrightarrow0.
\]

Conversely, a vector-space splitting of any such extension gives this
matrix form. Changing the splitting changes \(b\) by a coboundary.
Consequently
\[
 \operatorname{Ext}^1_{\operatorname{Rep}^{\mathrm{cts}}_k(\Gamma)}(k,k(1))
 \cong H^1(\Gamma,k(1)),
\]
where the category consists of finite-dimensional discrete continuous
\(k\)-representations. The order of the two endpoint characters is
part of this identification.

The coefficient inclusion \(u\mapsto uE_{21}\) and projection onto
the lower-left entry commute with the cochain differential. They induce
an injection with a specified left inverse
\[
 j:H^1(\Gamma,k(1))\lhook\joinrel\longrightarrow H^1(\Gamma,\mathfrak a_0).
\]
The class \(j([b])\) is represented by
\(\rho_{\epsilon b}(g)=\left(\begin{smallmatrix}1&0\\\epsilon b(g)&\omega(g)\end{smallmatrix}\right)\).

\begin{theorem}[Formal cyclotomic family]
\label{v4-arith:ic:local-cyclotomic:formal-family}
Set \(d=2\) and \((a_1,a_2)=(p,1+p)\) for odd \(p\).
For \(p=2\), set \(d=3\) and \((a_1,a_2,a_3)=(2,-1,5)\).
Choose the Kummer cocycles above and set
\[
 R_{\mathrm{gal}}=k[[t_1,\ldots,t_d]],\qquad
 u(g)=\sum_{i=1}^d t_i\kappa_{a_i}(g).
\]
With the maximal-ideal topology on \(R_{\mathrm{gal}}\), the matrices
\[
 \rho_{R_{\mathrm{gal}}}(g)=\begin{pmatrix}1&0\\u(g)&\omega(g)\end{pmatrix}
\]
define a continuous representation reducing to \(\rho_0\), with
determinant \(\omega\) and the standard determinant identification.
Its tangent map is injective, with image \(j(H^1(\Gamma,k(1)))\).
\end{theorem}

\begin{proof}
The cocycle equation for \(u\) is exactly the lower-left entry of
\(\rho_{R_{\mathrm{gal}}}(gh)=\rho_{R_{\mathrm{gal}}}(g)\rho_{R_{\mathrm{gal}}}(h)\).
The finitely many cocycles factor through a common finite Galois quotient,
so the representation is continuous for the maximal-ideal topology.
Its determinant and reduction follow from the diagonal entries.

Give each Artinian local \(k\)-algebra \(B\) with residue field \(k\)
the discrete topology. Use the same residual-marked deformation
definition on free rank-two \(B\)-modules. A continuous local
homomorphism \(R_{\mathrm{gal}}\to B\) specializes these matrices
to a deformation. This is compatible with homomorphisms between such
algebras and therefore defines a map from the functor of points of
\(\operatorname{Spf}R_{\mathrm{gal}}\) to the determinant-identified deformation groupoid.

A tangent homomorphism sends \(t_i\) to \(\epsilon x_i\).
Its class is \(j(\sum_i x_i[\kappa_{a_i}])\).
The Kummer basis and the left inverse of \(j\) prove injectivity.
\end{proof}

Changing \(\beta_{a_i}\) to \(\zeta_p^{r_i}\beta_{a_i}\) adds
\(dr_i\) to the corresponding cocycle. Set \(s=\sum_i t_i r_i\).
Conjugation by \(1-sE_{21}\) carries the old formal family to the new
one, since its lower-left entry changes by \(ds\).
This calculation holds over \(R_{\mathrm{gal}}\), without a first-order truncation.
The formal family therefore changes by an explicit isomorphism under
these root choices.

This family fixes both diagonal characters and a line and quotient.
It is not asserted to represent the full deformation problem.
At \(p=2\), the underlying residual representation is trivial, and its
unmarked form has no preferred decomposition into two lines.
The inclusion through \(E_{21}\) uses the stated frame or flag.
For odd \(p\), the two distinct character subspaces are intrinsic,
although their scalar coordinates still depend on the chosen bases.

\subsection{Dimensions and comparison maps}

The Kummer calculation determines one direct summand without local
duality. The full dimension table also uses finite-module local duality
and the local Euler characteristic theorem.
These are Milne, \emph{Arithmetic Duality Theorems},
Corollary~I.2.3, p.~28, and Theorem~I.2.8, p.~31.
For a finite-dimensional continuous \(k[\Gamma]\)-module \(N\), they imply
\[
 h^2(N)=h^0(N^\vee(1)),\qquad
 h^0(N)-h^1(N)+h^2(N)=-\dim_kN,
\]
where \(h^i(N)=\dim_kH^i(\Gamma,N)\) and
\(N^\vee=\operatorname{Hom}_k(N,k)\).

The field trace pairing identifies the underlying finite-module dual
with the additive module of \(N^\vee(1)\).
The Euler formula follows by substituting
\(|N|=p^{[k:\mathbb F_p]\dim_kN}\) into
\(\chi(\Gamma,N)=(\mathbb Z_p:|N|\mathbb Z_p)^{-1}\).
Both theorems apply at \(p=2\), since the base field has characteristic zero.
For \(N=k(r)\), they give
\[
 h^1(k(r))=1+\mathbf1_{\omega^r=1}+\mathbf1_{\omega^{1-r}=1}.
\]

The matrix decomposition and
\(\dim B^1(\Gamma,N)=\dim N-h^0(N)\) now give
\[
\begin{array}{c|ccc}
 &p\ge5&p=3&p=2\\\hline
 h^1(k)&2&2&3\\
 h^1(k(1))&2&2&3\\
 h^1(k(-1))&1&2&3\\
 h^0(\mathfrak a)&2&2&4\\
 h^0(\mathfrak a_0)&1&1&3\\
 h^1(\mathfrak a)&7&8&12\\
 h^1(\mathfrak a_0)&5&6&9\\
 \dim Z^1(\Gamma,\mathfrak a)&9&10&12\\
 \dim Z^1(\Gamma,\mathfrak a_0)&7&8&9\\
 h^2(\mathfrak a)&1&2&4\\
 h^2(\mathfrak a_0)&1&2&3
\end{array}
\]
These dimensions concern the unrestricted local continuous problem.
A crystalline, finite-flat, ordinary, or global Selmer condition requires
its own definition before imposing a further subspace.

The explicit formal family supplies a local Galois target for a
comparison. A specified nonzero Kummer class selects a one-parameter
subfamily. The unit condition selects a line for odd \(p\) and a plane
at \(p=2\). To identify an analytic residue line with one of these lines,
construct a map carrying the residue to the specified nonzero cohomology class.

A coherent extension requires further categorical data.
Suppose an exact \(k\)-linear functor from
\(\operatorname{Rep}^{\mathrm{cts}}_k(\Gamma)\) to an abelian category
has identified images of \(k\) and \(k(1)\).
It sends \(E_b\) to an extension with those endpoints. An exact functor between specified
stable categories instead sends its associated triangle to a triangle.
Either construction gives a map on extension classes.

An isomorphism
requires a calculation of that map, including its kernel and image.
A tangent vector alone does not define two coherent objects or a functor
between their categories. On \(\operatorname{Spec}k[t]\), for instance,
the tangent space at zero is \(k\), whereas
\(\operatorname{Ext}^1(k[t],k[t])=0\). For the residue module \(k_0\),
the resolution
\[
 0\longrightarrow k[t]\xrightarrow{t}k[t]\longrightarrow k_0\longrightarrow0
\]
instead gives \(\operatorname{Ext}^1(k_0,k_0)=k\).
Thus the endpoints of a coherent extension are indispensable.

A comparison from a complex representation category also requires
arithmetic coefficient data. There is no unital field map
\(k\to\mathbb C\), since it would carry \(p\cdot1_k=0\) to
\(p\cdot1_{\mathbb C}\ne0\). There is no unital map in the reverse
direction either, since \(p\) is invertible in \(\mathbb C\).
Reduction requires an integral coefficient ring and an integral
object or stable lattice. A residue-field map to \(k\) must then
induce the required map on extension classes. A global Galois comparison
additionally requires a ramification set, local conditions, and the
restriction maps into the local cohomology computed here.

Finally, scalar twisting gives
\[
 (1\oplus\omega)\otimes\omega^{-1}=\omega^{-1}\oplus1.
\]
This is \(1\oplus1\) only at \(p=2\). At \(p=3\), exchanging summands
identifies it with the original representation.
For any scalar character \(\chi\), conjugation gives
\(\operatorname{Ad}(\rho_0\otimes\chi)=\operatorname{Ad}(\rho_0)\),
while the determinant changes by \(\chi^2\).
Hence a scalar twist need not preserve a fixed-determinant fiber, and
no extra determinant twist enters the adjoint cochain formulas.


\subsection{The formal stack and its framed presentation}
\label{v4-arith:ic:w10-d:prop:EG-closed-point-s1}
The local deformation groupoids above retain the residual representation
and its automorphisms. At the split cyclotomic point, the
residual-marked determinant-identified groupoid has
\begin{equation}\label{v4-arith:ic:w10-d:eq:EG-tangent-s1}
 \pi_0\mathcal D_{\mathrm{det}}(k[\epsilon]/(\epsilon^2))
       =H^1_{\mathrm{cts}}(\Gamma,\operatorname{ad}^0(1\oplus\omega)).
\end{equation}
Here $\pi_0$ means the set of isomorphism classes, with its vector-space
structure from the cocycle description. Fixing only the determinant
character gives the same isomorphism classes at this split point,
but retains the larger automorphism group $H^0(\Gamma,\mathfrak a)$.
This distinction follows from the equivariant trace splitting through
$E_{11}$, including in characteristic two.

For the Emerton--Gee moduli problem, a framed lifting ring supplies a
versal presentation of the formal deformation groupoid.
Emerton--Gee, Proposition~3.6.3, p.~76, identifies its self-fiber product
using the change-of-basis group completed along the residual centralizer.
The pointed formal problem instead uses completion along the identity.
It does not replace the unframed groupoid by a formal scheme with
no automorphisms. A derived enhancement, wild-inertia conditions and
an identification with its full tangent complex require the corresponding
moduli comparison. The computed local complex retains this requirement.

\subsection{Tangent complexes and spectral extension classes}
\label{v4-arith:ic:w10-d:prop:spectral-ext1-s1}
A tangent direction describes deformation of one point. An extension
between two sheaves requires two objects and a category of morphisms.
To express the proposed comparison, suppose a derived stack
$\mathcal X$ over $k$ has been specified, with a cotangent complex
$\mathbb L_{\mathcal X}$ admitting the dual used below. At a point $x$,
set $T_x=R\operatorname{Hom}_k(x^*\mathbb L_{\mathcal X},k)$.
Its degree-zero cohomology is the first-order tangent space,
while degree minus one records infinitesimal automorphisms.
Write $\mathbb T^1_x=H^0(T_x)$ for this convention.

Suppose in addition that a spectral category
$\mathcal A_{\mathrm{sp}}$ has been constructed, with objects
$\mathcal O_{\mathrm{St}}$ and $\mathcal O_{\mathbf1}$ and specified
support. These names alone do not define sheaves.
For example, an ind-coherent category with nilpotent singular support
requires a derived stack for which its singularity space and the
specified nilpotent subset are defined. The desired comparison is
\begin{equation}\label{v4-arith:ic:w10-d:eq:spectral-ext1-s1}
 L_{\mathrm{cyc}}\longrightarrow
 \operatorname{Ext}^1_{\mathcal A_{\mathrm{sp}}}
          (\mathcal O_{\mathrm{St}},\mathcal O_{\mathbf1}),
 \qquad L_{\mathrm{cyc}}\subseteq\mathbb T^1_{x_{\mathrm{cyc}}}.
\end{equation}

Here the subspace $L_{\mathrm{cyc}}$, its dimension and the displayed
map are constructions still required. A cotangent calculation does
not identify this Ext group between two different sheaves.
\begin{remark}[The reverse spectral comparison]
\label{v4-arith:ic:w10-d:cor:spectral-ext1-s0}
For a point $x_{\mathrm{triv}}$ representing $1\oplus1$, the opposite
orientation asks for a specified subspace
$L_{\mathrm{triv}}\subseteq\mathbb T^1_{x_{\mathrm{triv}}}$ and a map
\[
 L_{\mathrm{triv}}\longrightarrow
 \operatorname{Ext}^1_{\mathcal A_{\mathrm{sp}}}
       (\mathcal O_{\mathbf1},\mathcal O_{\mathrm{St}}).
\]
The two endpoints and their support must be constructed in this category.
The explicit scalar-twist calculation above does not identify this
problem with the cyclotomic problem at odd primes. At $p=2$, equality
of the residual representations still does not choose the two sheaves
or the map. An Arthur parameter over $\mathbb C$ supplies neither
spectral comparison by itself.
\end{remark}


\section{Blocks, extensions and geometric support}
\label{v4-arith:ic:blocks}

A local categorical comparison must specify its block algebra and
its spectral support. The ordinary center, the full endomorphism
algebra of a generator and the deformation ring are different
objects. Their distinctions persist even when every simple object
has the expected parameter.

\subsection{Characteristic-zero Bernstein decomposition}

\begin{theorem}[Bernstein--Deligne]
\label{v4-arith:ic:w5-f:F1a3-AB:thm:bernstein}
 Let \(G\) be a connected reductive
\(p\)-adic group and fix an algebraically closed coefficient
field \(k\) of characteristic zero. Let \(\mathrm{Mod}^{\mathrm{sm}}_G\) denote its smooth
representations over \(k\). This category decomposes as a product:
\begin{equation}
\label{v4-arith:ic:w5-f:F1a3-AB:eq:bernstein-decomposition}
\mathrm{Mod}^{\mathrm{sm}}_{G}\;=\;\prod_{\mathfrak{s}\in\mathfrak{B}(G)}
\mathrm{Block}_{\mathfrak{s}},
\end{equation}
Here \(\mathfrak{B}(G)\) is the set of inertial equivalence classes
of supercuspidal supports. The decomposition is realized by a
central idempotent decomposition of the Bernstein center
\(Z(\mathrm{Mod}^{\mathrm{sm}}_{G})=\prod_{\mathfrak{s}}Z_{\mathfrak{s}}\)
into a product of integral domains.
\end{theorem}

\begin{proof}[Attribution]
Bernstein--Deligne (1984), in the proceedings
\emph{Representations of reductive groups over a local field}.
For the block decomposition and torsor of unramified twists modulo
self-twists, see also Helm, arXiv:1201.1874, Theorems~3.4--3.5,
pp.~7--8.
\end{proof}

\begin{proposition}[central-character fibers of Bernstein blocks]
\label{v4-arith:ic:w5-f:F1a3-AB:prop:block-chi-compat}

Let \(k\) be algebraically closed of characteristic
zero, and let \(\chi:Z(G)\to k^\times\) be a smooth
central character. For a Bernstein block
\(\mathcal B_{\mathfrak s}\), put
\[
 \mathcal B_{\mathfrak s}^{\chi}
 =\{M\in\mathcal B_{\mathfrak s}:
        z|_M=\chi(z)\,\mathrm{id}_M
        \text{ for every }z\in Z(G)\}.
\]
Then
\[
 \operatorname{Rep}^{\mathrm{sm},\chi}_k(G)
   \simeq \prod_{\mathfrak s}
                     \mathcal B_{\mathfrak s}^{\chi},
\]
with zero factors omitted.

Let \(\delta_z\in Z_{\mathfrak s}\) denote the
central operation supplied by \(z\in Z(G)\), and
let
\[
 I_{\mathfrak s,\chi}
   =(\delta_z-\chi(z):z\in Z(G))
       \subseteq Z_{\mathfrak s}.
\]
The ring \(Z_{\mathfrak s}/I_{\mathfrak s,\chi}\)
acts on \(\mathcal B_{\mathfrak s}^{\chi}\).
This is a fiber condition inside the block.
\end{proposition}

\begin{proof}
Each Bernstein projector commutes with the action
of \(Z(G)\). Thus every component of an object
with central character \(\chi\) again has that
central character. Conversely, a family of block objects reconstructs a smooth
representation by its direct sum, and the sum satisfies the same
scalar equations. This proves the
categorical decomposition.

The generators of \(I_{\mathfrak s,\chi}\)
act as zero on the indicated subcategory, so
the central action factors through the quotient.
No claim that taking the categorical center
commutes with this quotient is needed.
\end{proof}

For $G=GL_2(F)$, fix a uniformizer $\varpi$ and normalized
parabolic induction. In the unramified principal-series block write
$a=\eta_1(\varpi)$, $b=\eta_2(\varpi)$ for its two unramified
inducing characters, so $t=a+b$ and $d=ab$. The modulus character
in normalized induction is trivial on scalar matrices. Consequently
the central operation of $u\in\mathcal O_F^\times$ is $1$, and that
of $\varpi$ is $d$, in the ring $k[t,d,d^{-1}]$.

These equalities
hold as central operations: their values on all supercuspidal supports
are respectively $1$ and $ab$, and the block ring is reduced.
If $\chi$ is unramified, put $c=\chi(\varpi)\in k^\times$.
Then $I_{\mathfrak s,\chi}=(d-c)$ and the fiber has the quotient
scalar action $k[t,d,d^{-1}]/(d-c)$. It is nonzero, for instance on
any irreducible constituent of the normalized induction of the
unramified characters with values $1,c$ at $\varpi$.

If $\chi$ is ramified, choose $u\in\mathcal O_F^\times$ with
$\chi(u)\ne1$. The relation $1-\chi(u)$ is an invertible scalar,
so the quotient and the fixed-character fiber category are zero.
Thus fixing central character neither selects a whole Bernstein
block nor gives a nonzero fiber in every block.


The exact Bernstein projectors extend degreewise to ordinary
complexes of smooth characteristic-zero representations. They
preserve quasi-isomorphisms because they are exact, and therefore
act on the derived category.
This elementary extension does not identify that category with a
separately defined geometric Hecke category.

For \(GL_2(F)\), take algebraically closed characteristic-zero
\(\ell\)-adic coefficients with \(\ell\ne p\).
Use compatible reciprocity and induction normalizations.
The ordinary Bernstein center then identifies with scalar functions
on the continuous Weil parameter stack.
This includes residue primes two and three.

The four inertial cases retain different parameter rings:
irreducible inertia, two fixed characters, an interchanged pair,
and scalar inertia. They cannot be recovered from a Frobenius
matrix with its inertia action discarded.
The ring construction and its full coefficient hypotheses are in
Proposition~\ref{v4-arith:w10-e:prop:excursion-bernstein}.

That scalar image need not exhaust the center of the spectral
category. Scalar \(GL_2\) acts trivially on the parameter coordinates,
but equivariant coherent objects retain central weights.
After Ind-completion, projections to these weight factors give
central operations. A determinant twist changes weight by two.

The weight-zero projector separates this operation from scalar
multiplication. See
Proposition~\ref{v4-arith:w10-e:prop:scalar-centre-obstruction}.
Choosing one weight removes this particular obstruction, without
proving a categorical equivalence.

\subsection{The coefficient-prime block algebra}

Use the coefficient and reciprocity conventions of
Pa\v{s}k\={u}nas, \emph{The image of Colmez's Montreal functor}.
Here \(p\ge5\), \(G=GL_2(\mathbb Q_p)\), and
\(\mathcal O\) is a finite extension's integer ring with residue
field \(k\). Fix a full central character \(\zeta\).
Let \(\mathcal A_\zeta^{\mathrm{lfin}}\) be the smooth
\(\mathcal O\)-torsion representations which are locally of finite
length: each vector generates a finite-length \(\mathcal O[G]\)-module.

A block \(\mathcal B\) is a connected component of the extension
relation on irreducible objects, allowing extensions in both
directions. Let \(J_{\mathcal B}\) be the injective envelope
of the sum of its irreducibles. Its Pontryagin dual
\(P_{\mathcal B}=J_{\mathcal B}^{\vee}\) is a projective
generator of the corresponding compact category. Put
\[
 E_{\mathcal B}=\operatorname{End}(P_{\mathcal B})
    =\operatorname{End}_G(J_{\mathcal B})^{\mathrm{op}}.
\]

Endomorphisms use ordinary composition. This is a pseudocompact
ring, with topology given by its finite-length quotients.
The functor \(\tau\mapsto
\operatorname{Hom}(P_{\mathcal B},\tau^\vee)\) is an
anti-equivalence from the block to compact right
\(E_{\mathcal B}\)-modules. The right action is precomposition.
These compact modules and the locally finite torsion category
are different from the finite-length admissible Banach category.

After a coefficient extension splitting the irreducibles, the
four block forms in this theorem are as follows. The first is a
supersingular singleton. The second consists of
\[
 \operatorname{Ind}_B^G(\chi_1\otimes\chi_2\omega^{-1}),
 \qquad
 \operatorname{Ind}_B^G(\chi_2\otimes\chi_1\omega^{-1}),
 \qquad \chi_1\chi_2^{-1}\notin\{1,\omega,\omega^{-1}\}.
\]
The third is the singleton
\(\operatorname{Ind}_B^G(\chi\otimes\chi\omega^{-1})\).

The fourth consists of
\[
 \eta\circ\det,\qquad
 \mathrm{Sp}\otimes(\eta\circ\det),\qquad
 \operatorname{Ind}_B^G(\omega\otimes\omega^{-1})
                    \otimes(\eta\circ\det).
\]
Here \(\omega\) is the residual cyclotomic character under the
paper's class-field convention.
The associated semisimple Galois parameters are, respectively,
absolutely irreducible, \(\chi_1\oplus\chi_2\),
\(\chi\oplus\chi\), and \(\eta\oplus\eta\omega\).
In particular, the last is not scalar when \(p\ge5\).

\begin{theorem}[Pa\v{s}k\={u}nas, Theorems 1.5--1.7]
\label{v4-arith:ic:paskunas}
For these blocks the center of \(E_{\mathcal B}\) is the
fixed-determinant pseudodeformation ring
\(R_{\operatorname{tr}\bar\rho}^{\mathrm{ps},\zeta\varepsilon}\).
The algebra \(E_{\mathcal B}\) is finite over its center and,
on the stated complement of the reducible locus, is a matrix algebra.

For a supersingular singleton, the full algebra is
\(R_{\bar\rho}^{\zeta\varepsilon}\).
For the generic reducible block, let \(\rho_1,\rho_2\) be the
two nonsplit extensions of its distinct characters in opposite
directions. Their universal fixed-determinant deformations satisfy
\[
 E_{\mathcal B}\simeq
 \operatorname{End}_{G_{\mathbb Q_p}}
             (\rho_1^{\mathrm{un}}\oplus\rho_2^{\mathrm{un}}).
\]
The corresponding dualized and twisted Montreal images of
\(J_{\mathcal B}\) are the indicated universal representation
or sum of the two universal representations.
\end{theorem}

The compact-module convention is that of the paper's \S2.
Pontryagin duality reverses the endomorphism ring. Keeping this
opposite is necessary when comparing the noncommutative block algebras.

The determinant is \(\zeta\varepsilon\) in this theorem.
The DEG convention above uses \(\zeta\varepsilon^{-1}\) and its
specified kernel twist. A comparison must retain the dualizations,
twists and reciprocity conventions which relate the two functors.

The same paper's Theorem 1.4 gives a different positive equivalence.
Fix an absolutely irreducible admissible unitary non-ordinary
Banach representation \(\Pi\). Take the finite-length block with
every irreducible subquotient isomorphic to \(\Pi\).
It is equivalent to finite-length modules over the fixed-determinant
deformation ring of its associated characteristic-zero Galois representation.
It does not identify every torsion block with one commutative
deformation ring.

For a finite sum \(P=P_1\oplus P_2\) of nonzero projective covers
of distinct simples, \(\operatorname{End}(P)\) contains the
orthogonal projection idempotents. It cannot be a commutative local
ring. A dimension or regularity calculation for its center therefore
does not give that calculation for its full endomorphism algebra.

\subsection{Generic localization and first-order classes}

\begin{proposition}[extension vanishing over the fraction field]
\label{v4-arith:ic:w5-f:F1a3-AB:prop:Ext1-generic-p5}
Let \(R\) be a commutative Noetherian domain, let
\(K=\operatorname{Frac}(R)\), and let \(M,N\) be finite \(R\)-modules.
Then
\[
 \operatorname{Ext}^1_R(M,N)\otimes_R K
 \simeq\operatorname{Ext}^1_K(M\otimes_R K,N\otimes_R K)=0.
\]
In particular, the statement applies to a regular local domain.
It asserts vanishing at the generic point, not at every regular point.
\end{proposition}
\begin{proof}
Noetherianity gives a free resolution of \(M\) with finitely generated
terms: take a finite set of generators for each successive kernel.
Localization is exact, and homomorphisms from each finite free term
commute with localization. Localizing the resulting Hom complex
therefore computes the displayed extension group over \(K\).
Every short exact sequence of vector spaces over \(K\) splits.
\end{proof}

\begin{corollary}[arithmetic transport on a specified generic fiber]
\label{v4-arith:ic:w5-f:F1a3-AB:cor:Ext1-zero-generic-locus}
Suppose an arithmetic block category is exactly equivalent
to finite modules over a commutative Noetherian domain \(R\).
Assume its specified generic-fiber category corresponds to finite
\(K=\operatorname{Frac}(R)\)-modules.
Then generic-fiber extensions vanish. If a spectral locus is also
identified with this localization by compatible functors, the same
vanishing holds in its corresponding category.
\end{corollary}
\begin{proof}
An exact equivalence identifies extension classes. Apply
Proposition~\ref{v4-arith:ic:w5-f:F1a3-AB:prop:Ext1-generic-p5}
and the stated localization comparison. Neither regularity alone
nor flatness alone constructs the arithmetic or spectral equivalence.
\end{proof}

\subsection{Local self-extensions}

\begin{proposition}[Local self-extensions and conditional arithmetic transport]
\label{v4-arith:ic:w5-f:F1a3-AB:prop:Ext1-ordinary-controlled}
For a commutative Noetherian local ring \((A,\mathfrak m,k)\),
there is a natural isomorphism
\[
 \operatorname{Ext}^1_A(k,k)\cong
 \operatorname{Hom}_k(\mathfrak m/\mathfrak m^2,k).
\]
Take distinct simples \(S,S'\) in an arithmetic ordinary block,
and put \(A=A_\rho\). Assume a specified \(k\)-linear map identifies
\(\operatorname{Ext}^1_{\mathrm{Block}_\rho}(S,S')\) with a subspace
\(D_\rho\) of the displayed dual cotangent space.
Then its dimension is \(\dim_k D_\rho\). If in addition
\(\dim_k\mathfrak m_\rho/\mathfrak m_\rho^2\leq3\) is established,
this gives the bound \(\dim_k D_\rho\leq3\).
Breuil--Mezard multiplicities alone supply neither this class
comparison nor the bound on the embedding dimension.
\end{proposition}
\begin{proof}
Apply \(\operatorname{Hom}_A(-,k)\) to
\(0\to\mathfrak m\to A\to k\to0\).
The restriction \(\operatorname{Hom}_A(A,k)\to
\operatorname{Hom}_A(\mathfrak m,k)\) is zero, since
\(\mathfrak m k=0\), and \(\operatorname{Ext}^1_A(A,k)=0\).
Thus the connecting map identifies \(\operatorname{Hom}_A(\mathfrak m,k)\)
with \(\operatorname{Ext}^1_A(k,k)\). Every such homomorphism kills
\(\mathfrak m^2\), giving the displayed dual cotangent space.
Concretely, a functional gives the extension obtained by pushing out
\(0\to\mathfrak m\to A\to k\to0\) along \(\mathfrak m\to k\).

The arithmetic conclusion follows from the separately assumed class
comparison. There is no subtraction by a number of defining relations:
for \(A=k[\epsilon]/(\epsilon^2)\), the embedding dimension is one
and \(\operatorname{Ext}^1_A(k,k)=k\). This space is annihilated by
\(\mathfrak m\), and is not a nonzero free \(A/\mathfrak m^2\)-module.
Every simple module over a commutative local ring is its residue field.
Comparison with distinct arithmetic simples therefore requires the
specified additional map.
\end{proof}


The reduced total-quotient version also has a precise domain.
For a reduced Noetherian ring \(R\), its total quotient ring is
a finite product of fields. Localization of finite-module Ext
therefore gives vanishing of every positive extension degree
over that semisimple ring.
Reducedness cannot be replaced by the hypersurface condition.
For
\[
 R=k[[t,u,\epsilon]]/(\epsilon^2),
\]
the generic coefficient algebra is
\(K[\epsilon]/(\epsilon^2)\), with
\(K=\operatorname{Frac}(k[[t,u]])\).
Its simple module has a periodic free resolution with differential
\(\epsilon\), so \(\operatorname{Ext}^1(K,K)=K\).

For a complete local algebra over a discrete valuation ring
with uniformizer \(\varpi\), the relative first-order tangent is
\[
 \operatorname{Hom}_k
 \bigl(\mathfrak m/(\mathfrak m^2+\varpi R),k\bigr).
\]
The absolute self-extension space uses
\(\mathfrak m/\mathfrak m^2\).
For \(R=\mathcal O[[x]]\), these dimensions are one and two,
respectively. They coincide for a coefficient problem over \(k\),
not for every mixed-characteristic deformation problem.

A comparison with an extension between two arithmetic simples
must specify both endpoints, the functor or class map, and its
normalization. A commutative local ring has only its residue field
as a simple module. It does not label several arithmetic simple
pairs by itself.
Matching one extension dimension gives no higher-product comparison.
The polynomial and dual-number examples in
Section~\ref{v4-arith:ic:recognition} agree in degree one and differ
in every higher degree.

\subsection{Support and compact generation}

The localizing subcategory generated by \(k=k[t]/(t)\) in
\(D(k[t])\) disappears after inverting \(t\).
The object \(k[t]\) survives. Thus a family of image objects
can have correct pointwise descriptions and still fail to generate
its intended target.

For a quasi-smooth derived scheme \(X\), singular support lies
in its scheme of singularities, built from
\(H^1(T_X)\) over \(X^{\mathrm{cl}}\).
Its action is the degree-two Hochschild action.
It is not a locus inside \(X\) defined merely by naming a
degree-one Galois class. On derived Artin stacks, the automorphism
directions must also be retained when specifying cotangent amplitude
and the smooth-local definition of support.

\begin{definition}[Residual gerbe]\label{v4-arith:ic:residual-gerbe}
Let \(x\) be a point of a locally Noetherian algebraic stack \(X\).
Its \emph{residual gerbe} \(\mathcal G_x\) is the reduced locally
Noetherian substack whose underlying space consists of \(x\) alone.
The inclusion is a monomorphism.
This characterizes the gerbe up to its unique equivalence inside \(X\).
Existence and uniqueness are The Stacks Project,
Tags~06MT and~0H22.
\end{definition}

The stabilizers of rank-two parameters give concrete tests.
An irreducible representation has scalar centralizer.
A distinct split pair has a two-dimensional torus.
A scalar representation has \(GL_2\) centralizer.

For a nonsplit self-extension with nilpotent matrix \(N\),
the centralizer is \(\{aI+bN:a\ne0\}\). It includes scalar
automorphisms and a unipotent subgroup. A rational representative
neutralizes its residual gerbe. An arbitrary gerbe need not have
such a chosen representative.
Determinant twists shift scalar weights by two and do not generate
all weights from one of them.

A finite-stratification argument can prove compact generation when
the support-compatible recollement functors are supplied.
In particular, an extension functor preserves compact objects if
its right adjoint preserves filtered colimits.
Apply this fact at each stratum, then use the recollement triangles
to generate the ambient category.
Per-stratum compactness alone does not prove the needed compactness
after extension.

At residue primes two and three, a comparison requires the actual
block algebra and an enhanced functor into the specified spectral support category.
Full derived endomorphisms and target generation must then be compared.
One simple object does not imply semisimplicity.
A type, an irreducible representation, an inertial class, a
deformation point and an irreducible geometric component are
different objects. None is substituted for another in this criterion.



\subsection{A nilpotent local model}
\label{v4-arith:ic:w6-g:agdet-transition}

\begin{proposition}[the nilpotent cone in rank two]
\label{v4-arith:ic:w6-g:prop:codim-2}
Let $K$ be an algebraically closed field of characteristic zero.
The nilpotent cone in $\mathfrak{sl}_2(K)$ has dimension two.
Its nonzero locus is the regular nilpotent orbit, of codimension
one in $\mathfrak{sl}_2(K)$. The zero orbit has dimension zero.
\end{proposition}

\begin{proof}
Write
\[
 X=\begin{pmatrix}a&b\\c&-a\end{pmatrix}.
\]
Then $X^2=(a^2+bc)1$, so nilpotency is the equation
$a^2+bc=0$ in $\mathbb A_K^3$.
On $b\ne0$ it is parametrized by $a,b$, with $c=-a^2/b$.
The quadric is irreducible: a factorization into linear forms would
give a quadratic form of rank at most two, whereas $a^2+bc$ has
rank three.

Thus the cone has dimension two.
Every nonzero nilpotent matrix has one Jordan block and is
conjugate to $\begin{psmallmatrix}0&1\\0&0\end{psmallmatrix}$.
A conjugating matrix can be rescaled to determinant one because
$K$ is algebraically closed. The nonzero locus is therefore one
$SL_2(K)$-orbit of dimension two. The remaining orbit is $\{0\}$.
\end{proof}

\begin{remark}[the arithmetic local model]
\label{v4-arith:ic:w6-g:rem:codim-2-role}
This calculation does not determine the dimension of an arithmetic
Selmer space. Such a conclusion requires a constructed local-model
morphism and its relative dimension. A nonempty transition space
cannot be a subspace of both a closed locus and its open complement,
since their intersection is empty.
In singular-support theory the support condition lies in the scheme
of singularities of a quasi-smooth derived space.

It is not, by
definition, a locus where the underlying representations have
unipotent monodromy. The relevant constructions are
Arinkin--Gaitsgory, arXiv:1201.6343, \S\S2.3, 4.1 and 8.2.
\S11.1 defines the global nilpotent cone for the geometric
local-system stack. An arithmetic analogue requires its own
space, singular-support subset and functors.
\end{remark}

\subsection{Filtration length and geometric codimension}

\begin{example}[a Nygaard filtration on a point]
\label{v4-arith:ic:w6-g:thm:bl-2-convergence}
On the prism $(\mathbb Z_p,(p))$, take the Frobenius lift
$\varphi=\mathrm{id}$. The Nygaard submodules defined by
$\varphi(x)\in p^i\mathbb Z_p$ are
\[
 N^{\ge i}\mathbb Z_p=p^i\mathbb Z_p\quad(i\ge0),\qquad
 \operatorname{gr}_N^i\mathbb Z_p=\mathbb F_p.
\]
Every nonnegative graded degree is nonzero, although
$\operatorname{Spec}\mathbb F_p$ has dimension zero.
Thus a geometric codimension bound does not by itself truncate
the Nygaard filtration in degree two.
\end{example}

For a filtered complex $C$ with decreasing filtration $F$, the
usual cohomological spectral sequence begins with
\begin{equation}
\label{v4-arith:ic:w6-g:eq:bl-nygaard-ss}
 E_1^{a,b}=H^{a+b}(\operatorname{gr}_F^a C).
\end{equation}
Convergence of Equation~\eqref{v4-arith:ic:w6-g:eq:bl-nygaard-ss}
to $H^{a+b}(C)$ requires appropriate completeness,
exhaustiveness and convergence hypotheses on the filtered complex.
A comparison with a cellular filtration additionally requires a
filtered chain map and a proof of its effect on graded pieces.
The dimension calculation above supplies neither a map nor a
spectral-sequence degeneration statement.


\subsection{Lifting cycles and automorphic hypotheses}

Fixed-type cycle theorems supply positive geometric information
without identifying categorical generators.
In Emerton--Gee's cocharacter notation, the minimal regular type
is \(\mathrm{BT}=(1,0)\), corresponding to filtration degrees
\(\{-1,0\}\) above.
For \(GL_2(K)\), with \(K/\mathbb Q_p\) finite and \(p>2\),
Theorem 8.6.2 gives the geometric Breuil--M\'ezard formula
for potentially Barsotti--Tate types and arbitrary inertial type.

Write \(\mathcal X_2^{\mathbf k}\) for the component indexed by
the Serre weight \(\mathbf k\).
Call \(\mathbf k\) Steinberg when
\(k_{\sigma,1}-k_{\sigma,2}=p-1\) at every embedding.
In that case put
\(\widetilde k_{\sigma,1}=\widetilde k_{\sigma,2}=k_{\sigma,2}\).
The cycles in that formula are
\[
 \mathcal Z_{\mathbf k}=
 \begin{cases}
  [\mathcal X_2^{\mathbf k}],&
          \mathbf k\text{ not Steinberg},\\
  [\mathcal X_2^{\mathbf k}]
    +[\mathcal X_2^{\widetilde{\mathbf k}}],&
          \mathbf k\text{ Steinberg}.
 \end{cases}
\]

The special-fiber cycle for a type is the sum of these cycles
with the Jordan--H\"older multiplicities of its reduced type.
The crystalline and semistable versions use their respective
types. The additional component in the Steinberg case is part
of the theorem.

The general geometric formula of their Conjecture 8.2.2 has a
different scope. Their Theorem 8.1.4 supplies finite cycle support
for each fixed regular type. It does not assert smooth divisors
or transverse intersections. Nor does a cycle determine a
projective generator and all its morphisms.
The range \(2\le k\le p-1\) gives dual weights
\((0,k-1)\), including \((0,p-2)\) at its endpoint.
It does not reduce the full supersingular lifting problem to
two weights.

There are also positive modularity theorems with specific global
hypotheses. For example, Taylor's Theorem A in
\emph{On icosahedral Artin representations II} applies to an
irreducible continuous
\(\rho:G_{\mathbb Q}\to GL_2(\mathbb C)\) with
\(\det\rho(c)=-1\).
If its projective image is \(A_5\), impose two further conditions.
Its projective inertia image at three has odd order.
Its projective decomposition image at five is unramified of order two.
Then \(\rho\) is modular and its Artin \(L\)-function is entire.

This theorem does not give a local compact generator for every
residual \(A_5\) block.
Potential automorphy, ordinary modularity and local-global
compatibility theorems retain their residual, regularity,
cyclotomic-irreducibility and local hypotheses.
Potential automorphy alone supplies no tower of cyclic extensions
for descent.
Cyclic descent requires the actual invariance and descent data
at each step. The finite-field \(PSL_2\) and \(PGL_2\) image cases
and overlaps of lifting loci also remain in a proposed coverage
argument.

A categorical comparison can use these results after constructing
the relevant image objects and supports. For any collection of
overlapping lifting conditions, take the localizing closure of
the actual image objects, including their extensions.
A direct-sum decomposition additionally requires orthogonality.
The generator criterion then applies on this specified category.
It does not infer coverage from a list of arithmetic names.


\section{Geometric Hecke categories and spectral actions}
\label{v4-arith:ic:geometric}

Several geometric constructions realize a full representation category
or a distinguished subcategory. Their coefficient theories and
geometric spaces are part of their statements.
The equal-characteristic affine flag, the Weil--Deligne parameter
stack and the \(p\)-adic determinant-identified stack give different
recipients for an arithmetic comparison.

\subsection{Bundles and elementary modifications}

Every finite projective \(\mathbb Z\)-module is free.
A finite projective module is a direct summand of a finite free
module. A submodule of a free module over the principal ideal
domain \(\mathbb Z\) is free.
Thus the groupoid of rank-two bundles on \(\operatorname{Spec}\mathbb Z\)
has one isomorphism class, with automorphism group \(GL_2(\mathbb Z)\).
This is a statement about bundles over that base.
It does not identify their moduli functor on every test scheme
with a constant classifying stack.

Elementary modifications still occur. The inclusion
\(p\mathbb Z\oplus\mathbb Z\subset\mathbb Z^2\) has quotient
\(\mathbb F_p\) and becomes an isomorphism away from \(p\).
The existence of these lattice correspondences is compatible with
freeness of both endpoint bundles.
Their convolution, its sheaf realization and its scalar function
are separate constructions.

\subsection{The affine flag and Springer correspondences}

Let \(G\) be a split reductive group over an algebraically closed
field \(k\), and fix \(\ell\ne\operatorname{char}k\).
Use constructible \(\overline{\mathbb Q}_\ell\)-complexes on the
equal-characteristic affine flag spaces of \(G(k((t)))\).
Let \(I\) be an Iwahori subgroup, \(I^0\) its pro-unipotent radical,
and \(T=I/I^0\).
Write \(D_{II}\) for the \(I\)-equivariant derived category on
\(G(k((t)))/I\).
Write \(D_{I^0I}\) for the \(I^0\)-equivariant category on this
same space.
Finally, \(D_{I^0I^0}^{\mathrm{un}}\) is the category on
\(G(k((t)))/I^0\) with \(I^0\)-equivariance on the left and
unipotent \(T\)-monodromy on the right.

Let \(\check G\) be the Langlands dual group over
\(\overline{\mathbb Q}_\ell\), with Lie algebra \(\check{\mathfrak g}\).
Its Grothendieck and Springer resolutions are
\[
 \widetilde{\mathfrak g}
   =\{(\mathfrak b,x):x\in\mathfrak b\},
 \qquad
 \widetilde{\mathcal N}
   =\{(\mathfrak b,x):x\in\operatorname{rad}(\mathfrak b)\}.
\]
Set
\[
 \mathrm{St}=
 \widetilde{\mathfrak g}\times_{\check{\mathfrak g}}
                       \widetilde{\mathfrak g},
 \qquad
 \mathrm{St}'=
 \widetilde{\mathfrak g}\times_{\check{\mathfrak g}}
                       \widetilde{\mathcal N}.
\]
For a map to \(\check{\mathfrak g}\), the subscript
\(\mathcal N\) on coherent sheaves denotes set-theoretic
nilpotent support.

\begin{theorem}[Bezrukavnikov]\label{v4-arith:ic:affine-flag}
There are equivalences
\[
 \begin{aligned}
 D_{I^0I^0}^{\mathrm{un}}
   &\simeq D^b\mathrm{Coh}^{\check G}_{\mathcal N}(\mathrm{St}),\\
 D_{I^0I}
   &\simeq D^b\mathrm{Coh}^{\check G}(\mathrm{St}'),\\
 D_{II}
   &\simeq
 \mathrm{DGCoh}^{\check G}
 \bigl(\widetilde{\mathcal N}
       \times_{\check{\mathfrak g}}^{\mathbb L}
       \widetilde{\mathcal N}\bigr).
 \end{aligned}
\]
The first and third preserve convolution. The middle equivalence
preserves both convolution actions.
With the corresponding Iwahori--Whittaker conditions, one also has
the module-category equivalences
\[
 D^I_{\mathrm{IW}}\simeq
 D^b\mathrm{Coh}^{\check G}(\widetilde{\mathcal N}),
 \qquad
 D^{I^0,\mathrm{un}}_{\mathrm{IW}}\simeq
 D^b\mathrm{Coh}^{\check G}_{\widetilde{\mathcal N}}
                  (\widetilde{\mathfrak g}).
\]
They preserve the respective actions above.
\end{theorem}

These are Theorems 1--2 of
\emph{On two geometric realizations of an affine Hecke algebra}.
The Whittaker character is nondegenerate on the finite simple-root
quotients and trivial on the affine-root quotient, with the left
and right conventions of that paper.
The monodromic convolution category in the first line need not
contain a unit. The Iwahori category in the third line is unital.
The derived fiber product in that line retains its higher Tor.
The two other displayed Steinberg products have their ordinary
scheme descriptions under the stated resolution hypotheses.

In particular, the Springer-resolution category is a module category,
and the Iwahori category has a Steinberg carrier.
The Iwahori derived category is not the bounded derived category of
its perverse heart in general.
Passing to functions over a finite field also requires Frobenius
structures and the sheaf--function map.
These are substantive parts of any comparison with an arithmetic
Hecke category.

\subsection{Coherent Springer objects}

Let \(H\) be reductive over \(L=\overline{\mathbb Q}_\ell\) or
\(\mathbb C\).
Let \(\widehat{\mathcal N}\) be the formal completion of its Lie
algebra along the nilpotent cone, and let
\(\mu:\widetilde{\mathcal N}\to\widehat{\mathcal N}\) be the
Springer map followed by completion.
Use \(H_{\mathrm{gr}}\) for \(H\) together with the dilation action,
and let \(\mathcal L\) denote derived loops.
The coherent Springer object is
\[
 \mathcal S_H=
 (\mathcal L\mu)_*
 \mathcal O_{\mathcal L(\widetilde{\mathcal N}/H_{\mathrm{gr}})}
 \quad\text{on }\mathcal L(\widehat{\mathcal N}/H_{\mathrm{gr}}).
\]

\begin{theorem}[Ben-Zvi--Chen--Helm--Nadler]
\label{v4-arith:ic:coherent-springer}
The full derived endomorphism algebra of \(\mathcal S_H\)
is the affine Hecke algebra \(\mathcal H_H\), concentrated
in degree zero.
Tensoring with \(\mathcal S_H\) defines a fully faithful functor
\[
 D(\mathcal H_H)\longrightarrow
 \mathrm{IndCoh}\bigl(\mathcal L(\widehat{\mathcal N}/H_{\mathrm{gr}})\bigr)
\]
whose image is generated by \(\mathcal S_H\).
It is compatible with parabolic induction.
The anti-spherical module maps to the projection of the dualizing
object onto this Springer subcategory.

For \(q\in L^\times\) not a root of unity, let
\[
 \mathcal L^u_{q,H}
   =[\{(g,N):N\in\mathcal N,\ gNg^{-1}=qN\}/H]
\]
be the unipotent parameter stack and let \(\mathcal S_{q,H}\)
be the corresponding coherent specialization.
There is a fully faithful functor
\[
 \mathrm{Perf}(\mathcal H_{q,H})\longrightarrow
 \mathrm{Coh}(\mathcal L^u_{q,H})
\]
defined by \(\mathcal S_{q,H}\).
It takes the anti-spherical module to
\(\mathcal O_{\mathcal L^u_{q,H}}\) and preserves the parabolic
induction correspondences.
\end{theorem}

The assertions are Theorems 1.7 and 1.9 of
\emph{Coherent Springer theory and categorical Deligne--Langlands}.
At a non-root of unity the affine Hecke algebra has finite
cohomological dimension, so its perfect and bounded coherent
derived categories agree in this statement.
The formal derived loop carrier before specialization and the
classical reduced \(q\)-parameter stack are different objects.
In particular, the unspecialized anti-spherical image is the stated
projection of a dualizing object.

For \(GL_n\), the same work gives a larger positive comparison.
Let \(F\) be a nonarchimedean local field and
\(\mathcal L_{F,M}\) the stack of continuous Weil--Deligne parameters
for a Levi subgroup \(M\) of \(GL_n\), over \(\mathbb C\).
It retains formal deformation spaces of parameters.
Theorem 1.12 gives full embeddings
\[
 D(\mathrm{Rep}^{\mathrm{sm}}_{\mathbb C}(M(F)))
       \longrightarrow \mathrm{IndCoh}(\mathcal L_{F,M}).
\]

They send irreducible cuspidal representations to the parameter
skyscraper sheaves prescribed by local Langlands, and commute
with parabolic induction through its pull--push correspondence.
These properties characterize the supplied embeddings.
No equivalence with their whole ind-coherent targets is asserted.

\subsection{The spectral action on the Fargues--Fontaine curve}

Let \(F\) have residue cardinality \(q\), let \(\ell\) differ
from its residue characteristic, and let
\(L\) be a field over \(\mathbb Q_\ell(\sqrt q)\).
For a reductive group \(G/F\), write \(\widehat G\) for its dual.
The parameter stack
\([Z^1(W_F,\widehat G)_L/\widehat G]\)
uses continuous Weil cocycles and the specified Weil action on
\(\widehat G\).
It is not generally quasi-compact.

Fargues--Scholze, Corollary X.1.3, constructs an action
\[
 \mathrm{Perf}([Z^1(W_F,\widehat G)_L/\widehat G])
 \curvearrowright D_{\mathrm{lis}}(\mathrm{Bun}_G,L)^\omega.
\]
Here \(\mathrm{Bun}_G\) is the bundle stack on the
Fargues--Fontaine curve, and \(\omega\) denotes compact objects
of the indicated category of lisse sheaves.
The action is compactly supported: on each object it factors
through parameters trivial on some sufficiently small open
subgroup of wild inertia.
Its restrictions along the representation functors recover
the Hecke actions, functorially in every finite set of legs.

For integral coefficients, their Theorem X.0.1 also imposes
\(\ell\nmid|\pi_1(\widehat G)_{\mathrm{tor}}|\).
The rational action is a theorem. The equivalence formulated
in their Conjecture X.1.4 is a further assertion about a
Whittaker-generated functor and its right adjoint.
Neither statement changes \(D_{\mathrm{lis}}\) into coherent
sheaves on \(\mathrm{Bun}_G\), or identifies the parameter stack
with the coefficient-prime stack \(\mathcal X_\delta\).



\subsection{Factorization and Whittaker equivalences}

The critical-level comparison has a stronger factorization form.
Use a reductive group on a smooth characteristic-zero curve.
With the twisting conventions of
\emph{Proof of the geometric Langlands conjecture II},
Theorem 6.1.4 gives
\[
 \mathrm{KL}(G)_{\mathrm{crit}}
   \simeq
 \mathrm{IndCoh}^{*}
       (\mathrm{Op}_{\check G}^{\mathrm{mon\text{-}free}})
\]
as factorization categories.
Here \(\mathrm{KL}\) is the spherical critical Kac--Moody category.

A monodromy-free oper consists of a local system extending to the
disk, together with its oper structure on the punctured disk.
It is not an unrestricted monodromic local system.
Corollary 6.1.9 proves that this functor is strictly unital.
The regular-oper structure object maps to the critical vacuum.

At nondegenerate levels, Campbell--Dhillon--Raskin construct
the spherical and affine-flag DG equivalences
\[
 \mathrm{Whit}^{\mathrm{sph}}_\kappa
 \simeq\widehat{\check{\mathfrak g}}_{\check\kappa}
                 \text{-}\mathrm{mod}^{L^+\check G},
 \qquad
 \mathrm{Whit}^{\mathrm{aff}}_\kappa
 \simeq\widehat{\check{\mathfrak g}}_{\check\kappa}
                 \text{-}\mathrm{mod}^{\check I}.
\]
The left categories are twisted Whittaker D-modules on the
affine Grassmannian and flag variety. The right categories are
Harish--Chandra DG module categories for the displayed affine
Lie algebra and subgroup.
The dual level uses the critical shifts of their \S2.7.

Their good-level hypothesis has a concrete form.
On each simple factor an irrational level is good.
For a rational level
\(\kappa^\circ=(-h^\vee+a/b)\kappa_{\mathrm{bas}}\),
with \((a,b)=1\), the denominator \(b\) must be coprime to
\[
 \begin{array}{c|ccccccccc}
 \text{type}&A_n&B_n&C_n&D_n&E_6&E_7&E_8&F_4&G_2\\ \hline
 n(\mathfrak g)&1&2&2&2&6&6&30&6&6 .
 \end{array}
\]
This is Proposition 3.4.5.
These DG equivalences in that paper set aside the factorization
enhancement. They are not a theorem about arbitrary quantum
parameters on an arithmetic base.

Gaitsgory--Lysenko's small local equivalence has a third precise
recipient. Choose their geometric metaplectic datum on a smooth curve.
Let \(q\) be its restricted Weyl-invariant root-of-unity quadratic
form, and \(H\) the associated dual datum.
Fix the required trivialization at the marked point.

Retain the nondegeneracy
and sheaf-theory hypotheses of that datum. Nondegeneracy means
that \(q(\alpha)\) has order different from one for every
coroot \(\alpha\) of \(G\), as in Definition 2.3.4.
Their graded Hecke category is formed using the metaplectic
Satake action of \(\mathrm{Rep}(H)\) and restriction to its
Cartan \(T_H\).
Theorem 19.2.5 identifies it with the renormalized category of
factorization modules:
\[
 \mathrm{Hecke}^{\bullet}(\mathrm{Whit}_{q,x}(G))
   \simeq
 \bigl(\Omega_q^{\mathrm{small}}
           \text{-}\mathrm{FactMod}\bigr)^{\mathrm{ren}}.
\]

The renormalization retains the specified Artinian structure and
compact objects. Replacing it by the unrenormalized category
changes the statement.
The allowed sheaf theories include the stated D-module,
constructible and \(\ell\)-adic settings.
The factorization-category strengthening is distinguished in
their \S0.4.5 from the categorical theorem they formulate.

These are geometric positive comparisons with specified units,
levels and enhancements. A value \(p^{-s}\) in a Hecke spectrum
does not identify their quantum parameter with an arithmetic one.
An arithmetic use requires the coefficient and operation-preserving
map, including its collision and unit data.


\section{Normalized intertwining for unramified principal series}

An Eisenstein comparison must distinguish the normalization of an
intertwining operator from its passage to a category of representations.
For unramified principal series over a nonarchimedean field, the operator
can be computed on every locally constant vector.
The computation also determines its poles and its degenerations.

Let $F$ be a finite extension of $\mathbb Q_p$.
Write $\mathcal O$ for its ring of integers, $\varpi$ for a uniformizer,
and $q=|\mathcal O/\varpi\mathcal O|$.
Normalize $|\varpi|=q^{-1}$ and the additive Haar measure by
$\operatorname{vol}(\mathcal O)=1$.
Fix an additive character $\psi:F\to\mathbb C^\times$ that is trivial on
$\mathcal O$ and nontrivial on $\varpi^{-1}\mathcal O$.

The operator below does not depend on $\psi$.
This choice fixes the additive convention for subsequent local-factor comparisons.
All representations and function spaces have coefficients in $\mathbb C$.

Put $G=\operatorname{GL}_2(F)$, $K=\operatorname{GL}_2(\mathcal O)$,
and let $B$ be the upper triangular subgroup.
For unramified smooth characters $\chi_1,\chi_2:F^\times\to\mathbb C^\times$,
let $I(\chi_1,\chi_2)$ consist of locally constant functions $f:G\to\mathbb C$
satisfying
\[
 f\left(\begin{pmatrix}a&b\\0&d\end{pmatrix}g\right)
 =|a/d|^{1/2}\chi_1(a)\chi_2(d)f(g).
\]
The action is right translation, $(\pi(g)f)(h)=f(hg)$.
Thus $I(\chi_1,\chi_2)$ is normalized smooth induction.

Set
\[
 z=\frac{\chi_1(\varpi)}{\chi_2(\varpi)},\qquad
 w=\begin{pmatrix}0&1\\1&0\end{pmatrix},\qquad
 n(x)=\begin{pmatrix}1&x\\0&1\end{pmatrix}.
\]
For $|z|<1$, define
\begin{equation}
\label{v4-arith:ic:eis:raw-intertwiner}
 (M_zf)(g)=\int_F f(wn(x)g)\,dx.
\end{equation}
The target is $I(\chi_2,\chi_1)$.
The subscript $z$ specifies the ratio parameter.
The source and target characters remain part of the operator's data.

\subsection{The compact picture and ball averages}

Restriction to $K$ identifies every $I(\chi_1,\chi_2)$ with
\[
 V=C^\infty(X;\mathbb C),\qquad X=(B\cap K)\backslash K=\mathbb P^1(F).
\]
Here $C^\infty$ means locally constant.
The identification of $X$ sends a matrix to its bottom row, up to scalar.
It follows from $G=BK$ and the triviality of unramified characters on
$\mathcal O^\times$.

A primitive row has at least one unit coordinate.
For primitive representatives $u,v\in\mathcal O^2$, define
\[
 \rho([u],[v])=|\det(u,v)|.
\]
This function is independent of the primitive representatives and is
invariant under the right action of $K$.
Let $\nu$ be the $K$-invariant probability measure on $X$.

For $n\ge1$, reduction modulo $\varpi^n$ gives a finite set
\[
 X_n=\mathbb P^1(\mathcal O/\varpi^n\mathcal O),\qquad
 |X_n|=(q+1)q^{n-1}.
\]
Let $V_n\subset V$ be the functions that factor through $X_n$.
Put $V_0=\mathbb C\mathbf1$ and
$H=L^2(X,\nu)$, with inner product linear in the first variable.
The orthogonal projection $E_n:H\to V_n$ averages over each reduction fiber.

The projection $E_0$ is global averaging.
The fibers are the balls $\rho(x,y)\le q^{-n}$, each of measure
\[
 b_n=\frac{1}{(q+1)q^{n-1}}.
\]
Every locally constant function factors through some $X_n$.
Consequently $V=\bigcup_{n\ge0}V_n$ is dense in $H$.

Define mutually orthogonal finite-dimensional spaces
\[
 H_0=V_0,\qquad H_n=V_n\cap V_{n-1}^{\perp}\quad(n\ge1).
\]
Then
\[
 V=\bigoplus_{n\ge0}^{\mathrm{alg}}H_n,
 \qquad H=\widehat{\bigoplus}_{n\ge0}H_n.
\]
In particular, $\dim H_1=q$ and
$\dim H_n=(q+1)(q-1)q^{n-2}$ for $n\ge2$.

\begin{lemma}[Integral kernel and meromorphic continuation]
\label{v4-arith:ic:eis:kernel}
For $|z|<1$, the integral in
\eqref{v4-arith:ic:eis:raw-intertwiner} converges absolutely.
It defines a smooth $G$-intertwiner.
In the compact picture it is
\begin{equation}
\label{v4-arith:ic:eis:projection-expansion}
 M_z=(1+q^{-1})E_0+(z-q^{-1})\sum_{n\ge1}z^{n-1}E_n.
\end{equation}
The series converges in operator norm on $H$.

On the orthogonal summands it acts by
\[
 M_z|_{H_0}=c(z),\qquad
 c(z)=\frac{1-q^{-1}z}{1-z},\qquad
 M_z|_{H_n}=\frac{(z-q^{-1})z^{n-1}}{1-z}\quad(n\ge1).
\]
These formulas give a meromorphic continuation on every $V_N$.
On every nonzero $V_N$, its only pole in $z\in\mathbb C^\times$
is the simple pole at $z=1$.
Here meromorphic continuation means continuation on each finite level,
not a family of bounded operators on $H$ for all parameters.
\end{lemma}

\begin{proof}
Let $x_0=[(0,1)]$.
For $|x|\le1$, the matrix $wn(x)$ belongs to $K$.
For $|x|>1$, the factorization
\[
 wn(x)=
 \begin{pmatrix}-x^{-1}&1\\0&x\end{pmatrix}
 \begin{pmatrix}1&0\\x^{-1}&1\end{pmatrix}
\]
shows that the inducing factor on $|x|=q^n$ is $q^{-n}z^n$.
The shell has additive measure $(1-q^{-1})q^n$.
Since compact-picture functions are bounded, the tail is dominated by
$a\sum_{n\ge1}|z|^n$ for a finite constant $a$.
The same argument applies to each right translate of $f$.

The integral commutes with right translations.
Changing $x$ to $x+y$ proves invariance under left multiplication by $n(y)$.
For $a,d\in F^\times$, use
\[
 wn(x)\begin{pmatrix}a&0\\0&d\end{pmatrix}
 =\begin{pmatrix}d&0\\0&a\end{pmatrix}wn(xd/a).
\]
The substitution $u=xd/a$ contributes $|a/d|$.

The resulting factor is $|a/d|^{1/2}\chi_2(a)\chi_1(d)$,
which is the target induction law.
Right local constancy is preserved by the integral.
Thus $M_z$ is the stated smooth intertwiner.

The chart $[(1,x)]$, $x\in\mathcal O$, has measure
$q(q+1)^{-1}dx$.
The chart $[(y,1)]$, $y\in\varpi\mathcal O$, has the same density.
For example, reduction modulo $\varpi^N$ divides either chart into
sets of additive measure $q^{-N}$ and $\nu$-measure
$((q+1)q^{N-1})^{-1}$.
This proves the density assertion on locally constant functions.

In the second chart, the substitution $y=x^{-1}$ has Jacobian
$dx=|y|^{-2}dy$.
On $|y|=q^{-n}$, the resulting density in the raw integral is
$q^nz^n\,dy$.
Thus its kernel, relative to $\nu$, has value $1+q^{-1}$ when
$\rho(x_0,y)=1$, and value $(1+q^{-1})q^nz^n$ when
$\rho(x_0,y)=q^{-n}$.
The diagonal has measure zero.
Right $K$-equivariance gives the same radial kernel about every point.

For $\rho(x,y)=q^{-m}$, telescope the kernel as
\[
 (1+q^{-1})\left[
 1+\sum_{n=1}^{m}\bigl(q^nz^n-q^{n-1}z^{n-1}\bigr)\right].
\]
Integration of the $n$th ball indicator is $b_nE_n$.
Since $(1+q^{-1})b_n=q^{-n}$, its coefficient is
$(z-q^{-1})z^{n-1}$.
Absolute convergence justifies the sum and proves
\eqref{v4-arith:ic:eis:projection-expansion}.
The projection norms are one, so the series also converges in operator norm.

On $H_0$, every $E_n$ acts as the identity.
On $H_m$, $m\ge1$, the projection $E_n$ is zero for $n<m$ and
is the identity for $n\ge m$.
Summing the geometric series gives the displayed eigenvalues.
Each $V_N$ is a finite sum of these spaces.

The rational formulas therefore give its meromorphic continuation.
At $z=1$, every numerator has nonzero value $1-q^{-1}$.
The pole is simple on every nonzero summand.
\end{proof}

\subsection{Normalization, composition, and domains}

The spherical vector is $\mathbf1\in H_0$.
The scalar in the preceding lemma equals the unramified local ratio
\[
 c(z)=\frac{L(0,\chi_1\chi_2^{-1})}
 {L(1,\chi_1\chi_2^{-1})},\qquad
 L(u,\chi_1\chi_2^{-1})=(1-zq^{-u})^{-1}.
\]
Normalizing the spherical vector therefore requires division by $c(z)$.
Define the meromorphic family
\begin{equation}
\label{v4-arith:ic:eis:normalized-intertwiner}
 R_z=c(z)^{-1}M_z.
\end{equation}

\begin{theorem}[The normalized unramified intertwiner]
\label{v4-arith:ic:eis:normalized-theorem}
\label{v4-arith:ic:w6-d:prop:local-FE}
The compact-picture eigenvalues of $R_z$ are
\begin{equation}
\label{v4-arith:ic:eis:normalized-eigenvalues}
 r_0(z)=1,\qquad
 r_n(z)=z^{n-1}\frac{z-q^{-1}}{1-q^{-1}z}\quad(n\ge1).
\end{equation}
The operator is holomorphic on each $V_N$ away from $z=q$.
It has a simple pole on every $H_n$ with $n\ge1$ at $z=q$.
The value at $z=1$ is $\operatorname{id}_V$.

The value at $z=q^{-1}$ is $E_0$.
For $z\notin\{q,q^{-1}\}$, the smooth maps satisfy
\[
 R_{z^{-1}}R_z=\operatorname{id}_V,
 \qquad R_zR_{z^{-1}}=\operatorname{id}_V.
\]
Here the second operator uses the swapped inducing characters.

If $\chi_1$ and $\chi_2$ are unitary, $R_z$ extends uniquely to a
unitary intertwiner between their induced Hilbert representations on $H$.
Its Hilbert domain is all of $H$.
In particular, the equal-character point $z=1$ is regular.
\end{theorem}

\begin{proof}
Division of the preceding eigenvalues by $c(z)$ gives
\eqref{v4-arith:ic:eis:normalized-eigenvalues}.
This calculation cancels the apparent singularity at $z=1$.
At $z=q$, the denominator has a simple zero and the numerator is nonzero
on every $H_n$, $n\ge1$.
Substitution gives the stated values at $1$ and $q^{-1}$.

For $n\ge1$, direct cancellation gives
\[
 r_n(z)r_n(z^{-1})
 =\frac{(z-q^{-1})(1-q^{-1}z)}
 {(1-q^{-1}z)(z-q^{-1})}=1.
\]
The identity also holds on $H_0$.
Every smooth vector has finitely many components, so both compositions
are defined on all of $V$ at the stated parameters.
Intertwining holds first on the convergence domain of the raw integral.

On each smooth vector and each fixed group element, the compact-picture
formulas are rational in the two nonzero unramified character values.
For a fixed $g\in G$, right translation sends $V_N$ into some $V_M$
independently of the unramified character values.
Indeed, an open subgroup of $K\cap gK_Ng^{-1}$ fixes the translated
vectors, where $K_N=1+\varpi^NM_2(\mathcal O)$ for $N\ge1$.
On the resulting finite partition of $K$, the Iwasawa factors have
fixed diagonal valuations.

Their character values are Laurent monomials.
The integral intertwining identity therefore continues meromorphically.
This also proves intertwining at every regular specialization.

If the inducing characters are unitary, $|z|=1$.
Then $|z-q^{-1}|=|1-q^{-1}z|$, so $|r_n(z)|=1$ for every $n$.
The orthogonal decomposition of $H$ gives a unique unitary extension.
Its inverse has parameter $z^{-1}$.

For completeness, the induced action preserves the stated Hilbert norm.
In the open Bruhat chart put $F_f(x)=f(wn(x))$.
For unitary inducing characters, the chart measures above give
\[
 \|f|_K\|_H^2=\frac{q}{q+1}\int_F|F_f(x)|^2\,dx.
\]
Right translation by $n(y)$ gives $F_f(x+y)$.

Right translation by $\operatorname{diag}(a,d)$ gives
\[
 |d/a|^{1/2}\chi_1(d)\chi_2(a)F_f(xd/a).
\]
Right translation by $w$, away from $x=0$, gives
\[
 |x|^{-1}\chi_1(-x^{-1})\chi_2(x)F_f(x^{-1}).
\]
Translations, dilations, and inversion preserve the displayed norm
with these factors and their respective Jacobians.
These elements generate $G$.

Each smooth vector has an open stabilizer. Density and unitarity
therefore imply strong continuity of the extended action.
The smooth equivariance of $R_z$ extends by density to these
induced unitary representations.
\end{proof}

For nonunitary parameters, the same compact-picture formulas also specify
an operator on a fixed Hilbert space.
This assertion does not equip the induced representation with an invariant
positive inner product.

\begin{proposition}[Closed compact-picture realization]
\label{v4-arith:ic:eis:maximal-domain}
For $z\ne q$, write $h=\sum_{n\ge0}h_n$ with $h_n\in H_n$.
The closure of $R_z|_V$ on $H$ is the normal operator with domain
\[
 \operatorname{Dom}(\overline R_z)=
 \left\{h\in H:\sum_{n\ge0}|r_n(z)|^2\|h_n\|^2<\infty\right\},
 \qquad \overline R_zh=\sum_{n\ge0}r_n(z)h_n.
\]
The space $V$ is a core.
This operator is bounded exactly when $|z|\le1$.
The composition assertions of the preceding theorem hold on $V$.
On the unit circle they hold on all of $H$.
\end{proposition}

\begin{proof}
The orthogonal sum of scalar operators with the displayed maximal domain
is closed and normal.
Its adjoint has coefficients $\overline{r_n(z)}$ and the same domain.
Truncating the expansion after $n=N$ converges in its graph norm.
Thus this maximal operator is the closure of its restriction to $V$.

Boundedness is equivalent to $\sup_n|r_n(z)|<\infty$.
Formula~\eqref{v4-arith:ic:eis:normalized-eigenvalues} proves boundedness for
$|z|\le1$.
For $|z|>1$ and $z\ne q$, the nonzero factor
$(z-q^{-1})/(1-q^{-1}z)$ is fixed and $|z|^{n-1}$ is unbounded.
\end{proof}

\subsection{Reducibility and the unitary axis}

The poles of a normalization and the reducibility of the induced
representation are different loci.
They can both be determined in the present unramified family.

\begin{proposition}[Reducibility locus]
\label{v4-arith:ic:eis:reducibility}
The smooth representation $I(\chi_1,\chi_2)$ is reducible exactly when
\[
 \chi_1\chi_2^{-1}=|\cdot|\quad\text{or}\quad
 \chi_1\chi_2^{-1}=|\cdot|^{-1}.
\]
Equivalently, $z\in\{q^{-1},q\}$.
For $z=q^{-1}$, the regular intertwiner $R_z$ has image $H_0$
in $I(\chi_2,\chi_1)$ and kernel in $I(\chi_1,\chi_2)$ equal to $\bigoplus_{n\ge1}^{\mathrm{alg}}H_n$.
For $z=q$, the swapped regular intertwiner $R_{z^{-1}}$ has image
$H_0$ in $I(\chi_1,\chi_2)$.
The stated kernel at $z=q^{-1}$ and image at $z=q$ are nonzero
proper invariant subspaces of $I(\chi_1,\chi_2)$.
\end{proposition}

\begin{proof}
The kernel and image assertions follow from
\eqref{v4-arith:ic:eis:normalized-eigenvalues} and the regular intertwining
identity.
They prove reducibility at the two displayed ratios.
Away from these ratios, apply Bernstein--Zelevinsky's irreducibility
criterion to the two characters of $\operatorname{GL}_1(F)$.
This is \emph{Induced representations of reductive $p$-adic groups I},
Theorem~4.2, p.~454.
The induction conventions are in \S1.8(a), p.~444,
and \S\S2.2--2.3, p.~446.

For the upper triangular parabolic, their modulus on the unipotent
radical is $|a/d|$, so their square-root factor is the one used here.
Every character of $\operatorname{GL}_1(F)$ is irreducible and cuspidal.
The criterion excludes exactly the relations
$\chi_j=|\cdot|\chi_i$ for distinct $i,j$.
Thus it proves irreducibility on the stated complement.
Only this direction of their criterion is used.
\end{proof}

For the symmetric family
$\chi_1=|\cdot|^{s/2}$, $\chi_2=|\cdot|^{-s/2}$, one has $z=q^{-s}$.
The raw integral converges for $\Re(s)>0$.
Its poles occur at
\[
 s\in\frac{2\pi i}{\log q}\mathbb Z.
\]
On every smooth finite level, the normalized operator is regular on
$\Re(s)\ge0$.

Its poles occur at
\[
 s\in-1+\frac{2\pi i}{\log q}\mathbb Z,
\]
and reducibility occurs at
\[
 s\in\left(1+\frac{2\pi i}{\log q}\mathbb Z\right)
 \cup\left(-1+\frac{2\pi i}{\log q}\mathbb Z\right).
\]
Neither reducibility locus meets the unitary axis $\Re(s)=0$.
These statements concern local unramified induction.
Poles of a global Eisenstein series require a separate global construction.

\begin{example}[The first nonconstant vectors]
Let $q=2$ and $z=-1$.
The three reduction fibers of $X\to\mathbb P^1(\mathbb F_2)$ have equal measure.
On $V_1$, the raw integral and normalized operator are
\[
 M_{-1}=\frac14
 \begin{pmatrix}-1&2&2\\2&-1&2\\2&2&-1\end{pmatrix},
 \qquad
 R_{-1}=\frac13
 \begin{pmatrix}-1&2&2\\2&-1&2\\2&2&-1\end{pmatrix}.
\]
The spherical scalar is $c(-1)=3/4$.

The normalized matrix fixes constants and negates their orthogonal complement.
Its square is the identity.
At $z=1/2$, the normalized matrix is the rank-one averaging projection.
At $z=1$, its continued value is the identity matrix.
\end{example}

\subsection{The automorphic comparison}

A categorical Eisenstein comparison requires categories and functors in
addition to these operators.
One must specify the coefficient field, the local parameter spaces,
and the global automorphic quotient with its central character and level.
The global category must carry its Hecke actions and its Eisenstein subcategory.
A comparison functor must identify the relevant derived mapping objects
and determine its essential image.
The shared global induction parameter must also be imposed on local families.

The construction above supplies the full smooth intertwiner and its Hilbert
extension only for unramified nonarchimedean induction.
Ramified characters and the archimedean place require their own operators,
normalizations, and analytic domains.
A global restricted product requires compatible distinguished vectors
and an explicit global family.
The categorical comparison further requires functors implementing these
maps and their Weyl coherences.
No equivalence of automorphic categories follows from the compact-picture
Hilbert isomorphism alone.




\subsection{The place-indexed operator notation}
\begin{definition}[the raw place-indexed operator]
\label{v4-arith:ic:w6-d:def:M-unnormalised}
For a finite local field $F_v$ and unramified characters
$\chi=(\chi_1,\chi_2)$, use normalized smooth induction
\[
 I_v(s,\chi)=\operatorname{Ind}_{B(F_v)}^{GL_2(F_v)}
       (\chi_1|\cdot|^{s/2}\otimes\chi_2|\cdot|^{-s/2}).
\]
Put
$\alpha_1=\chi_1|\cdot|^{s/2}$,
$\alpha_2=\chi_2|\cdot|^{-s/2}$ and
$z=\alpha_1(\varpi)/\alpha_2(\varpi)$.
With $\chi^w=(\chi_2,\chi_1)$, the notation is
\begin{equation}
\label{v4-arith:ic:w6-d:eq:M-unnormalised}
 M_v(s,\chi)=M_z:
 I_v(s,\chi)=I(\alpha_1,\alpha_2)
 \longrightarrow I(\alpha_2,\alpha_1)=I_v(-s,\chi^w).
\end{equation}
The integral is defined for $|z|<1$ and continued on each smooth
finite level by Lemma~\ref{v4-arith:ic:eis:kernel}.
\end{definition}
\begin{definition}[the normalized place-indexed operator]
\label{v4-arith:ic:w6-d:def:N-normalised}
On the same source and target, define
\begin{equation}
\label{v4-arith:ic:w6-d:eq:N-normalised}
 N_v(s,\chi)=R_z=c(z)^{-1}M_z,
 \qquad c(z)=\frac{1-q^{-1}z}{1-z}.
\end{equation}
Cancellation and continuation at $z=1$ are part of
Theorem~\ref{v4-arith:ic:eis:normalized-theorem}. The expression is not
a pointwise product of undefined values there. Its poles, domains
and unitary specializations are those proved in that theorem.
\end{definition}

\section{Algebraic parameters for unramified induction}

The normalized unramified intertwiner varies with the inducing characters.
A category of coefficient complexes records this variation and its Weyl symmetry.
Its relation to smooth induction is a functor defined by tensoring
with universal smooth induction.
The resulting category concerns local algebraic parameters.
An automorphic comparison requires further constructions.

Retain the nonarchimedean field $F$, residue cardinality $q$, group
$G=\operatorname{GL}_2(F)$, and compact subgroup $K$ of the unramified calculation.
In particular, $|\varpi|=q^{-1}$ and $\operatorname{vol}(\mathcal O)=1$.
Let $a,b$ denote the values at $\varpi$ of the two unramified characters.
Their ratio is $z=a/b$.
To make both normalized intertwining directions regular and invertible, put
\[
 A=\mathbb C[a^{\pm1},b^{\pm1},(a-qb)^{-1},(b-qa)^{-1}],
 \qquad T=\operatorname{Spec}A.
\]

Thus $T$ omits precisely the two divisors $z=q,q^{-1}$.
The Weyl group $W=\{1,w\}$ acts through the ring involution
$\sigma(a)=b$, $\sigma(b)=a$.
The involution exchanges the two inverted divisors.
At a complex point $(a_0,b_0)$, define
$\chi_{a_0}(x)=a_0^{v_F(x)}$, where $v_F(\varpi)=1$.
The corresponding representation is the normalized induction
$I(\chi_{a_0},\chi_{b_0})$.

The ratio alone determines the normalized operator in the compact picture.
It does not determine the representation.
Replacing $(a_0,b_0)$ by $(ta_0,tb_0)$ twists the representation by
$\chi_t\circ\det$.
For example, the central element $\varpi\operatorname{id}$ acts by
$a_0b_0$ before the twist and by $t^2a_0b_0$ afterwards.
Both characters must therefore remain in the parameter space.

\subsection{Coefficient complexes and their unit}

For a ring $D$, write $\mathcal D(D)$ for the stable $\infty$-category
of unbounded cochain complexes of left $D$-modules, with quasi-isomorphisms inverted.
The differential has degree $1$.
We use $P[1]^n=P^{n+1}$ with differential $-d_P$.
For commutative $D$, tensor products are derived over $D$.
Their symmetry has sign $(-1)^{ij}$ on degrees $i,j$.
Let $\operatorname{Perf}(D)$ be a small skeleton of the full subcategory
generated by $D[0]$ under finite sums, shifts, cones, and retracts.

\begin{proposition}[The coefficient category]
\label{v4-arith:ic:eis:parameter-perf}
The category $\operatorname{Perf}(A)$ is stable, idempotent complete,
and symmetric monoidal under $\otimes_A^{\mathbf L}$.
There is a canonical equivalence
\[
 \operatorname{Ind}(\operatorname{Perf}(A))\simeq\mathcal D(A).
\]
The compact objects of $\mathcal D(A)$ are exactly $\operatorname{Perf}(A)$.
The monoidal unit is $A[0]$.
The functor
\[
 e:\mathcal D(\mathbb C)\longrightarrow\mathcal D(A),
 \qquad U\longmapsto A\otimes_{\mathbb C}U
\]
is a colimit-preserving symmetric monoidal pointing and preserves compact objects.
\end{proposition}

\begin{proof}
The category of derived modules is stable and presentable.
For a commutative ring it identifies symmetric monoidally with
modules over the associated discrete ring spectrum.
This is Lurie, \emph{Higher Algebra} (2017),
Theorem~7.1.2.13, pp.~1212--1213.
Its shifts of $A$ detect zero objects, since
$\operatorname{Hom}_{\mathcal D(A)}(A[-n],P)=H^n(P)$.

These shifts are compact because cohomology commutes with filtered colimits of modules.
The module compact-generation theorem identifies their finite
stable retract closure with all compact objects and gives the Ind equivalence.
See Lurie, \emph{Higher Algebra} (2017), Definition~7.2.4.1
and Proposition~7.2.4.2, pp.~1260--1261.
For ordinary rings these are bounded complexes of finitely generated projective modules
(The Stacks Project, Tag 07LQ, Lemma~15.80.1 and Proposition~15.80.3).

Tensoring two such complexes gives another such complex.
The associativity, unit, and symmetry are those of derived tensor product.
Extension of scalars preserves colimits and takes a finite complex of
finite-dimensional vector spaces to a finite free complex.
Its tensor and unit identities prove the assertion about $e$.
\end{proof}

A point of $T$ gives the specialization functor
$\mathbb C\otimes_{A,(a_0,b_0)}^{\mathbf L}(-)$.
The point module $A/(a-a_0,b-b_0)$ is also perfect.
Its length-two Koszul resolution follows because $a-a_0,b-b_0$ form a regular sequence.
Indeed, the quotient by the first element is a localization of
$\mathbb C[b^{\pm1}]$, and the second element is nonzero there.

The inverted elements remain nonzero at the chosen point.
This point module differs from the monoidal unit $A$.
For two distinct points, its tensor product with the other point module is zero,
as follows by localizing at their disjoint supports.
The unit cannot have this property.

\subsection{The universal smooth representation}

Let $X=(B\cap K)\backslash K$ and $V=C^\infty(X;\mathbb C)$,
where $B$ is the upper triangular subgroup and functions are locally constant.
The compact-picture calculation gives
$V=\bigoplus_{n\ge0}H_n$ and $V_N=\bigoplus_{n=0}^NH_n$,
with $H_0=\mathbb C\mathbf1$ and finite-dimensional $H_n$.
Put $E=A\otimes_{\mathbb C}V$ and $E_N=A\otimes V_N$.
All these modules are in cochain degree zero.

\begin{proposition}[Induction over the parameter ring]
\label{v4-arith:ic:eis:parameter-induction}
There is an $A$-linear smooth action $\pi_A$ of $G$ on $E$.
Its specialization at $(a_0,b_0)$ is $I(\chi_{a_0},\chi_{b_0})$.
The normalized intertwiner defines an $A$-linear isomorphism
\[
 R:E\longrightarrow E^{\mathrm{sw}},\qquad
 R|_{A\otimes H_0}=1,\qquad
 R|_{A\otimes H_n}=z^{n-1}\frac{z-q^{-1}}{1-q^{-1}z}\quad(n\ge1),
\]
where $E^{\mathrm{sw}}$ has the inducing parameters $b,a$.
Its inverse has the displayed coefficients with $z$ replaced by $z^{-1}$.
The coefficient module $E$ is free and is not compact in $\mathcal D(A)$.
Every $E_N$ is compact and $K$-stable.
\end{proposition}

\begin{proof}
Use the universal unramified characters
$\chi_a(x)=a^{v_F(x)}$ and $\chi_b(x)=b^{v_F(x)}$ in the induction law.
The factor $|x/y|^{1/2}$ is a fixed positive real scalar.
For $g\in G$, extend a function from $K$ to $BK=G$ by this law,
then translate it on the right by $g$.
Independence of the Iwasawa factorization follows because the
characters and modulus are trivial on $B\cap K$.
For each fixed $g,N$, compactness and local constancy give a finite partition
on which all diagonal valuations are constant.

The resulting coefficients are Laurent monomials in $a,b$.
Thus $\pi_A(g)E_N\subset E_M$ for some finite $M$.
Right translation satisfies the group law before restriction to $K$,
and consequently defines an action on $E$.
Each vector has an open stabilizer in $K$.
Specializing the induction law proves the fiber assertion.

The stated coefficients belong to $A$ and are units in $A$.
The normalized intertwiner calculation proves the intertwining identity
and $r_n(z^{-1})r_n(z)=1$.
These identities apply over $A$: for fixed $g,N$ they are finite matrix
identities in rational functions of $a,b$.
They hold at every regular complex specialization, hence hold in $A$.
This proves the asserted isomorphism and inverse.

Each $E_N$ is finite free.
The decomposition of $V$ makes $E$ free of infinite rank.
If $E$ were compact, its identity would factor through some term of
$E=\operatorname{colim}_NE_N$.
That would express an infinite-rank free module as a retract of a finite-rank module,
which is impossible after tensoring with the fraction field of $A$.
The compact-picture $K$-action preserves each $V_N$.
The full $G$-action need not preserve these finite levels.
\end{proof}

Let $\mathcal D_{\mathrm{sm}}(G,A)$ denote the derived $\infty$-category
of smooth $A$-linear $G$-representations.
Here smooth means that every vector has an open stabilizer, and $G$ acts trivially on $A$.
Tensoring with the flat module $E$ gives an actual functor
\[
 \mathcal F:\mathcal D(A)\longrightarrow\mathcal D_{\mathrm{sm}}(G,A),
 \qquad P\longmapsto E\otimes_A P.
\]
It is well-defined on quasi-isomorphisms since $E$ is free over $A$.

Its specialization is
$I(\chi_{a_0},\chi_{b_0})\otimes_{\mathbb C}
(\mathbb C\otimes_A^{\mathbf L}P)$.
The spherical map $A\to E$, $c\mapsto c\mathbf1$, is $K$-equivariant.
It is generally not $G$-equivariant when $A$ has trivial $G$-action.
For example, the central element $\varpi\operatorname{id}$ acts on $E$ by $ab$.
Thus neither the spherical vector nor the representation $E$
replaces the coefficient category's monoidal unit.

\subsection{Weyl descent}

The involution $\sigma$ acts on coefficient complexes by restriction of scalars.
On the underlying complex the new scalar action is $c\cdot m=\sigma(c)m$.
Applying this operation twice is the identity.
It preserves derived tensor products and the unit, with their usual canonical identifications.
This is a specified coherent action of $W$ on $\mathcal D(A)$.
Its descent category is the homotopy limit $\mathcal D(A)^{hW}$.

\begin{proposition}[An explicit descent category]
\label{v4-arith:ic:eis:parameter-descent}
Let $D=A\rtimes W=A\oplus Au$, with multiplication given by
$u^2=1$ and $uc=\sigma(c)u$ for $c\in A$.
There is an equivalence
\[
 \mathcal D(A)^{hW}\simeq\mathcal D(D).
\]
This category is stable and presentable.
An object is compact if and only if its underlying $A$-complex is perfect.
Its symmetric monoidal product is derived tensor over $A$, with diagonal Weyl action.
Its compact monoidal unit is $A$ with $u(c)=\sigma(c)$.
\end{proposition}

\begin{proof}
A strict $D$-module is an $A$-module with a semilinear involution $\theta$:
$\theta(cm)=\sigma(c)\theta(m)$ and $\theta^2=1$.
The same description applies to complexes with a chain involution.
To identify the derived category with coherent descent, consider the
evaluation functor from the homotopy limit to $\mathcal D(A)$.
It is conservative and has the free-equivariant left adjoint
$P\mapsto P\oplus\sigma^*P$ with its permutation action.

Colimits are computed on underlying complexes since $\sigma^*$ preserves them.
In particular, evaluation preserves geometric realizations of split simplicial objects.
The monadicity theorem therefore applies
(Lurie, \emph{Higher Algebra} (2017), Theorem~4.7.3.5, p.~685).
Its monad is $P\mapsto P\oplus\sigma^*P$, with multiplication from $W$.

This is also the monad of restriction from $D$-modules to $A$-modules,
whose left adjoint is $D\otimes_A^{\mathbf L}(-)$.
Restriction is conservative and preserves geometric realizations.
The two monadic descriptions prove the equivalence, including all higher coherences.

As a left or right $A$-module, $D$ is free of rank two.
Thus every perfect $D$-complex becomes perfect over $A$.
Conversely, a $D$-complex $M$ is a retract of $D\otimes_A M$.
The action map is a retraction of the $D$-linear chain map
\[
 m\longmapsto\tfrac12(1\otimes m+u\otimes\theta(m)).
\]

The relations $uc=\sigma(c)u$ and $\theta^2=1$ verify linearity directly.
If the underlying $A$-complex is perfect, extension of scalars makes
$D\otimes_A^{\mathbf L}M$ perfect over $D$.
The retract is therefore perfect.
The compact-object characterization in the coefficient proposition applies to $D$ as well.

The diagonal involution on $M\otimes_A^{\mathbf L}N$ is
$\theta_M\otimes\theta_N$.
Semilinearity makes it well-defined, and the usual tensor constraints are equivariant.
The module $A$ with involution $\sigma$ is the unit and is perfect over $A$.
This proves its compactness after descent.
\end{proof}

The involution is not free on $T$.
Every point $(a_0,a_0)$ is fixed and belongs to $T$.
The modules $A$ with involutions $\sigma$ and $-\sigma$ specialize there
to the trivial and sign representations of $W$.
These are not isomorphic: an intertwining scalar would satisfy $c=-c$.
Consequently descent retains stabilizer representations.
It cannot be replaced by the category of modules on the coarse orbit space
through ordinary pullback, which has trivial stabilizer action on every fiber.

The representation $E$ has its own compatible descent datum.
Let $S:E\to E$ apply $\sigma$ to coefficients and fix $V$.
Define $\Theta=SR$, using the common compact-picture module.
Then
\[
 \Theta(c f)=\sigma(c)\Theta(f),\qquad
 \Theta^2=1,\qquad \Theta\pi_A(g)=\pi_A(g)\Theta.
\]
Indeed, $S\pi_A(g)=\pi_A^{\mathrm{sw}}(g)S$ and
$R\pi_A(g)=\pi_A^{\mathrm{sw}}(g)R$.

Conjugating $R$ by $S$ replaces $z$ by $z^{-1}$, so its inverse is obtained.
These identities prove all three assertions.
Moreover, $\Theta(c\mathbf1)=\sigma(c)\mathbf1$.
For an equivariant coefficient complex $(P,\theta)$, the formula
$(E\otimes_A P,\Theta\otimes\theta)$ therefore defines equivariant smooth induction.
This is a construction of a functor, without an assertion of full faithfulness or essential surjectivity.

\subsection{A common complex parameter}

For the symmetric inducing characters $|\cdot|^{s/2},|\cdot|^{-s/2}$,
the two parameters are $a=q^{-s/2}$ and $b=q^{s/2}$.
Here $q^s=\exp(s\log q)$, with the real logarithm of $q$.
The ratio $z=q^{-s}$ has period $2\pi i/\log q$.
The ordered pair $(a,b)$ has period $4\pi i/\log q$.
A shift by the smaller period changes both character values by $-1$.
This explains why the ratio does not encode the whole family.

Let $\Omega=\{s\in\mathbb C:|\operatorname{Re}s|<1\}$ and
$B_\Omega=\mathcal O(\Omega)$, treated as an ordinary commutative ring.
There is a ring map
\[
 \phi_q:A\longrightarrow B_\Omega,\qquad
 a\longmapsto e^{-s\log q/2},\qquad b\longmapsto e^{s\log q/2}.
\]
The inverted factors have no zeros in $\Omega$ because
$|q^{-s}|=q^{-\operatorname{Re}s}$ lies strictly between $q^{-1}$ and $q$.
Their reciprocals are holomorphic, which proves that the map is defined.

The involution $\tau(h)(s)=h(-s)$ satisfies $\phi_q\sigma=\tau\phi_q$.
Derived extension of scalars gives a compact-preserving functor
\[
 B_\Omega\otimes_A^{\mathbf L}(-):\mathcal D(A)\longrightarrow\mathcal D(B_\Omega)
\]
and an equivariant version on descent categories.
The latter sends $(P,\theta)$ to the complex with involution
$h\otimes p\mapsto\tau(h)\otimes\theta(p)$.
All finite-level matrices of $E$ and $R$ pull back to holomorphic matrices.

Evaluating at $s$ gives exactly the smooth representation and normalized
operator at $z=q^{-s}$.
On $s\in i\mathbb R$, the previously constructed Hilbert extension is unitary.
No topology on $B_\Omega$ or Hilbert completion is part of this scalar-extension functor.

For a finite set of finite places $S$, allow independent local character
parameters and put $A_S=\bigotimes_{v\in S,\mathbb C}A_v$.
The common symmetric family gives a single map
\[
 \phi_S:A_S\longrightarrow B_\Omega,\qquad
 \bigotimes_v f_v\longmapsto\prod_v\phi_{q_v}(f_v).
\]
It imposes the same variable $s$ at every place.
For $S\subset S'$, scalar extension along $A_S\to A_{S'}$ supplies
compatible coefficient functors, and $\phi_{S'}$ restricts to $\phi_S$.

Thus these maps also define a map from the filtered union of the rings $A_S$.
They give actual base-change functors for independent local parameters.
The compatible Weyl action is the simultaneous swap at all places.
An individual local swap generally does not preserve the common family.

For example, with $q_2=2$ and $q_3=3$, the independent unitary ratios
$(z_2,z_3)=(1,-1)$ do not come from any common $s$.
The first equation forces $s=2\pi i n/\log2$ for an integer $n$.
The second would give $2n\log3=(2m+1)\log2$ for an integer $m$,
after changing signs of the indices if necessary.
Exponentiation contradicts unique factorization of rational numbers.
Thus tensoring independent unitary axes does not impose the global parameter.

After common base change, put $E_{v,\Omega}=B_\Omega\otimes_{A_v}E_v$.
Write $R_{v,\Omega}$ for the scalar extension of $R_v$.
The local semilinear involution $\Theta_v=S_vR_v$ induces
\[
 \Theta_{v,\Omega}(h\otimes f)=\tau(h)\otimes\Theta_v(f).
\]
This formula is balanced because $\phi_{q_v}\sigma_v=\tau\phi_{q_v}$.

It squares to the identity, commutes with $G_v$, and fixes $\mathbf1_v$.
For finite $S$, put $E_{S,\Omega}=\bigotimes_{v\in S,B_\Omega}E_{v,\Omega}$.
The $B_\Omega$-linear operator $\bigotimes_{v\in S}R_{v,\Omega}$
has target the family with every inducing pair swapped.
The semilinear involution on $E_{S,\Omega}$ is instead defined by
\[
 \Theta_S\Bigl(\bigotimes_{v\in S}f_v\Bigr)
     =\bigotimes_{v\in S}\Theta_{v,\Omega}(f_v).
\]

For $S=\varnothing$, set $E_{S,\Omega}=B_\Omega$ and $\Theta_S=\tau$.
For nonempty $S$, moving a scalar between two tensor factors changes
both resulting expressions by its image under $\tau$.
Thus the displayed map is balanced and $\tau$-semilinear.
Applying it twice gives the identity on pure tensors, hence on the whole module.

It commutes with the action of $\prod_{v\in S}G_v$.
Insertion of $\mathbf1_v$ defines compatible maps as $S$ grows.
These maps intertwine $\Theta_S$ and $\Theta_{S'}$ because each
$\Theta_{v,\Omega}$ fixes the inserted vector.
This also holds for the insertion from the empty set, by semilinearity.

The algebraic filtered colimit is the restricted tensor product of these
smooth families over finite places, relative to $\mathbf1_v$.
It therefore inherits a semilinear involution from the maps $\Theta_S$.
It carries the restricted product group relative to $K_v$:
outside finitely many places an element belongs to $K_v$ and fixes $\mathbf1_v$.
Every tensor and every group action is computed at some finite stage.
No infinite product of operator norms is needed for this algebraic construction.

\subsection{The automorphic comparison problem}

The common family above fixes the two global inducing characters to be trivial
before the powers of the absolute value are inserted.
For general inducing characters, their local components must arise from
specified Hecke characters, and the Weyl action must also swap that character pair.
Independent local characters do not by themselves supply these global data.

An analytic coefficient category still requires a choice of sheaf theory,
topology, growth conditions, and derived completion on a complex parameter domain.
The ring category $\mathcal D(B_\Omega)$ specifies none of those choices.
An automorphic target further requires the adelic quotient, central character,
level, archimedean factor, and an Eisenstein functor with analytic domains.
To compare that target with the parameter category, one must construct
the functor on complexes and compare its derived mapping objects and essential image.

Compactness of coefficient complexes does not prove admissibility or
finite generation over local Hecke algebras after induction.
The local smooth family, its coherent Weyl descent, and common-parameter
base change provide the input maps for this comparison.
They do not identify a category of automorphic representations.



\begin{remark}[global normalizations and the pairing comparison]
\label{v4-arith:ic:w6-d:rem:W6C-relation}
The local operator $N_v(s,\chi)=R_z$ divides the raw operator by its
spherical scalar $c(z)$. A global product requires Hecke characters,
compatible spherical insertions, the archimedean and ramified factors,
and an analytic continuation on a specified global family.
For the trivial level-one Eisenstein datum, put
\[
 \Lambda(u)=\pi^{-u/2}\Gamma(u/2)\zeta(u),\qquad
 \xi(u)=\tfrac12u(u-1)\Lambda(u).
\]
The classical scattering factor is $\Lambda(2u-1)/\Lambda(2u)$,
with symmetric induction variable $s=2u-1$.

The entire-completion ratio instead satisfies
\[
 \frac{\xi(2u-1)}{\xi(2u)}
   =\frac{u-1}{u}\frac{\Lambda(2u-1)}{\Lambda(2u)}.
\]
The two phases are different. For a general inducing character,
the global epsilon product retains its character and additive
convention. It is not universally $1$.

An operator, a scalar phase and a Plancherel density are different data.
The Koszul--Petersson transport requires its actual pairing map
and measure identity. If an even measure on the full axis has density
$\tfrac12 w(t)dt$, pushforward under $t\mapsto|t|$ has density
$w(t)dt$ off zero. This folding identity does not determine $w$.
\end{remark}

\begin{definition}[local degeneration and the regular unitary locus]
\label{v4-arith:ic:w6-d:def:exceptional}
For the symmetric unramified family at a finite place, define
\begin{equation}
\label{v4-arith:ic:w6-d:eq:exceptional}
 \Sigma_v^{\mathrm{exc}}=
 \left(1+\frac{2\pi i}{\log q_v}\mathbb Z\right)
 \cup\left(-1+\frac{2\pi i}{\log q_v}\mathbb Z\right).
\end{equation}
These are reducibility parameters. The normalized operator has poles
only on the second set and is a projection at the first set. The
non-exceptional unitary locus is all of $i\mathbb R$. The point $0$
is regular and $1$ is not on the unitary axis.
\end{definition}
\begin{remark}[local degeneration and global residues]
\label{v4-arith:ic:w6-d:rem:exceptional-form}
A residue of a global Eisenstein series is not a pole of every normalized
local operator. Its parameter, global quotient and central character
must be specified. The local kernel/image sequences at $z=q^{\pm1}$
are given by Proposition~\ref{v4-arith:ic:eis:reducibility}. Comparison of
their extension classes with arithmetic deformation data remains a
separate coefficient-compatible construction.
\end{remark}


\section{The regular unramified family}

The normalized unramified intertwiner loses invertibility at two character ratios.
Its cone records the resulting kernels and quotients, including those created by specialization.
To define this cone across both ratios, fix the scalar used to clear the denominator.

Let $F$ be a finite extension of $\mathbb Q_p$ with ring of integers $\mathcal O$,
uniformizer $\varpi$, and residue cardinality $q$.
Set $|\varpi|=q^{-1}$ and normalize additive Haar measure by
$\operatorname{vol}(\mathcal O)=1$.
Put $G=\operatorname{GL}_2(F)$, $K=\operatorname{GL}_2(\mathcal O)$,
and let $B$ be the upper triangular subgroup.
The coefficient ring and ratio are
\[
 R=\mathbb C[a^{\pm1},b^{\pm1}],\qquad z=a/b.
\]

Every coefficient ring in this local-family construction is an ordinary commutative $\mathbb C$-algebra.
It has no analytic topology.
The universal unramified characters are
$\chi_a(x)=a^{v_F(x)}$ and $\chi_b(x)=b^{v_F(x)}$, with $v_F(\varpi)=1$.

Let $E=I_R(\chi_a,\chi_b)$ be the smooth normalized induction.
Its elements are locally constant functions $f:G\to R$ satisfying
\[
 f\left(\begin{pmatrix}u&v\\0&t\end{pmatrix}g\right)
   =|u/t|^{1/2}a^{v_F(u)}b^{v_F(t)}f(g).
\]
The square root is positive real, and $G$ acts by right translation.
Restriction to $K$ identifies its coefficient module with $R\otimes_\mathbb C V$, where
\[
 X=(B\cap K)\backslash K=\mathbb P^1(F),\qquad
 V=C^\infty(X;\mathbb C).
\]

Here $C^\infty$ means locally constant.
Indeed, $G=BK$, and the inducing factor is trivial on $B\cap K$.
For a fixed $g\in G$, compactness partitions $K$ into finitely many pieces
with constant diagonal valuations in an Iwasawa decomposition of $kg$.
Right translation therefore has Laurent polynomial coefficients on every finite space of functions.

This constructs the action over $R$ and proves its compatibility with arbitrary coefficient extension.
Every vector is fixed by an open subgroup of $K$.
Define $E^\mathrm{sw}=I_R(\chi_b,\chi_a)$ in the same way.
These are different $G$-actions on the same coefficient module.

Let $\nu$ be the $K$-invariant probability measure on $X$.
For $N\ge1$, let $V_N$ be the functions pulled back from
$X_N=\mathbb P^1(\mathcal O/\varpi^N\mathcal O)$, and put $V_0=\mathbb C\mathbf1$.
Averaging on reduction fibers defines projections $P_N:V\to V_N$.
Let $P_0$ be global averaging, and define
\[
 H_0=V_0,\qquad H_n=\ker(P_{n-1}|_{V_n})\quad(n\ge1).
\]

Thus $V=\bigoplus_{n\ge0}H_n$ and $V_N=\bigoplus_{n=0}^N H_n$.
Write $H_+=\bigoplus_{n\ge1}H_n$.
The dimensions are $\dim H_0=1$, $\dim H_1=q$, and
$\dim H_n=(q+1)(q-1)q^{n-2}$ for $n\ge2$.
These are $K$-stable spaces.
They need not be $G$-stable.

The unramified integral calculation fixes the normalization.
With
$w=\left(\begin{smallmatrix}0&1\\1&0\end{smallmatrix}\right)$ and
$n(x)=\left(\begin{smallmatrix}1&x\\0&1\end{smallmatrix}\right)$, it gives
\[
 M_zf(g)=\int_F f(wn(x)g)\,dx,\qquad
 c(z)=\frac{1-q^{-1}z}{1-z},\qquad R_z=c(z)^{-1}M_z.
\]
The integral is initially defined at complex specializations with $|z|<1$.
Its rational continuation, on each $V_N$, satisfies
\begin{equation}\label{v4-arith:ic:w10-d:local:eq:normalized}
 R_z|_{H_0}=1,\qquad
 R_z|_{H_n}=z^{n-1}\frac{z-q^{-1}}{1-q^{-1}z}\quad(n\ge1).
\end{equation}
For clarity, the compact-picture calculation underlying this formula is
\[
 M_z=(1+q^{-1})P_0+(z-q^{-1})\sum_{n\ge1}z^{n-1}P_n.
\]

A reduction fiber at depth $n$ has measure
$((q+1)q^{n-1})^{-1}$.
For primitive rows $u,v\in\mathcal O^2$, set $\rho([u],[v])=|\det(u,v)|$.
The radial integral kernel has successive values
$(1+q^{-1})q^nz^n$ at distance $q^{-n}$.
Subtracting successive values gives the displayed projection coefficients.

On $H_m$, $m\ge1$, only the terms $n\ge m$ remain.
Their geometric sum is $(z-q^{-1})z^{m-1}/(1-z)$.
On constants the sum is $c(z)$, which proves \eqref{v4-arith:ic:w10-d:local:eq:normalized}.
The integral commutes with right translation and exchanges the inducing characters
by the substitution $x\mapsto xd/u$ for a diagonal matrix $\operatorname{diag}(u,d)$.
Its intertwining identity consequently continues as a rational identity on finite spaces of vectors.

Set
\[
 \alpha=1-q^{-1}z,\qquad \beta=z-q^{-1},\qquad
 h=(1-q^{-1}z)(1-q^{-1}z^{-1})=z^{-1}\alpha\beta.
\]
The chosen regular map is
\begin{equation}\label{v4-arith:ic:w10-d:local:eq:D}
 D=\alpha R_z=(1-z)M_z:E\longrightarrow E^\mathrm{sw},\qquad
 D|_{R\otimes H_0}=\alpha,\quad
 D|_{R\otimes H_n}=z^{n-1}\beta.
\end{equation}
The equalities with meromorphic operators mean equality of rational families.
The right-hand coefficients define $D$ over all of $R$, including $z=1$.
Multiplication by an additional nonunit would change the cone.
Thus the scalar in \eqref{v4-arith:ic:w10-d:local:eq:D} is part of the construction.

\begin{example}\label{v4-arith:ic:w10-d:local:ex:first}
Let $q=2$ and $N=1$.
The three residue directions have equal measure.
In their characteristic-function basis, put $J=(1)_{i,j=1}^3$.
The projection onto constants is $P=J/3$, so
\[
 D_1=\alpha P+\beta(1-P)
 =\frac12\begin{pmatrix}
 z&1-z&1-z\\1-z&z&1-z\\1-z&1-z&z
 \end{pmatrix},\qquad
 \det D_1=(1-z/2)(z-1/2)^2.
\]

At $z=1/2$ its value is $J/4$, with kernel the two-dimensional sum-zero space.
At $z=2$ its value is $\frac32(1-P)$, with kernel the constants.
At $z=1$ its value is $\frac12\operatorname{id}$.
The matrix with swapped inducing characters is $D_1(z^{-1})$, and
\[
 D_1(z^{-1})D_1(z)=(1-z/2)(1-z^{-1}/2)\operatorname{id}.
\]

Indeed, both matrices preserve constants and their sum-zero complement.
Multiplying the two eigenvalues on either space gives the displayed scalar.
The two distinct kernels determine the two support summands of the general cone.
\end{example}

\begin{proposition}\label{v4-arith:ic:w10-d:local:prop:regular}
The map $D$ is an $R$-linear smooth $G$-intertwiner on all of $E$.
Let $D^\mathrm{sw}:E^\mathrm{sw}\to E$ be the same regular construction with $(a,b)$ replaced by $(b,a)$.
Then
\begin{equation}\label{v4-arith:ic:w10-d:local:eq:composition}
 D^\mathrm{sw} D=h\operatorname{id}_E,\qquad DD^\mathrm{sw}=h\operatorname{id}_{E^\mathrm{sw}}.
\end{equation}
On $H_0$ the coefficient of $D^\mathrm{sw}$ is $\beta/z$.
On $H_n$, $n\ge1$, it is $z^{-n}\alpha$.
After extension to $R[h^{-1}]$, the normalized operator is
$R_z=\alpha^{-1}D$ and its inverse is $(1-q^{-1}z^{-1})^{-1}D^\mathrm{sw}$.
\end{proposition}
\begin{proof}
Every vector has only finitely many $H_n$ components, so \eqref{v4-arith:ic:w10-d:local:eq:D} defines a map on all smooth vectors.
For fixed $g,N$, both sides of the intertwining identity lie in some $R\otimes V_M$.
Their coefficients are Laurent polynomials.
The rational intertwining identity established above proves their equality in the fraction field of $R$, hence in $R$.

Replacing $z$ by $z^{-1}$ gives the stated opposite coefficients.
Their products are $h$ on every summand, which proves \eqref{v4-arith:ic:w10-d:local:eq:composition}.
Inverting $h$ inverts both $\alpha$ and $\beta$, and the normalized composition follows.
\end{proof}

The ordered pair cannot be replaced by its ratio in the representation data.
The central element $\varpi\operatorname{id}$ acts by $ab$.
Replacing $(a,b)$ by $(ta,tb)$ leaves $z$ fixed and tensors the representation
with $\chi_t\circ\det$.
The swap in Proposition~\ref{v4-arith:ic:w10-d:local:prop:regular} always exchanges the two characters.

\section{The cone and its coefficient extensions}

Use cochain complexes, with differentials of degree $1$ and
$P[1]^j=P^{j+1}$ with differential $-d_P$.
The cone of the map between degree-zero representations is
\[
 C=\operatorname{Cone}(D)=
 [\,E\xrightarrow{D}E^\mathrm{sw}\,],
 \qquad C^{-1}=E,\quad C^0=E^\mathrm{sw}.
\]
Let $\mathcal D_{\mathrm{sm}}(G,R)$ denote the derived category of smooth
$R$-linear representations of $G$.
Its objects are complexes, with quasi-isomorphisms inverted.

The action on coefficients is trivial.
The displayed complex is an object of this category.
Its terms are free as $R$-modules.

\begin{theorem}\label{v4-arith:ic:w10-d:local:thm:cone}
The map $D$ is injective over $R$.
Writing $Q_{\mathrm{loc}}=\operatorname{coker} D$, the augmentation $C\to Q_{\mathrm{loc}}[0]$ is a $G$-equivariant quasi-isomorphism.
As an $(R,K)$-module,
\begin{equation}\label{v4-arith:ic:w10-d:local:eq:Q}
 Q_{\mathrm{loc}}\cong (R/(\alpha))\otimes H_0
       \ \oplus\ (R/(\beta))\otimes H_+.
\end{equation}
Its support is the disjoint union of the two divisors $z=q$ and $z=q^{-1}$.
It is annihilated by $h$.
The two support summands in \eqref{v4-arith:ic:w10-d:local:eq:Q} are $G$-stable.

For every commutative $R$-algebra $S$, the complex
\[
 C_S=[\,S\otimes_R E\xrightarrow{1\otimes D}S\otimes_R E^\mathrm{sw}\,]
\]
computes the derived coefficient extension $S\otimes_R^{\mathbf L}C$.
As $S$-modules with $K$-action, its cohomology is
\begin{align*}
 H^{-1}(C_S)&=(\operatorname{ann}_S\alpha)\otimes_\mathbb C H_0
              \oplus(\operatorname{ann}_S\beta)\otimes_\mathbb C H_+,\\
 H^0(C_S)&=(S/\alpha S)\otimes_\mathbb C H_0
              \oplus(S/\beta S)\otimes_\mathbb C H_+.
\end{align*}
There is no cohomology in other degrees.
Here $\alpha,\beta,z$ denote their images in $S$.
\end{theorem}
\begin{proof}
The ring $R$ is a domain.
Every eigenvalue in \eqref{v4-arith:ic:w10-d:local:eq:D} is nonzero, so multiplication is injective on each free summand.
The direct sum remains injective.
Since $z$ is a unit, the image on $R\otimes H_n$ is
$\beta(R\otimes H_n)$ for $n\ge1$.
This proves \eqref{v4-arith:ic:w10-d:local:eq:Q} and the quasi-isomorphism.

The identity
\begin{equation}\label{v4-arith:ic:w10-d:local:eq:bezout}
 \alpha+q^{-1}\beta=1-q^{-2}
\end{equation}
shows that the two defining ideals are comaximal.
Both quotient summands are nonzero, so their supports are exactly the stated divisors.
Scalar multiplication by $h=z^{-1}\alpha\beta$ kills them.
The Chinese remainder idempotents of $R/(\alpha\beta)$ commute with $G$.
Their images therefore give the $G$-stable support summands.
This argument makes no claim that individual $H_n$ are $G$-stable.

A bounded complex of free coefficient modules computes derived extension of scalars.
Tensoring \eqref{v4-arith:ic:w10-d:local:eq:D} with $S$ gives the displayed complex.
The factors $z^{n-1}$ remain units.
Computing kernels and cokernels of scalar multiplication on each direct summand gives the formulas.
\end{proof}

For a coefficient map $S\to T$, the maps $s\otimes f\mapsto\overline s\otimes f$
define a $G$-equivariant chain map $C_S\to C_T$.
They compose strictly, and
$T\otimes_S^{\mathbf L}C_S\simeq C_T$.
These are the specialization maps of this family.
In particular, evaluation at a point $x=(a_0,b_0)\in(\mathbb C^\times)^2$
gives
\[
 C_x=[\,I(\chi_{a_0},\chi_{b_0})
           \xrightarrow{D_x}I(\chi_{b_0},\chi_{a_0})\,].
\]

This statement includes nonflat coefficient maps.
The global equality $\ker D=0$ does not imply $\ker D_x=0$.
The new kernel is $\operatorname{Tor}^R_1(\mathbb C_x,Q_{\mathrm{loc}})$.
Distinct closed points have no specified specialization map between their fibers.
The maps above arise from actual coefficient homomorphisms.

\section{The representations on the exceptional divisors}

The support decomposition permits a full $G$-equivariant description.
Let $J=C^\infty(B\backslash G;\mathbb C)$ be unnormalized induction of the trivial character,
with right translation.
Its constant functions form the trivial representation.
Define the Steinberg representation here by the quotient
\[
 \operatorname{St}=J/\mathbb C\mathbf1.
\]
Only this definition is needed below.

Its compact-picture coefficient space is $V/H_0$.
For any invertible coefficient $t$, write $\chi_t(x)=t^{v_F(x)}$.
Define free $R$-modules with smooth $G$-action by
\[
 L_\alpha=R\otimes(\chi_{q^{1/2}b}\circ\det),\qquad
 L_\beta=(R\otimes_\mathbb C\operatorname{St})\otimes(\chi_{q^{1/2}a}\circ\det).
\]
The first notation means a rank-one free module on which $g$ acts by
$(q^{1/2}b)^{v_F(\det g)}$.
The second combines the indicated scalar with the usual action on $\operatorname{St}$.

\begin{theorem}\label{v4-arith:ic:w10-d:local:thm:representation}
There is an isomorphism of smooth $R$-linear $G$-representations
\begin{equation}\label{v4-arith:ic:w10-d:local:eq:GQ}
 Q_{\mathrm{loc}}\cong L_\alpha/\alpha L_\alpha
       \oplus L_\beta/\beta L_\beta.
\end{equation}
Consequently, in $\mathcal D_{\mathrm{sm}}(G,R)$,
\begin{equation}\label{v4-arith:ic:w10-d:local:eq:koszul}
 C\simeq
 [\,L_\alpha\xrightarrow{\alpha}L_\alpha\,]
 \oplus
 [\,L_\beta\xrightarrow{\beta}L_\beta\,],
\end{equation}
with both source terms in degree $-1$.
The equivalence is specified by the augmentations of both sides to \eqref{v4-arith:ic:w10-d:local:eq:GQ}.
It remains an equivalence after every derived coefficient extension.
\end{theorem}
\begin{proof}
On $\beta=0$, the ordered parameters satisfy $b=qa$.
The target induction has ordered pair $(b,a)=(qa,a)$.
Put $\lambda=q^{1/2}a$.
Its normalized induction factor is
\[
 |u/t|^{1/2}(qa)^{v_F(u)}a^{v_F(t)}
   =\lambda^{v_F(ut)}.
\]

Thus $E^\mathrm{sw}|_{\beta=0}$ is $J$ tensored with $\chi_\lambda\circ\det$.
The image of $D$ there is precisely the constants, since
$\alpha|_{\beta=0}=1-q^{-2}$ is nonzero.
Taking the quotient gives $L_\beta/\beta L_\beta$.
This computation holds over the divisor ring $R/(\beta)$.

On $\alpha=0$, the ordered pair is $(a,b)=(qb,b)$.
The source $E|_{\alpha=0}$ is $J$ tensored with
$\chi_{q^{1/2}b}\circ\det$, by the same induction-law calculation.
The target has ratio $q^{-1}$.
The regular normalized intertwiner from this target to the source is $P_0$,
by \eqref{v4-arith:ic:w10-d:local:eq:normalized}.

Its image is the constant line in the source, and its kernel is $H_+$.
The map $D$ from the source to the target has image $H_+$,
since all its positive-level coefficients are units on this divisor.
Therefore $P_0$ identifies the cokernel with $L_\alpha/\alpha L_\alpha$.
Together with the $G$-stable support decomposition, this proves \eqref{v4-arith:ic:w10-d:local:eq:GQ}.

Multiplication by $\alpha$ on $L_\alpha$ and by $\beta$ on $L_\beta$ is injective.
The two displayed complexes are free coefficient resolutions of the corresponding summands of $Q_{\mathrm{loc}}$.
Their augmentations and $C\to Q_{\mathrm{loc}}$ give the claimed equivalence.
All terms of the two resolutions and of $C$ are coefficient-flat.
Hence these resolutions also compute every derived coefficient extension.
\end{proof}


\begin{proposition}\label{v4-arith:ic:w10-d:local:prop:fibers}
At a complex point $x=(a_0,b_0)$, put $z_0=a_0/b_0$.
If $z_0\notin\{q,q^{-1}\}$, then $C_x$ is acyclic.
At the two exceptional ratios, define
\[
 W_x=\begin{cases}
 \chi_{q^{1/2}b_0}\circ\det,&z_0=q,\\
 \operatorname{St}\otimes(\chi_{q^{1/2}a_0}\circ\det),&z_0=q^{-1}.
 \end{cases}
\]
Then $H^{-1}(C_x)\cong W_x\cong H^0(C_x)$ as smooth $G$-representations.

In the full category $\mathcal D_{\mathrm{sm}}(G,\mathbb C)$,
\[
 C_x\simeq W_x[1]\oplus W_x.
\]
In the compact picture, the kernel and cokernel spaces are $H_0$ at $z_0=q$
and $H_+$ at $z_0=q^{-1}$.
The cokernel means the quotient of the target by the complementary image.
\end{proposition}
\begin{proof}
Outside the two divisors every coefficient of $D_x$ is nonzero, so the map is invertible on $V$.
At $z_0=q$, the $H_0$ coefficient is zero and all positive-level coefficients are nonzero.
At $z_0=q^{-1}$ the reverse holds.
This computes kernels and cokernels as $K$-spaces.

Applying coefficient extension to \eqref{v4-arith:ic:w10-d:local:eq:koszul} identifies the full derived $G$-objects.
On the vanishing divisor the relevant Koszul differential becomes zero.
The other differential becomes an invertible scalar, so its complex is contractible.
The surviving representation is exactly $W_x$.
\end{proof}

\begin{remark}[The central character]\label{v4-arith:ic:w10-d:local:rem:center}
The action on $E$, $E^\mathrm{sw}$, $C$, and $Q_{\mathrm{loc}}$ satisfies
\[
 \pi(t\operatorname{id})=(ab)^{v_F(t)}\operatorname{id}\qquad(t\in F^\times).
\]
One may require this identity strictly on every object and complex.
It defines a smaller abelian category of smooth representations and its derived category.
The cone $C$, its augmentation to $Q_{\mathrm{loc}}$, and its coefficient specializations belong to that category,
with the coefficient character specialized when necessary.

The lifts $L_\alpha$ and $L_\beta$ have central characters
$(qb^2)^{v_F(t)}$ and $(qa^2)^{v_F(t)}$, respectively.
These equal $(ab)^{v_F(t)}$ only after imposing their respective divisor equations.
Thus \eqref{v4-arith:ic:w10-d:local:eq:koszul} is an equivalence in the full category stated in Theorem~\ref{v4-arith:ic:w10-d:local:thm:representation}.
The universal lifts therefore do not prove a splitting inside the
strict universal-central-character category.

At a fixed exceptional
complex point, however, the same actual complex $C_x$ does split in
the strictly specialized category. The next construction supplies a
projective resolution and the two quasi-isomorphisms that establish
this assertion. Universal marked-family compatibility remains a
separate question.
\end{remark}

\section{Strict central character at a complex point}

At an exceptional unramified character ratio, the regular intertwining complex has two nonzero cohomology groups.
A splitting must preserve the central character on every representation and every morphism.
The lattice tree supplies a projective resolution inside that category and removes the degree-two obstruction.

Use the same local field and operator as above. More explicitly, let $F$ be a finite extension of $\mathbb Q_p$.
Write $\mathcal O$ for its ring of integers, $\varpi$ for a uniformizer,
and $q$ for its residue cardinality.
Normalize $v_F(\varpi)=1$ and $|\varpi|=q^{-1}$.

Set $G=\operatorname{GL}_2(F)$, with center $Z=F^\times\operatorname{id}$.
In this fixed-point splitting construction, representations are complex vector spaces with smooth group action:
each vector has an open stabilizer.
All tensor products are algebraic over $\mathbb C$.

Fix $a_0,b_0\in\mathbb C^\times$ with $z_0=a_0/b_0\in\{q,q^{-1}\}$.
Define $\chi_t(u)=t^{v_F(u)}$ for $t\in\mathbb C^\times$.
Let $E_x=I(\chi_{a_0},\chi_{b_0})$ be normalized induction from the upper triangular subgroup $B$.
Thus its locally constant functions satisfy
\[
 f\left(\begin{pmatrix}u&v\\0&t\end{pmatrix}g\right)
   =|u/t|^{1/2}\chi_{a_0}(u)\chi_{b_0}(t)f(g),
 \qquad (\pi(g)f)(h)=f(hg),
\]
with the positive square root.

Write $E_x^{\mathrm{sw}}=I(\chi_{b_0},\chi_{a_0})$.
Both representations have central character
\[
 \zeta(t\operatorname{id})=(a_0b_0)^{v_F(t)}.
\]
Let $\mathcal A_\zeta$ be the abelian category of smooth complex $G$-representations
on which every $t\operatorname{id}$ acts exactly by this scalar.
Its morphisms are $G$-equivariant linear maps.

Kernels and cokernels in smooth representations preserve the stated identity.
Write $\mathcal D^b(\mathcal A_\zeta)$ for its bounded derived category.
Use cochain complexes with differential of degree $1$ and
$P[1]^j=P^{j+1}$ with differential $-d_P$.

The regular intertwiner is the specialization of
$D_z=(1-q^{-1}z)R_z=(1-z)M_z$, with additive measure
$\operatorname{vol}(\mathcal O)=1$ in the raw intertwining integral $M_z$.
The source and target are the ordered inducing pairs and their swaps.
The established compact-picture formula is
\[
 D_z|_{H_0}=1-q^{-1}z,\qquad
 D_z|_{H_n}=z^{n-1}(z-q^{-1})\quad(n\ge1).
\]
Here $H_0$ is the constant line on $\mathbb P^1(F)$.

For $n\ge1$, $H_n$ consists of functions on
$\mathbb P^1(\mathcal O/\varpi^n\mathcal O)$ with zero average on every fiber of reduction to level $n-1$.
At level zero, reduction means global averaging.
These spaces form the algebraic direct sum of all locally constant functions.
The regular map $D_x:E_x\to E_x^{\mathrm{sw}}$ is a smooth $G$-intertwiner.

The complex under consideration is
\begin{equation}\label{v4-arith:ic:w10-d:fixed-center:eq:cone}
 C_x=[\,E_x\xrightarrow{D_x}E_x^{\mathrm{sw}}\,],
 \qquad C_x^{-1}=E_x,\quad C_x^0=E_x^{\mathrm{sw}}.
\end{equation}
Its kernel and cokernel are both isomorphic to
\begin{equation}\label{v4-arith:ic:w10-d:fixed-center:eq:W}
 W_x=\begin{cases}
 \chi_{q^{1/2}b_0}\circ\det,&z_0=q,\\
 \operatorname{St}\otimes(\chi_{q^{1/2}a_0}\circ\det),&z_0=q^{-1},
 \end{cases}
\end{equation}
where $\operatorname{St}=C^\infty(B\backslash G;\mathbb C)/\mathbb C\mathbf1$ uses unnormalized induction of the trivial character.
The operator and these kernel and quotient identifications are the input to the splitting problem.
The following construction applies to every two-term complex in $\mathcal A_\zeta$.

\subsection{Removing the fixed central character}

Set
\[
 \lambda=\begin{cases}q^{1/2}b_0,&z_0=q,\\q^{1/2}a_0,&z_0=q^{-1}.
 \end{cases}
 \qquad \eta=\chi_\lambda\circ\det.
\]
In either case $\lambda^2=a_0b_0$.
The character $\eta$ is smooth, and its restriction to $Z$ equals $\zeta$.
Let $H=G/Z=\operatorname{PGL}_2(F)$, with the quotient topology.

\begin{lemma}\label{v4-arith:ic:w10-d:fixed-center:lem:twist}
Twisting by $\eta^{-1}$ and descending through $G\to H$ gives an exact equivalence
\[
 \mathcal A_\zeta\longrightarrow\operatorname{Rep}_{\mathrm{sm},\mathbb C}(H).
\]
Its inverse inflates an $H$-representation and twists by $\eta$.
Under this equivalence, $W_x$ becomes the trivial representation when $z_0=q$
and the Steinberg quotient when $z_0=q^{-1}$.
\end{lemma}
\begin{proof}
For $V\in\mathcal A_\zeta$, every central element acts trivially on $V\otimes\eta^{-1}$.
Its action therefore factors uniquely through $H$.
If a vector has open stabilizer $U$ in $G$, the image of $UZ$ is open in $H$ and fixes that vector.
Thus the descended action is smooth.

Conversely, inflation of a smooth $H$-action is smooth, and the twist gives central character $\zeta$.
The two constructions are inverse on objects and morphisms.
Their underlying vector-space functors are exact.
Formula~\eqref{v4-arith:ic:w10-d:fixed-center:eq:W} gives the final assertion.
\end{proof}

It remains to construct projective resolutions for smooth representations of $H$.
The quotient by the full center is essential: the stabilizers of lattice classes in $H$ are compact.

\subsection{The lattice tree}

An $\mathcal O$-lattice in $F^2$ is the $\mathcal O$-span of an $F$-basis.
Vertices are homothety classes $[L]$ of such lattices.
Two classes are adjacent if they have representatives satisfying
\[
 \varpi L\subset M\subset L,\qquad
 \dim_{\mathcal O/\varpi\mathcal O}(M/\varpi L)=1.
\]
This is a symmetric relation on classes.
Indeed, the reverse inclusion can be represented by
$\varpi M\subset\varpi L\subset M$.
The group $H$ acts on this graph by its action on lattice classes.

\begin{lemma}\label{v4-arith:ic:w10-d:fixed-center:lem:tree}
This graph is a tree.
Each vertex has $q+1$ neighbors.
\end{lemma}
\begin{proof}
The neighbors of $[L]$ are exactly the lines in $L/\varpi L$, so there are $q+1$.
Fix the root $[L_0]$, where $L_0=\mathcal O e_1\oplus\mathcal O e_2$.
Each class has a unique representative $L\subset L_0$ with $L\not\subset\varpi L_0$.
Multiplication by a unit does not change a lattice, which proves uniqueness.
Choose a primitive vector $u\in L$ and complete it to a basis $u,v$ of $L_0$.

The image of $L$ in $L_0/\mathcal O u$ is an ideal $\varpi^n\mathcal O$.
Subtracting an $\mathcal O u$ component from a lift of its generator gives
\[
 L=\mathcal O u\oplus\mathcal O\varpi^n v,\qquad n\ge0.
\]
The integer $n$ is the valuation of the determinant of this normalized lattice relative to $L_0$.
It is independent of the chosen basis and will be its height.
Height zero gives the root.

For $n>0$, the line generated by $\varpi^n v$ in $L/\varpi L$
gives the neighboring lattice
$\mathcal O\varpi u\oplus\mathcal O\varpi^n v$.
After normalization, its height is $n-1$.
Every other line is generated by $u+c\varpi^n v$ for one residue $c\in\mathcal O/\varpi\mathcal O$.
Its neighboring lattice is
\[
 \mathcal O(u+c\varpi^n v)\oplus\mathcal O\varpi^{n+1}v,
\]
which is normalized and has height $n+1$.

Every neighbor of the root has height one.
Thus every positive-height vertex has exactly one neighbor of lower height.
Repeated lower-height steps connect it to the root.
A cycle would have a maximal-height vertex with two neighbors of lower height, which is impossible.
\end{proof}

An element of $H$ can exchange the endpoints of an edge.
Insert a midpoint in each edge to obtain the barycentric subdivision $T$.
Orient each new edge from its lattice vertex to its midpoint.
The $H$-action preserves these two vertex types and therefore preserves the chosen orientations.

Let $v_0=[L_0]$ and $v_1=[\mathcal O e_1\oplus\mathcal O\varpi e_2]$.
Let $m_0$ be the midpoint between them.
The stabilizer $K$ of $v_0$ is the image of $\operatorname{GL}_2(\mathcal O)$ in $H$.
It is compact and open.

The stabilizer $I$ of the incidence $(v_0,m_0)$ fixes both $v_0$ and $v_1$.
In $K$, it is the inverse image of the stabilizer of a line modulo $\varpi$, hence compact and open.
The class of
\[
 s=\begin{pmatrix}0&1\\\varpi&0\end{pmatrix}
\]
exchanges $v_0$ and $v_1$ and satisfies $s^2=1$ in $H$.
The midpoint stabilizer is $J=I\sqcup sI$ and is compact and open.

The action is transitive on lattice vertices, on midpoints, and on the oriented edges of $T$.
The first assertion follows by sending one lattice basis to another.
Transitivity on incident pairs follows because $K$ acts transitively on lines modulo $\varpi$.

\begin{example}[The first ball when $q=2$]\label{v4-arith:ic:w10-d:fixed-center:ex:star}
The first ball about $v_0$ is a star with three leaves $v_1,v_2,v_3$.
Write $m_i$ for the midpoint of $v_0v_i$.
The subdivision has seven vertices and six oriented edges
$e_i=(v_0,m_i)$ and $f_i=(v_i,m_i)$, for $1\le i\le3$.
Its boundary map and augmentation are
\[
 \partial e_i=[m_i]-[v_0],\qquad
 \partial f_i=[m_i]-[v_i],\qquad
 \epsilon[v]=1.
\]
Define a linear map $h$ from vertices to edge chains by
\[
 h[v_0]=0,\qquad h[m_i]=e_i,\qquad h[v_i]=e_i-f_i.
\]

Then
\[
 \partial h=\operatorname{id}-[v_0]\epsilon,\qquad h\partial=\operatorname{id}.
\]
Both identities follow by applying the displayed formulas to each basis vector.
Thus $\partial$ has rank six and its image is $\ker\epsilon$.
This finite ball is not $H$-stable.
The same finite-path calculation proves exactness for arbitrary finite chains in the entire tree.
\end{example}

\subsection{Projective cellular chains}

For a compact open subgroup $U\subset H$, let $\operatorname{Rep}_{\mathrm{sm},\mathbb C}(U)$ denote its smooth complex representations.
An object is projective if every morphism from it to a quotient lifts through the quotient map.
The following property depends on the coefficient field $\mathbb C$.

\begin{lemma}\label{v4-arith:ic:w10-d:fixed-center:lem:compact}
Every smooth complex representation of $U$ is projective in
$\operatorname{Rep}_{\mathrm{sm},\mathbb C}(U)$.
\end{lemma}
\begin{proof}
The orbit of any vector under $U$ is finite because its stabilizer is open and $U$ is compact.
The orbit spans a finite-dimensional $U$-submodule.
A basis of this submodule has a common open stabilizer.
Its normal core is an open normal subgroup, so the action factors through a finite quotient of $U$.

For an invariant subspace, average a linear projection over the finite quotient to obtain an equivariant projection.
Its kernel is an invariant complement.
Induction on dimension therefore decomposes this submodule into irreducibles.
Hence every smooth $U$-module is a sum of finite-dimensional irreducible submodules.

Choose a maximal direct sum of such submodules.
If it does not contain the whole module, a finite-dimensional submodule contributes a simple summand outside it.
That simple summand has zero intersection with the maximal direct sum, a contradiction.
Thus the whole module is a direct sum of finite-dimensional irreducibles.

Each such irreducible $\tau$ is projective.
To verify the lifting property explicitly, let $Y\twoheadrightarrow Z$ be an epimorphism and let $f:\tau\to Z$ be $U$-equivariant.
Choose lifts in $Y$ of the images of a basis of $\tau$ and extend them to a linear map $\ell:\tau\to Y$.
The $U$-orbits of these finitely many lifts span a finite-dimensional submodule $Y_0$.

Choose an open normal subgroup $N\subset U$ acting trivially on both $Y_0$ and $\tau$.
Then
\[
 \widetilde f(v)=\frac1{|U/N|}\sum_{u\in U/N}u\,\ell(u^{-1}v)
\]
is well-defined, $U$-equivariant, and lifts $f$.
Only finite-dimensional smooth data are averaged.
Finally, morphisms out of a direct sum can be lifted on its summands and then assembled.
A direct sum of projectives is therefore projective.
\end{proof}

For $M\in\operatorname{Rep}_{\mathrm{sm},\mathbb C}(U)$, define compact induction
$\operatorname{c\!\operatorname{-}Ind}_U^H M$ by finite sums of symbols $[g,m]$, linear in $m$, with relations
\[
 [gu,m]=[g,um]\qquad(g\in H,\ u\in U).
\]
The group acts by $h[g,m]=[hg,m]$.
This is the algebraic finite-support model on $H/U$, valid because $U$ is open.
It is smooth: the stabilizer of $[g,m]$ contains $gU_mg^{-1}$, where $U_m$ is an open stabilizer of $m$ in $U$.
A finite sum has the intersection of finitely many such open stabilizers.

The map determined by $[g,m]\mapsto g f(m)$ gives the adjunction
\[
 \operatorname{Hom}_H(\operatorname{c\!\operatorname{-}Ind}_U^H M,V)\cong\operatorname{Hom}_U(M,V|_U).
\]
Its inverse restricts a morphism to the symbols $[1,m]$.
Restriction is exact, so compact induction takes projectives to projectives.
Lemma~\ref{v4-arith:ic:w10-d:fixed-center:lem:compact} applies to every compact open cell stabilizer of $T$.

For a discrete set $Y$, write $C_c(Y;\mathbb C)$ for its finitely supported functions,
with basis vectors denoted $[y]$.
Let $T_0$ and $T_1$ be the vertices and the chosen oriented edges of $T$.
For a smooth $H$-representation $V$, define
\[
 P_1(V)=C_c(T_1;\mathbb C)\otimes V,\qquad
 P_0(V)=C_c(T_0;\mathbb C)\otimes V,
\]
with diagonal action $g([c]\otimes v)=[gc]\otimes gv$.
For an oriented edge $e=(u,m)$, define
\[
 \partial([e]\otimes v)=([m]-[u])\otimes v,
 \qquad \epsilon([u]\otimes v)=v.
\]
Both maps commute with $H$, and $\epsilon\partial=0$.

\begin{theorem}\label{v4-arith:ic:w10-d:fixed-center:thm:resolution}
For every smooth complex $H$-representation $V$, the sequence
\begin{equation}\label{v4-arith:ic:w10-d:fixed-center:eq:resolution}
 0\longrightarrow P_1(V)\xrightarrow{\partial}P_0(V)
 \xrightarrow{\epsilon}V\longrightarrow0
\end{equation}
is exact and its first two terms are projective smooth $H$-representations.
Consequently,
\[
 \operatorname{Ext}^j_{\operatorname{Rep}_{\mathrm{sm},\mathbb C}(H)}(V,U)=0
 \qquad(j\ge2)
\]
for all smooth $V,U$.
\end{theorem}
\begin{proof}
Fix a root vertex $o$ of $T$.
Let $p_v$ be the finite edge chain along the unique path from $o$ to $v$.
Its coefficient is $-1$ wherever the path opposes the chosen edge orientation.
Then $\partial p_v=[v]-[o]$.

The map $[v]\otimes w\mapsto p_v\otimes w$ gives
$\partial h=\operatorname{id}-([o]\otimes-)\epsilon$.
Thus every chain of augmentation zero is a boundary.
The augmentation is surjective because $w=\epsilon([o]\otimes w)$.
A nonzero finite-support edge cycle would have a leaf in the finite subtree spanned by its support.

At that leaf, its boundary equals the nonzero coefficient of its unique incident supported edge, up to sign.
This is impossible, so $\partial$ is injective.
These arguments also apply to vector coefficients and use only finite sums.
The root-path contraction proves exactness and is not claimed equivariant.

For a chosen cell $c$ with stabilizer $U$, the map
\[
 \operatorname{c\!\operatorname{-}Ind}_U^H(V|_U)\longrightarrow C_c(Hc;\mathbb C)\otimes V,
 \qquad [g,v]\longmapsto[gc]\otimes gv
\]
is an equivariant isomorphism.
The relation $[gu,v]=[g,uv]$ verifies that it is well-defined.
Its inverse sends $[gc]\otimes w$ to $[g,g^{-1}w]$ and is independent of the representative $g$.
The three cell orbits therefore give
\[
 P_1(V)\cong\operatorname{c\!\operatorname{-}Ind}_I^H(V|_I),\qquad
 P_0(V)\cong\operatorname{c\!\operatorname{-}Ind}_K^H(V|_K)\oplus\operatorname{c\!\operatorname{-}Ind}_J^H(V|_J).
\]
They are smooth projectives by Lemma~\ref{v4-arith:ic:w10-d:fixed-center:lem:compact} and the compact-induction adjunction.
Applying $\operatorname{Hom}(-,U)$ to this projective resolution computes $\operatorname{Ext}$ and leaves no terms in degrees $j\ge2$.
\end{proof}

The constant-coefficient resolution above needs no finite generation or admissibility hypothesis.
Its projective terms can have infinite multiplicity.
Schneider--Stuhler develop cellular resolutions with coefficient
systems of compact-subgroup invariants in
\emph{Representation theory and sheaves on the Bruhat--Tits building}.
See Section~II.2, pp.~121--122, and Theorem~II.3.1, p.~123.
Here the rank-one tree and constant coefficients give the stated resolution directly.

\subsection{The obstruction and the splitting maps}

Lemma~\ref{v4-arith:ic:w10-d:fixed-center:lem:twist} transports the resolution and projectivity into $\mathcal A_\zeta$.
For $V\in\mathcal A_\zeta$, apply \eqref{v4-arith:ic:w10-d:fixed-center:eq:resolution} to its descended twist $V\otimes\eta^{-1}$,
then inflate and tensor every term and map with $\eta$.
This gives a projective resolution
\[
 0\longrightarrow P_1^\zeta(V)\xrightarrow{d}P_0^\zeta(V)
 \xrightarrow{\epsilon}V\longrightarrow0
\]
inside $\mathcal A_\zeta$.
In particular, $\operatorname{Ext}^j_{\mathcal A_\zeta}(V,U)=0$ for all $j\ge2$.
Every central element acts by $\zeta$ on both projective terms.

\begin{theorem}\label{v4-arith:ic:w10-d:fixed-center:thm:split}
Let $A\xrightarrow{D}B$ be a morphism in $\mathcal A_\zeta$.
Place $A$ in degree $-1$ and $B$ in degree $0$, and set
$U=\ker D$, $V=\operatorname{coker} D$.
There is an isomorphism
\[
 [\,A\xrightarrow{D}B\,]\simeq U[1]\oplus V
 \qquad\text{in }\mathcal D^b(\mathcal A_\zeta).
\]
It can be represented by two explicit quasi-isomorphisms through the cellular projective resolution of $V$.
\end{theorem}
\begin{proof}
Write $P_1\xrightarrow{d}P_0\xrightarrow{\epsilon}V$ for the resolution just constructed.
Let $p:B\twoheadrightarrow V$ and $i:U\hookrightarrow A$ be the quotient and kernel maps.
Since $P_0$ is projective, choose a morphism $f_0:P_0\to B$ with $pf_0=\epsilon$.
The identity $pf_0d=\epsilon d=0$ shows that $f_0d$ takes values in $\operatorname{im}D$.

Projectivity of $P_1$ lifts that map through the epimorphism $A\twoheadrightarrow\operatorname{im}D$.
Thus there is $f_1:P_1\to A$ satisfying
\begin{equation}\label{v4-arith:ic:w10-d:fixed-center:eq:lifts}
 Df_1=f_0d.
\end{equation}
All objects, quotient maps, and lifts lie in $\mathcal A_\zeta$.

Let $P=[P_1\xrightarrow{d}P_0]$, in degrees $-1,0$.
Equation~\eqref{v4-arith:ic:w10-d:fixed-center:eq:lifts} defines a chain map
\[
 \Phi:U[1]\oplus P\longrightarrow[\,A\xrightarrow{D}B\,],
 \qquad
 \Phi^{-1}(u,v)=i(u)+f_1(v),\quad \Phi^0=f_0.
\]
The superscripts denote cochain degrees.
Since $d$ is injective, $H^{-1}(U[1]\oplus P)=U$ and $H^0(U[1]\oplus P)=V$.
The map $\Phi$ induces the identity on these groups, using $pf_0=\epsilon$.

It is therefore a quasi-isomorphism.
The augmentation of $P$ gives another quasi-isomorphism
\[
 \Psi:U[1]\oplus P\longrightarrow U[1]\oplus V,
 \qquad \Psi^{-1}(u,v)=u,\quad \Psi^0=\epsilon.
\]
The roof
\[
 [\,A\xrightarrow{D}B\,]\xleftarrow{\ \Phi\ }U[1]\oplus P
 \xrightarrow{\ \Psi\ }U[1]\oplus V
\]
proves the asserted derived isomorphism.
The compact averaging construction in Lemma~\ref{v4-arith:ic:w10-d:fixed-center:lem:compact}, followed by compact induction,
constructs the required lifts after choices of irreducible summands and linear lifts.
No averaging over the noncompact group is used.
\end{proof}

The associated obstruction is the class of the exact two-extension
\[
 0\longrightarrow U\xrightarrow{i}A\xrightarrow{D}B\xrightarrow{p}V\longrightarrow0
\]
in $\operatorname{Ext}^2_{\mathcal A_\zeta}(V,U)$.
Equivalently, it is the connecting morphism $V\to U[2]$ in the cohomological truncation triangle.
Theorem~\ref{v4-arith:ic:w10-d:fixed-center:thm:resolution} makes this group zero, and the roof above realizes the splitting by actual maps.
The construction proves existence and depends on the choices of lifts.
No canonical splitting or compatibility between different parameters is asserted.

\begin{proposition}[The exceptional fibers]\label{v4-arith:ic:w10-d:fixed-center:prop:application}
For the exceptional pair fixed in \eqref{v4-arith:ic:w10-d:fixed-center:eq:cone}, there is an isomorphism
\[
 C_x\simeq W_x[1]\oplus W_x
 \qquad\text{in }\mathcal D^b(\mathcal A_\zeta).
\]
Its degree-two obstruction vanishes in the strictly fixed-central-character category.
\end{proposition}
\begin{proof}
The kernel and cokernel identifications \eqref{v4-arith:ic:w10-d:fixed-center:eq:W} are morphisms inside $\mathcal A_\zeta$.
Apply Theorem~\ref{v4-arith:ic:w10-d:fixed-center:thm:split} to the actual map $D_x$.
The resulting roof uses the projective resolution twisted by the stated $\eta$,
so every term has central character $\zeta$.
\end{proof}

This is a pointwise statement over $\mathbb C$.
The compact projectivity proof uses division by finite group orders and does not establish a theorem over arbitrary coefficient rings.
The character $\eta$ is fixed at the chosen complex point.
Universal-central-character splitting over the ordered-character ring
requires a separate construction. So do canonical Weyl-compatible
pointwise splittings and a comparison with global Eisenstein categories.


\section{Normal equations and specialization maps}
The opposite operator provides a direct comparison of these two cohomology groups.
Put $T=R/(h)$.
Equation~\eqref{v4-arith:ic:w10-d:local:eq:composition} gives a well-defined map
\[
 \overline {D^\mathrm{sw}}:\operatorname{coker}(D_T)\longrightarrow\ker(D_T),
 \qquad [y]\longmapsto D^\mathrm{sw}_Ty.
\]
It is an isomorphism of smooth $T$-linear $G$-representations.

To prove this, decompose $T$ by the two comaximal ideals.
On $\alpha=0$, its surviving coefficient is $\beta/z=1-q^{-2}$ on $H_0$.
On $\beta=0$, its surviving coefficients are $z^{-n}\alpha$ on $H_n$, $n\ge1$.
All are units in the corresponding divisor ring.
The same isomorphism holds after further coefficient extension from $T$.

There is also a comparison determined by a chosen equation of one divisor.
For $t=\alpha$ or $t=\beta$, let $\overline R=R/(t)$.
For $x\in\ker D_{\overline R}$, choose a lift $\widetilde x\in E$.
The vector $D\widetilde x$ is divisible by $t$ in the free module $E^\mathrm{sw}$.

Define
\begin{equation}\label{v4-arith:ic:w10-d:local:eq:bockstein}
 \delta_t:\ker D_{\overline R}\longrightarrow\operatorname{coker} D_{\overline R},
 \qquad x\longmapsto[\,D\widetilde x/t\,].
\end{equation}
Changing the lift by $tu$ changes the quotient by $Du$, so the map is well-defined.
Equivariance follows by applying $D\pi(g)=\pi^\mathrm{sw}(g)D$ to the chosen lift.
This is the connecting map for the coefficient sequence
$0\to R/(t)\xrightarrow{t}R/(t^2)\to R/(t)\to0$ applied to $C$.
It also identifies the kernel calculated from the free resolution of $Q_{\mathrm{loc}}$
with its $t$-torsion using the resolution $[R\xrightarrow{t}R]$ of $R/(t)$.

The map $\delta_\alpha$ has coefficient $1$ on $H_0$.
The map $\delta_\beta$ has coefficient $z^{n-1}$ on $H_n$, $n\ge1$, evaluated on $\beta=0$.
Thus both are isomorphisms.
They are compatible with subsequent extension along the respective divisor.

They satisfy
\[
 \delta_t\overline {D^\mathrm{sw}}=(h/t)|_{t=0}\,\operatorname{id},
 \qquad
 (h/\alpha)|_{\alpha=0}=1-q^{-2},\quad
 (h/\beta)|_{\beta=0}=q(1-q^{-2}).
\]
Replacing the local equation $t$ by $ut$, with $u$ a unit, multiplies $\delta_t$ by $u^{-1}$ on the divisor.
The normal equation is therefore part of this specialization comparison.

\section{Finite levels and matrix factorizations}

The finite spaces $V_N$ retain the integral's compact-level refinement.
Set $E_N=R\otimes V_N$ and $E_N^\mathrm{sw}=R\otimes V_N$, with their respective induced compact actions.
Restriction gives $D_N:E_N\to E_N^\mathrm{sw}$ and $D_N^\mathrm{sw}:E_N^\mathrm{sw}\to E_N$.
Inclusions $V_N\hookrightarrow V_{N+1}$ and averaging maps $P_N:V_{N+1}\to V_N$
commute with both operators because they preserve each $H_n$.

They therefore induce maps of the cones and of their coefficient extensions.
The inclusions give
\[
 C=\operatorname*{colim}_N C_N,\qquad C_N=[E_N\xrightarrow{D_N}E_N^\mathrm{sw}].
\]
Filtered colimits of modules are exact, so kernels, cokernels, and the displayed coefficient extensions commute with this colimit.

For $N\ge1$, let $K_N=1+\varpi^NM_2(\mathcal O)$, and set $K_0=K$.
Then $E^{K_N}=E_N$ and $(E^\mathrm{sw})^{K_N}=E_N^\mathrm{sw}$.
Compact averaging is defined over every $\mathbb C$-algebra and is exact on smooth representations.
Indeed, any finite set of vectors is fixed by an open subgroup, and averaging reduces to a finite sum with invertible denominator.

It follows that $C_N=C^{K_N}$ and $\operatorname{coker} D_N=Q_{\mathrm{loc}}^{K_N}$.
Full $G$-translation can move $E_N$ into a larger $E_M$.
The finite complexes carry $K$-actions, whereas their colimit carries the full $G$-action.

A finite-free matrix factorization of $h\in R$ consists of finite-free modules $M_0,M_1$
and maps $d_0:M_0\to M_1$, $d_1:M_1\to M_0$ satisfying
$d_1d_0=h\operatorname{id}$ and $d_0d_1=h\operatorname{id}$.
Thus $(D_N,D_N^\mathrm{sw})$ is a finite-free $K$-equivariant matrix factorization of $h$.
Its rank is $\dim V_N$ in each parity.
The full pair $(D,D^\mathrm{sw})$ has infinite-rank free terms.

\begin{proposition}\label{v4-arith:ic:w10-d:local:prop:mf}
Each finite-level matrix factorization $(D_N,D_N^\mathrm{sw})$ is contractible through $K$-equivariant maps.
Its cokernel is a projective $R/(h)$-module.
The ordinary two-term cone $C_N$ is nevertheless nonzero in the derived category for every $N\ge0$.
\end{proposition}
\begin{proof}
Put $c=1-q^{-2}$.
Write $d_0=D_N$, $d_1=D_N^\mathrm{sw}$.
A contraction consists of maps $s_0:M_0\to M_1$ and $s_1:M_1\to M_0$
with $ds+sd=\operatorname{id}$ on the two-parity module.
On $H_0$ take
\[
 s_1=c^{-1},\qquad s_0=q^{-1}z/c.
\]

Then $s_1d_0+d_1s_0=(\alpha+q^{-1}\beta)/c=1$.
On $H_n$, $n\ge1$, take
\[
 s_1=q^{-1}z^{-(n-1)}/c,\qquad s_0=z^n/c.
\]
Here $s_1d_0+d_1s_0=(q^{-1}\beta+\alpha)/c=1$ as well.
All coefficients commute, so the identity also holds on the other parity.
These are Laurent coefficients and preserve the $K$-stable decomposition.

The ring $R/(h)$ is $R/(\alpha)\times R/(\beta)$ by \eqref{v4-arith:ic:w10-d:local:eq:bezout}.
The cokernel is free of rank one on the first component and rank $\dim V_N-1$ on the second.
It is therefore projective over their product.
As an $R$-module it is nonzero, since the first component is nonzero.
This cokernel is $H^0(C_N)$, which proves the last assertion.
\end{proof}

The contraction above is asserted only for the coefficient modules with $K$-action.
Its diagonal formulas do not establish a contraction commuting with the full $G$-action.
A matrix factorization uses both opposite maps and satisfies $d^2=h$.
The ordinary cone uses only $D$ and has the nonzero cohomology \eqref{v4-arith:ic:w10-d:local:eq:Q}.
These are distinct homological constructions.

The regular intertwiner therefore defines a smooth equivariant degeneration complex on the entire ordered-character torus.
Its cohomology consists of a character family on one divisor and a Steinberg quotient family on the other.
Its derived fibers and their specialization comparisons are given above.
A category comparison for global Eisenstein representations still requires a global target and a functor
whose derived mapping objects and essential image are determined.
The local degeneration complex supplies no such equivalence by itself.


\section{Local Yoneda classes and finite exterior products}
\label{v4-arith:ic:w10-d:X1}

A local nonsplit principal series defines a degree-one extension between
its Steinberg quotient and its constant subrepresentation. At two distinct
finite places, the exterior product of these extensions has degree two.
An infinite restricted tensor introduces a second condition: its
distinguished local vectors must be fixed by the prescribed compact
subgroups at almost every place. These two requirements distinguish a
finite exterior product of Steinberg representations from an adelic
extension of principal series.

\subsection{The local sequence and its invariant vectors}

Fix the coefficient field $\mathbb C$. For a finite prime $p$, put
$G_p=GL_2(\mathbb Q_p)$, $K_p=GL_2(\mathbb Z_p)$, and let $B_p$ be the upper triangular
Borel subgroup. Give these groups their usual $p$-adic topologies.
Write $X_p=B_p\backslash G_p\simeq\mathbb P^1(\mathbb Q_p)$, with its compact
quotient topology and right $G_p$-action. The identification sends a matrix
to the line spanned by its bottom row.

Let
\[
 E_p=C^{\mathrm{lc}}(X_p,\mathbb C),\qquad
 (\rho_p(g)f)(x)=f(xg).
\]
Here $C^{\mathrm{lc}}$ means locally constant functions, without an
analytic completion. This is unnormalized smooth induction of the trivial
character of $B_p$. With normalized induction and $|p|_p=p^{-1}$, its
ordered inducing characters are $|\cdot|_p^{-1/2}$ and
$|\cdot|_p^{1/2}$. Thus, if symmetric induction is written with characters
$(|\cdot|_p^{s/2},|\cdot|_p^{-s/2})$, this sequence occurs at $s=-1$.

The constants give the exact sequence
\begin{equation}\label{v4-arith:ic:w10-d:adelic:eq:local}
 0\longrightarrow\mathbb C\xrightarrow{i_p}E_p
   \longrightarrow\operatorname{St}_p\longrightarrow0,
 \qquad \operatorname{St}_p=E_p/\mathbb C\mathbf1_p.
\end{equation}
This is the usual rank-one Steinberg quotient.
See Casselman, \emph{Introduction to the Theory of Admissible
Representations of $p$-Adic Reductive Groups}, Section~8.1, p.~70. Every term has trivial central character.

We use $\operatorname{Rep}^{\mathrm{sm}}_{\mathbb C}(G_p)$ for the abelian category of all smooth complex
representations. Its full subcategory on which the center acts trivially
is identified with $\operatorname{Rep}^{\mathrm{sm}}_{\mathbb C}(PGL_2(\mathbb Q_p))$. Extension groups below specify
which of these categories is used.

\begin{lemma}\label{v4-arith:ic:w10-d:adelic:lem:local}
The sequence~\eqref{v4-arith:ic:w10-d:adelic:eq:local} is nonsplit. Moreover,
\[
 \operatorname{Hom}_{G_p}(E_p,\mathbb C)=0,\qquad
 \operatorname{Hom}_{G_p}(\operatorname{St}_p,\mathbb C)=0,\qquad
 \operatorname{St}_p^{K_p}=0.
\]
The same statements hold after passage to $PGL_2(\mathbb Q_p)$.
\end{lemma}
\begin{proof}
Consider a $G_p$-invariant linear functional $\ell:E_p\to\mathbb C$.
Projective $p$-adic balls are translates of one another under $G_p$.
Affine transformations relate the balls in $\mathbb Q_p$.
Inversion relates them to the balls containing infinity. A ball is a disjoint union
of $p$ smaller balls.

Consequently the common value $c$ of $\ell$ on
ball indicators satisfies $c=pc$, so $c=0$. Locally constant functions
on $X_p$ are finite linear combinations of these indicators, since $X_p$
is compact and the balls form a clopen basis. Hence $\ell=0$.
A splitting of~\eqref{v4-arith:ic:w10-d:adelic:eq:local} would provide such a functional with
$\ell(\mathbf1_p)=1$. A morphism $\operatorname{St}_p\to\mathbb C$ would pull back to a
morphism $E_p\to\mathbb C$, proving the first two assertions.

Every line in $\mathbb Q_p^2$ has a primitive integral representative that extends
to an integral basis. Thus $K_p$ acts transitively on $X_p$ and
$E_p^{K_p}=\mathbb C\mathbf1_p$. Invariants under a compact open subgroup
are exact for smooth complex representations. Indeed, normalized Haar
averaging of a smooth vector is a finite sum, and averaging a lift of an
invariant vector gives an invariant lift.

Applying invariants to
\eqref{v4-arith:ic:w10-d:adelic:eq:local} gives $\operatorname{St}_p^{K_p}=0$. This averaging argument is also
Casselman, Proposition~2.1.7, p.~20.
All actions factor through $PGL_2(\mathbb Q_p)$, and the same proof applies there.
\end{proof}

Write $\alpha_p$ for the Yoneda class of~\eqref{v4-arith:ic:w10-d:adelic:eq:local}. It is a
specified nonzero element of $\operatorname{Ext}^1(\operatorname{St}_p,\mathbb C)$. No residue operator is
needed to define it, and no dimension assertion about that group is
required below.

\subsection{Two finite places and the exterior degree}

Let $p\ne q$ be finite primes. External tensor products over $\mathbb C$ are
written $\boxtimes$ and carry the separate actions of the two groups.
Tensoring the two local sequences gives the exact sequence
\begin{equation}\label{v4-arith:ic:w10-d:adelic:eq:two}
\begin{split}
 0\longrightarrow\mathbb C\boxtimes\mathbb C
 &\xrightarrow{d_2}
 (E_p\boxtimes\mathbb C)\oplus(\mathbb C\boxtimes E_q)
 \xrightarrow{d_1}E_p\boxtimes E_q\\
 &\longrightarrow\operatorname{St}_p\boxtimes\operatorname{St}_q\longrightarrow0,
\end{split}
\end{equation}
where
\[
 d_2=(i_p\boxtimes1,-1\boxtimes i_q),\qquad
 d_1(x,y)=(1\boxtimes i_q)x+(i_p\boxtimes1)y.
\]
Unlike the cochain indexing of the local cone, the following totals use homological degrees and differentials of degree $-1$.

These signs are the tensor differential for homological complexes,
$\partial(x\otimes y)=\partial x\otimes y+(-1)^{\deg x}x\otimes\partial y$.
In particular $d_1d_2=0$. The two complexes
$[\mathbb C\to E_p]$ and $[\mathbb C\to E_q]$, in homological degrees $1,0$,
have homology only in degree zero, with values $\operatorname{St}_p$ and $\operatorname{St}_q$.
Tensoring over a field preserves this assertion. It proves the exactness
of~\eqref{v4-arith:ic:w10-d:adelic:eq:two}.

The sequence~\eqref{v4-arith:ic:w10-d:adelic:eq:two} is a Yoneda two-extension, not a short exact
sequence with these endpoints. Its class is the exterior Yoneda product
of $\alpha_p$ and $\alpha_q$, with the displayed tensor sign convention.
The following argument proves that this class is nonzero without
replacing either local representation by a finite-dimensional model.

\begin{proposition}\label{v4-arith:ic:w10-d:adelic:prop:finite}
Let $S$ be a nonempty finite set of finite primes, $n=|S|$, and
$\operatorname{St}_S=\boxtimes_{p\in S}\operatorname{St}_p$. Put $G_S=\prod_{p\in S}G_p$.
In $\operatorname{Rep}^{\mathrm{sm}}_{\mathbb C}(G_S)$ one has
\[
 \operatorname{Ext}^r_{G_S}(\operatorname{St}_S,\mathbb C)=0\quad(0\le r<n).
\]
The exterior product of the classes $\alpha_p$ is nonzero in
$\operatorname{Ext}^n_{G_S}(\operatorname{St}_S,\mathbb C)$. It is represented by the total tensor
of the $n$ sequences~\eqref{v4-arith:ic:w10-d:adelic:eq:local}.
Both assertions also hold in the category where each local center acts
trivially, namely $\operatorname{Rep}^{\mathrm{sm}}_{\mathbb C}(\prod_{p\in S}PGL_2(\mathbb Q_p))$.
\end{proposition}
\begin{proof}
For either of the groups $G_p$ and $PGL_2(\mathbb Q_p)$, compact induction
$\operatorname{c-Ind}_K^G\mathbb C$ from a compact open subgroup is projective. A $G$-map from this compact induction to $V$ is determined
by the image of the indicator of $K$.
That image can be any element of $V^K$.
Thus its Hom functor is $V\mapsto V^K$, which is exact by averaging.

Every smooth vector is fixed by a compact open subgroup. Taking one such compact induction for each vector gives enough
projectives. Resolutions can therefore use sums of these compact inductions.

Choose such a resolution $P_{p,\bullet}\to\operatorname{St}_p$, with homological
indices at least zero. Products of compact open subgroups are compact
open, and
\[
 \left(\operatorname{c-Ind}_K^G\mathbb C\right)\boxtimes
 \left(\operatorname{c-Ind}_L^H\mathbb C\right)
 \simeq\operatorname{c-Ind}_{K\times L}^{G\times H}\mathbb C.
\]
Both sides are finite-support functions on the product of the relevant
discrete coset spaces. Hence the total external tensor of the chosen
resolutions is a projective resolution of $\operatorname{St}_S$.

Let $C_{p,\bullet}=(P_{p,\bullet})_{G_p}$ be its coinvariant complex
(or use $PGL_2(\mathbb Q_p)$ in the fixed-central-character category).
For any representation $V$, a map $V\to\mathbb C$ factors uniquely through
$V_G$. It follows that
\begin{equation}\label{v4-arith:ic:w10-d:adelic:eq:coinvariant}
 \operatorname{Ext}^r_G(\operatorname{St}_p,\mathbb C)
 \simeq\operatorname{Hom}_{\mathbb C}(H_r(C_{p,\bullet}),\mathbb C).
\end{equation}
Here the algebraic dual is exact over the field $\mathbb C$.

No topological
or continuous dual is used. Coinvariants of an external tensor for a
product group are the tensor of the coinvariants. Consequently the same
argument identifies the product-group Ext with the dual homology of
$\bigotimes_{p\in S}C_{p,\bullet}$.

Lemma~\ref{v4-arith:ic:w10-d:adelic:lem:local} and~\eqref{v4-arith:ic:w10-d:adelic:eq:coinvariant} imply
$H_0(C_{p,\bullet})=0$. The K\"unneth formula for bounded-below complexes
of vector spaces gives
\[
 H_r\left(\bigotimes_{p\in S}C_{p,\bullet}\right)
 \simeq\bigoplus_{\substack{(r_p)\in\mathbb Z_{\ge0}^{S}\\\sum_pr_p=r}}
       \bigotimes_{p\in S}H_{r_p}(C_{p,\bullet}).
\]
For completeness, one can split each vector-space complex as its homology
with zero differential plus a contractible complex. Tensoring these
splittings proves the formula. For $r<n$, every summand has some $r_p=0$
and is zero. This proves the vanishing assertion.

Each nonzero class $\alpha_p$ is a nonzero functional on
$H_1(C_{p,\bullet})$. Choose a homology element on which it is nonzero.
The product functional is nonzero on the tensor of these elements,
which lies in degree $n$. It represents the exterior Yoneda class.

Lift the projective resolution's augmentation to $E_p$.
Its degree-one boundary then factors through $\mathbb C$,
giving a chain map to $[\mathbb C\to E_p]$. The resulting degree-one cocycle represents $\alpha_p$.
Tensoring these chain maps gives exactly the total extension and the
product cocycle, with the tensor signs already specified.
Thus the exterior class is nonzero.

This proof uses genuinely compact open subgroups even for $GL_2$.
It does not treat compact-mod-center subgroups as compact, and it makes
no assertion about vanishing above degree $n$ in the unrestricted
$GL_2$ category.
\end{proof}

\section{The restricted tensor condition}

Let
\[
 G_f=GL_2(\mathbb A_f)=\prod_p'(G_p,K_p).
\]
Its restricted-product topology has neighborhoods of the identity that
contain $\prod_{p\notin S}K_p$ for some finite $S$. A usual algebraic
restricted tensor representation uses nonzero distinguished vectors
$e_p\in V_p^{K_p}$ for all but finitely many $p$. Its vectors are finite
linear combinations of tensors whose factors equal $e_p$ almost
everywhere. For $g=(g_p)\in G_f$, both $g_p\in K_p$ and the distinguished
factor condition hold almost everywhere, making the coordinate action
well-defined.

\begin{proposition}\label{v4-arith:ic:w10-d:adelic:prop:noall}
The local representations $\operatorname{St}_p$ do not admit a nonzero restricted
tensor representation with the standard hyperspecial pointings at almost
every finite prime. More generally, suppose a smooth $G_f$-representation $V$ has
factorizations $V\simeq\operatorname{St}_p\otimes W_p$ for every finite $p$.
If $G_p$ acts only on the first factor, then $V=0$.
\end{proposition}
\begin{proof}
The required pointing space is zero at every prime by
Lemma~\ref{v4-arith:ic:w10-d:adelic:lem:local}. For the second assertion, a vector in a smooth
$G_f$-representation is fixed by an open subgroup. That subgroup contains
$K_p$ for every $p$ outside a finite set. But the stipulated factorization
implies
\[
 V^{K_p}=\operatorname{St}_p^{K_p}\otimes W_p=0.
\]
This equality follows from averaging on the first tensor factor, or from
writing any tensor as a finite sum with linearly independent second
factors. Thus the vector is zero.
\end{proof}

Choosing Iwahori-fixed vectors at infinitely many places does not supply
these hyperspecial pointings. The restricted product relative to those
Iwahori subgroups is a different subgroup and topology. No passage to
an unspecified completed tensor product is part of this construction.

For a finite set $S$, a genuine smooth adelic representation is
\[
 \operatorname{St}_S^{(f)}=\left(\boxtimes_{p\in S}\operatorname{St}_p\right)
               \boxtimes\left(\bigotimes_{p\notin S}'\mathbb C\right),
\]
where each omitted place uses the trivial representation and vector $1$.
It is inflation from $G_S$ under the continuous projection $G_f\to G_S$.
The extension of Proposition~\ref{v4-arith:ic:w10-d:adelic:prop:finite} inflates with it and remains
nonzero: restriction along the inclusion $G_S\hookrightarrow G_f$ is
exact and recovers the original Yoneda class. This representation has
Steinberg factors precisely at the specified finite set. It is not an
all-place Steinberg representation.

\section{An adelic principal-series extension}

The local principal series $E_p$ do have compatible distinguished vectors,
namely the constant functions $\mathbf1_p$. These vectors are fixed by
all of $G_p$. For finite $S$, write $E_S=\boxtimes_{p\in S}E_p$, inflated
to $G_f$ with trivial factors outside $S$. If $S\subset T$, define
\[
 E_S\longrightarrow E_T,\qquad
 x\longmapsto x\otimes\bigotimes_{p\in T\setminus S}\mathbf1_p.
\]

The superscript $\mathrm{res}$ below means restricted tensor product. It does not name an automorphic residual spectrum. These maps are injective and $G_f$-equivariant. The canonical associativity
and symmetry identifications make their composites agree for
$S\subset T\subset U$.

Put $E^{\mathrm{res}}=\operatorname{colim}_S E_S$.
Each vector lies in some $E_S$ and is smooth at the finitely
many primes of $S$. The full groups outside $S$ fix it.
Thus the colimit representation is smooth. Its central character is trivial.

\begin{proposition}\label{v4-arith:ic:w10-d:adelic:prop:global}
Let $X=\prod_p\mathbb P^1(\mathbb Q_p)$ with its compact product topology.
There is a $G_f$-equivariant isomorphism
\[
 E^{\mathrm{res}}\simeq C^{\mathrm{lc}}(X,\mathbb C)
\]
for the coordinatewise right action. The exact sequence
\begin{equation}\label{v4-arith:ic:w10-d:adelic:eq:global}
 0\longrightarrow\mathbb C\longrightarrow E^{\mathrm{res}}
       \longrightarrow Q_{\mathrm{ad}}\longrightarrow0,
 \qquad Q_{\mathrm{ad}}=E^{\mathrm{res}}/\mathbb C,
\end{equation}
is nonsplit in $\operatorname{Rep}^{\mathrm{sm}}_{\mathbb C}(G_f)$ and in its trivial-central-character category.
Both $E^{\mathrm{res}}$ and $Q_{\mathrm{ad}}$ are admissible: their invariants under
every compact open subgroup are finite-dimensional.
\end{proposition}
\begin{proof}
For finite $S$, products of functions identify $E_S$ with
$C^{\mathrm{lc}}(\prod_{p\in S}X_p,\mathbb C)$. Surjectivity follows by covering
clopen subsets of the compact finite product by finitely many clopen
rectangles and refining this cover to a finite partition. Evaluations separate every finite-dimensional space of functions.
For each space involved, choose a basis and finitely many
evaluation points giving an injective matrix.
The tensor of these matrices is injective, proving injectivity.

Every locally constant function on $X$ has a finite cover by basic open
cylinders on which it is constant. Such a finite cover exists by
compactness. The union of the finite coordinate sets defining the
cylinders is finite. Two points equal in these coordinates belong to
the same cylinders and therefore have equal function values.

The function
thus depends on only finitely many coordinates. This proves surjectivity
from the colimit, and injectivity follows from the finite-stage statement.
The insertion maps become pullback along coordinate projections.
Each $g\in G_f$ acts by the coordinate homeomorphisms $x_p\mapsto x_pg_p$,
so the isomorphism is equivariant. Smoothness also follows from the
finite-coordinate description established above.

A compact open subgroup contains $U_S\times\prod_{p\notin S}K_p$
for some finite $S$. Here $U_S$ is a product of local compact
open subgroups at the places of $S$. Invariance under
the tail $K_p$ forces a finite-coordinate function to be independent of
those tail coordinates, because $K_p$ is transitive on $X_p$.
Its invariants therefore lie in $E_S^{U_S}$.

Each local compact open
subgroup has open orbits on the compact space $X_p$, hence finitely many
orbits. It follows that $E_S^{U_S}$ is finite-dimensional.
This proves admissibility of $E^{\mathrm{res}}$. Compact-open invariants
are exact by the same finite averaging argument, so its quotient $Q_{\mathrm{ad}}$
is also admissible.

A splitting of~\eqref{v4-arith:ic:w10-d:adelic:eq:global} would give a $G_f$-invariant retraction
$E^{\mathrm{res}}\to\mathbb C$. Its restriction to the embedded $E_{\{p\}}$
would be a $G_p$-invariant retraction of the constants, contradicting
Lemma~\ref{v4-arith:ic:w10-d:adelic:lem:local}. The same argument applies when the full center
acts trivially.
\end{proof}

The quotient $Q_{\mathrm{ad}}$ retains mixed finite-place terms. To describe them, set
$F_{-1}E^{\mathrm{res}}=0$ and, for $m\ge0$, let $F_mE^{\mathrm{res}}$
be the sum of the images of $E_T$ over finite sets $T$ with $|T|\le m$.
This is an increasing, exhaustive filtration by $G_f$-subrepresentations,
with $F_0E^{\mathrm{res}}=\mathbb C$.

\begin{proposition}\label{v4-arith:ic:w10-d:adelic:prop:filtration}
For every $m\ge0$ there is a canonical isomorphism
\[
 \operatorname{gr}_m^F E^{\mathrm{res}}
 \simeq\bigoplus_{\substack{T\text{ finite}\\|T|=m}}\operatorname{St}_T^{(f)},
\]
where the empty-set factor is $\mathbb C$. The induced exhaustive filtration
$Q_{\mathrm{ad},m}=F_mE^{\mathrm{res}}/\mathbb C$ of $Q_{\mathrm{ad}}$ starts with $Q_{\mathrm{ad},0}=0$ and has these
same associated graded pieces for $m\ge1$.
\end{proposition}
\begin{proof}
At a finite stage $S$, use the two-step filtration
$\mathbb C\mathbf1_p\subset E_p$ in every tensor factor. Take a class in $\boxtimes_{p\in T}\operatorname{St}_p$,
where $T\subset S$ has size $m$. Lift its factors to $E_p$
and insert constants elsewhere. This defines a class in the
$m$th associated graded.

Changing a lift adds a tensor from the preceding
filtration term, so this map is canonical and equivariant.
It is an isomorphism: choose arbitrary vector-space complements to the
constants in each $E_p$ and expand the finite tensor product. The
summands with precisely $m$ complement factors give the asserted graded
piece. The complements prove bijectivity but do not define the map.

The canonical maps are compatible with the constant insertions as $S$
grows. Taking the filtered colimit gives the displayed direct sum, since
each vector involves only finitely many places. Quotienting by the
constant subrepresentation gives the statement for $Q_{\mathrm{ad}}$.
\end{proof}

\Needspace{6\baselineskip}
In particular there is a nonsplit degree-one extension
\[
 0\longrightarrow\mathbb C\longrightarrow F_1E^{\mathrm{res}}
   \longrightarrow\bigoplus_p\operatorname{St}_{\{p\}}^{(f)}\longrightarrow0.
\]
Pulling it back to any single summand gives~\eqref{v4-arith:ic:w10-d:adelic:eq:local}, inflated
to $G_f$. This degree-one construction has a direct sum as its quotient.
The full quotient $Q_{\mathrm{ad}}$ additionally contains every finite mixed tensor
through its filtration. Neither quotient is the all-place tensor
$\bigotimes_p'\operatorname{St}_p$.

Indeed, the local quotient maps $E_p\to\operatorname{St}_p$ send the distinguished
constant vectors to zero. Using those zero vectors as insertions makes every finite
tensor map zero after a new prime is added.
The resulting colimit is zero.
Thus the quotient maps cannot turn~\eqref{v4-arith:ic:w10-d:adelic:eq:global} into a nonzero
all-place Steinberg extension.

\section{Archimedean and automorphic scope}

All preceding Ext groups belong to explicitly stated categories of
smooth finite-adelic or finite-product representations. Tensoring with
the trivial representation of $GL_2(\mathbb R)$ gives a specified
archimedean factor and an action of $GL_2(\mathbb A)$, without adding another
extension class. A nontrivial archimedean extension requires its own
smooth topological representation category or Harish--Chandra module
category, its central character, and its tensor conventions. An
additional degree-one exterior factor contributes an additional degree
when its exterior product is defined. It does not erase the finite-place
degrees above.

A representation of $GL_2(\mathbb A)$ is not thereby a realization on
$GL_2(\mathbb Q)\backslash GL_2(\mathbb A)$. Such a realization requires a specified
automorphic function space or enhanced category, central quotient and
support conditions, and a comparison functor. The symbols
$\operatorname{Ext}^1_{[GL_2]}$ alone do not specify that category. The constructions
here give genuine local and adelic group representations and their
Yoneda extensions. They neither identify a residue with an extension
class nor supply an arithmetic or automorphic comparison.



\section{Galois topology and the geometric recipient}
\label{v4-arith:ic:topology}

For a finite field \(\mathbb F_q\), its absolute Galois group is
procyclic, generated topologically by arithmetic Frobenius.
For a nonarchimedean local field \(F\), the corresponding exact
sequence is
\[
 1\longrightarrow I_F\longrightarrow G_F
   \longrightarrow G_{\mathbb F_q}\longrightarrow1.
\]
Its kernel is inertia. The tame quotient of inertia has its
prime-to-residue-characteristic Tate twist, whereas wild inertia
is a pro-\(p\) group. Neither is this Frobenius quotient.

Precomposition with a quotient \(G\to Q\) embeds representations
of \(Q\) among those representations of \(G\) on which the
kernel acts trivially. It does not descend an arbitrary
representation. Similarly, a representation of a product of
groups has commuting component images, with continuity for
the chosen topologies. Independently chosen local representations
need not define such a product representation.

Unramified twists make the distinction visible without wild
ramification. Two different unramified characters of \(W_F\)
have the same trivial inertia restriction and different
Frobenius eigenvalues. Thus inertia alone does not determine
the full parameter. On a general test scheme, a mapping stack
also retains torsors and their descent. A conjugacy set on
field-valued points does not define that stack.

\subsection{Finite cyclic and ordinary loops}

Let \(k\) be algebraically closed of characteristic zero.
For the finite cyclic group \(C_n\),
\[
 \operatorname{Map}(BC_n,B\mathbb G_m)
 \simeq\bigsqcup_{\zeta\in\mu_n}B\mathbb G_m.
\]
Indeed, a character of \(C_n\) is determined by an \(n\)th root
of unity. Its automorphism group is \(\mathbb G_m\).
This is distinct from ordinary derived loops
\[
 \mathcal L(B\mathbb G_m)
   =\operatorname{Map}(B\mathbb Z,B\mathbb G_m)
   \simeq\mathbb G_m\times B\mathbb G_m.
\]
It is also distinct from an infinitesimal neighborhood of the
identity diagonal.

If \(n\mid m\), the quotient \(C_m\to C_n\) induces
\[
 \operatorname{Map}(BC_n,X)
       \longrightarrow\operatorname{Map}(BC_m,X)
\]
by precomposition. The arrows reverse those of the groups.
A continuous homomorphism from \(\widehat{\mathbb Z}\) to a
discrete group has finite image, since a compact discrete
subgroup is finite. A topology on the coefficient group can
change that condition. The ambient prestack and topology
must therefore be fixed before a profinite loop comparison.

\subsection{Selmer deformation functors}

Fix one finite residual field \(k\), one coefficient ring \(W(k)\),
a residual representation, determinant convention and a finite
ramification set. Let \(\mathcal F_{\mathrm{glob}}\) and
\(\mathcal F_v\) be the continuous deformation functors on the
same category of simplicial Artin coefficient rings.
Give every local condition an actual map
\(\mathcal L_v\to\mathcal F_v\). Then the Selmer functor is
\[
 \mathcal F_{\mathrm{Sel}}
 =\mathcal F_{\mathrm{glob}}
       \times^h_{\prod_v\mathcal F_v}\prod_v\mathcal L_v.
\]

Its construction and tangent fiber are proved in
Theorem~\ref{v4-arith:coefficients:thm:selmer-tangent}.
A cochain quotient notation is available when compatible
subcomplexes have actually been chosen.
Local stacks over unrelated residue-characteristic bases cannot
be inserted into this product without coefficient comparisons.

Framing retains cocycles. A determinant identification gives
the fiber of the trace map on the full shifted cochain complexes.
Fixed determinant value modulo arbitrary strict \(GL_2\)
conjugacy has the different gauge convention described above.
Automorphisms and obstructions remain in their respective
degrees. At coefficient characteristic two, real decomposition
groups also require their specified archimedean or modified
Tate condition.

Formal pullback notation proves neither representability nor
a cotangent-amplitude bound. For example, the dual numbers
over a characteristic-zero field have a two-term cotangent
complex and a nonzero obstruction tangent at their closed point.
Such an obstruction is allowed in a quasi-smooth derived scheme.
Likewise, smoothness of \(X\) does not make its derived loop
space classical: for \(X=\mathbb A^1_k\),
\[
 \mathcal O(\mathcal LX)
   \simeq k[t]\otimes\Lambda(dt),
 \qquad |dt|=-1.
\]

Cotangent comparisons on finite products use the pullbacks of
each factor's cotangent complex over the common base.
Changes of charts use the transitivity triangle.
An inverse-limit formula or an ind-coherent refinement requires
its own continuity theorem.


\section{Restriction to two places}
\label{v4-arith:ic:two-places}

Unramified matrices describe a local deformation problem.
Global restriction imposes further relations between places and
can have several lifts over the same local pair. The following
example retains absolute irreducibility of the residual
representation and fixes its determinant.


Fix distinct finite primes $p,q$ and a coefficient prime
$p_0\ne p,q$. Let $k$ be a finite field of characteristic $p_0$,
and put $W=W(k)$. Choose a finite set $S$ of places of $\mathbb Q$
containing $p,q,p_0,\infty$. Write $G_{\mathbb Q,S}$ for the Galois
group of the maximal extension unramified outside $S$.

Fix a continuous representation
$\bar\rho:G_{\mathbb Q,S}\to GL_2(k)$ which is absolutely
irreducible, and let $\bar D$ be its determinant polynomial law.
Here a determinant of dimension two is a multiplicative homogeneous
polynomial law of degree two on the group algebra, as in
Definition~\ref{v4-arith:ic:w5-h:def:chenevier-det-trace}.
Its characteristic coefficients, rather than a trace alone, define
the object in all residual characteristics.

Choose decomposition-group maps
$G_{\mathbb Q_p},G_{\mathbb Q_q}\to G_{\mathbb Q,S}$.
Choose closed local and global deformation conditions preserved by
restriction. These may include a fixed determinant character lifting
$\det\bar\rho$. Let $A$ be a finite local $W$-algebra with residue field $k$.
Write $\mathcal D_{\mathrm{glob}}(A)$ and $\mathcal D_v(A)$
for the corresponding continuous determinant lifts, where $v=p,q$.
The local residual restrictions need not be absolutely irreducible.

Chenevier's representability theorem gives complete Noetherian local
rings $R^{\mathrm{Ch}},R_p^{\mathrm{Ch}},R_q^{\mathrm{Ch}}$ for
these functors. The precise hypotheses are those of
Proposition~\ref{v4-arith:ic:w5-h:prop:chenevier-moduli-via-det-trace}.
Its source is Chenevier, Proposition~3.3 and Example~3.4, p.~42,
and Example~3.6 and Proposition~3.7, p.~43.
Each closed condition must be represented by the quotient used here.
These are formal deformation rings over one coefficient ring $W$.
An identification with functions on a local or global spectral stack
is additional geometric data.

\subsection{A pair of indistinguishable local restrictions}

\begin{proposition}[nonuniqueness with absolutely irreducible residual representation]
\label{v4-arith:ic:bcglue-core:prop:two-prime-nonuniqueness}
There are distinct continuous rank-two determinant deformations over
$\mathbb Z/8\mathbb Z$ with the same absolutely irreducible reduction
and fixed determinant character. Their restrictions at $3$ and $5$ coincide.
\end{proposition}

\begin{proof}
Let $K$ be the splitting field of $X^3-X-1$ over $\mathbb Q$.
The polynomial has no root modulo $2$, hence is irreducible over
$\mathbb Q$. Its discriminant is $-23$, so its Galois group is
$S_3$. An integral standard representation of this group is generated
by
\[
 A=\begin{pmatrix}0&-1\\1&-1\end{pmatrix},\qquad
 B=\begin{pmatrix}0&1\\1&0\end{pmatrix}.
\]

Direct multiplication gives $A^3=B^2=1$ and $BAB=A^{-1}$.
Modulo $2$, the two eigenlines of $A$ over $\mathbb F_4$ are
interchanged by $B$. Thus this residual representation is absolutely
irreducible. Inflate the integral representation to
$\rho:G_{\mathbb Q}\to GL_2(\mathbb Z_2)$.

Let $\chi$ be the quadratic character of
$L=\mathbb Q(\sqrt{241})$.
For $v=3,5$, the polynomial $X^2-241$ has the simple root $1$
modulo $v$. A root $u_r$ modulo $v^r$ has a unique lift modulo $v^{r+1}$
in its residue class. The correction $u_r+c v^r$ solves an
equation with invertible coefficient $2u_r$ modulo $v$.
Taking the limit gives a root in $\mathbb Z_v$.

Hence $\chi|_{G_{\mathbb Q_v}}=1$ for both places.
Set $\rho'=\chi\rho$. Then
\[
 \bar\rho'=\bar\rho,\qquad
 \det\rho'=\chi^2\det\rho=\det\rho.
\]
The restrictions of $\rho$ and $\rho'$ at $3$ and $5$ agree.
Both are unramified outside
$S=\{2,3,5,23,241,\infty\}$.

The unique quadratic subfield of $K$ is $\mathbb Q(\sqrt{-23})$,
which differs from $L$. Consequently $K\cap L=\mathbb Q$, and
the composite extension has Galois group $S_3\times C_2$.
Choose $g$ whose image is $(1,-1)$. Then
\[
 \det(T-\rho(g))=T^2-2T+1,\qquad
 \det(T-\rho'(g))=T^2+2T+1.
\]

The traces are $2$ and $6$ modulo $8$.
Thus the two determinant laws modulo $8$ lift the same residual
law and are distinct. Their local restrictions and fixed determinant
character agree, as asserted.
\end{proof}

The example lies in a single finite-ramification deformation problem.
A more restrictive arithmetic condition could exclude one lift.
Uniqueness on such a condition would require its own proof.

\section{Restriction, closed image and fibers}
\label{v4-arith:ic:bcglue-core:prismatic}

The local product functor is represented by
\[
 B_{p,q}=R_p^{\mathrm{Ch}}\widehat\otimes_W R_q^{\mathrm{Ch}}
 =\varprojlim_r
 (R_p^{\mathrm{Ch}}/\mathfrak m_p^r)
       \otimes_W(R_q^{\mathrm{Ch}}/\mathfrak m_q^r).
\]
The completion is for the sum of the maximal ideals.
Restriction is a natural transformation
$\mathcal D_{\mathrm{glob}}\to\mathcal D_p\times\mathcal D_q$,
so its map on representing rings has the opposite direction.

\begin{proposition}[the formal restriction diagram]
\label{v4-arith:ic:bcglue-core:thm:prismatic-nygaard-transition}

Restriction induces a continuous local $W$-algebra homomorphism
\begin{equation}
\label{v4-arith:ic:bcglue-core:eq:chenevier-restriction-map}
 r_{p,q}^{\sharp}:B_{p,q}\longrightarrow R^{\mathrm{Ch}},
 \qquad x\otimes y\longmapsto r_p^{\sharp}(x)r_q^{\sharp}(y).
\end{equation}
For the map in Equation~\eqref{v4-arith:ic:bcglue-core:eq:chenevier-restriction-map},
put $J_{p,q}=\ker r_{p,q}^{\sharp}$ and
$C_{p,q}=B_{p,q}/J_{p,q}$. Then
\[
 B_{p,q}\twoheadrightarrow C_{p,q}\hookrightarrow R^{\mathrm{Ch}},
 \qquad
 \operatorname{Spf}R^{\mathrm{Ch}}\longrightarrow
 \operatorname{Spf}C_{p,q}\lhook\joinrel\longrightarrow
 \operatorname{Spf}B_{p,q}.
\]
For a continuous local map $b:B_{p,q}\to A$, the fiber over the
prescribed local pair is represented over $A$ by
$R^{\mathrm{Ch}}\widehat\otimes_{B_{p,q},b}A$.
\end{proposition}

\begin{proof}
Restricting the universal determinant gives maps
$r_v^{\sharp}:R_v^{\mathrm{Ch}}\to R^{\mathrm{Ch}}$.
Their images commute, and their $W$-actions agree.
The completed tensor-product property gives
$r_{p,q}^{\sharp}$ with the displayed formula.
The kernel is closed. The quotient is complete and injects into
$R^{\mathrm{Ch}}$.

Since the rings are profinite, its image is
closed. The formal-spectrum diagram reverses these maps.
Let $A'$ be a finite local $A$-algebra.
A map from the completed base change to $A'$ is a map
$R^{\mathrm{Ch}}\to A'$ agreeing with the prescribed $B_{p,q}$-action through $A$.
This proves the fiber assertion.
\end{proof}

\begin{definition}[the closed image of two-prime restriction]
\label{v4-arith:ic:bcglue-core:def:cross-prime-chenevier}
The formal scheme $\operatorname{Spf}C_{p,q}$ is the closed image
of restriction inside the local deformation product.
A proposed two-prime Selmer algebra $R_{p,q}^{\mathrm{Ch}}$
requires a specified $B_{p,q}$-algebra structure and comparison
$R_{p,q}^{\mathrm{Ch}}\to C_{p,q}$ before the two can be identified.
The quotient $C_{p,q}$ itself imposes the kernel relations only.
\end{definition}

\begin{example}[a flat closed-image map with a nonliftable point]
\label{v4-arith:ic:bcglue-core:ex:closed-image}
Let $B=k[[x]]$ and $R=k[[t]]$, with $x\mapsto t^2$.
Every series in $t$ is uniquely $f(t^2)+t g(t^2)$.
Thus $R$ is free of rank two over $B$, and the map is injective.
Its kernel quotient is $B$ itself.

Nevertheless, the local point
$B\to k[\epsilon]/(\epsilon^2)$, $x\mapsto\epsilon$, has no lift
to $R$. A local image of $t$ is $a\epsilon$, whose square is zero.
Flatness does not turn the closed image into the image functor
on these Artinian points.
\end{example}

\begin{example}[an ordinary fiber with nontrivial derived structure]
\label{v4-arith:ic:bcglue-core:ex:derived-fibre}
For $B=k[[x]]\to R=k$, $x\mapsto0$, the fiber over the residue
point has ordinary coordinate ring $k$.
The free resolution $[B\xrightarrow{x}B]$, in cohomological
degrees $-1,0$, gives
\[
 k\otimes_B^{\mathbb L}k\simeq k\oplus k[1],\qquad
 H^{-1}(k\otimes_B^{\mathbb L}k)=k.
\]
Thus the derived fiber retains information lost by the ordinary
fiber. The same calculation proves that this map is not flat.
\end{example}

\begin{remark}[the information retained by restriction]
\label{v4-arith:ic:bcglue-core:rem:why-chenevier}
The kernel quotient records equations satisfied by local pairs.
The fibers record their global lifts.
Proposition~\ref{v4-arith:ic:bcglue-core:prop:two-prime-nonuniqueness}
shows that those lifts need not be unique, so
$r_{p,q}^{\sharp}$ need not be surjective.
In a derived deformation theory the fiber uses derived completed
tensor product. None of these constructions identifies a Nygaard
filtration or supplies a comparison with prismatic sheaves.
\end{remark}

\section{The pairwise pullback and its central actions}
\label{v4-arith:ic:bcglue-core:local}

\begin{definition}[completed external product and pullback]
\label{v4-arith:ic:bcglue-core:def:diagonal-pullback}
Let $M_p,M_q$ be finite modules over the local deformation rings.
Their completed external product is the finite $B_{p,q}$-module
\[
 M_{p,q}=\varprojlim_r
 (M_p/\mathfrak m_p^rM_p)\otimes_W
 (M_q/\mathfrak m_q^rM_q).
\]
The module obtained by ordinary restriction pullback is
\begin{equation}
\label{v4-arith:ic:bcglue-core:eq:diagonal-pullback}
 M_{p,q}^{\mathrm{res}}
 =R^{\mathrm{Ch}}\widehat\otimes_{B_{p,q}}M_{p,q}
 =\varprojlim_r
 \frac{R^{\mathrm{Ch}}\otimes_{B_{p,q}}M_{p,q}}
 {\mathfrak m_{\mathrm{glob}}^r
       (R^{\mathrm{Ch}}\otimes_{B_{p,q}}M_{p,q})}.
\end{equation}
The algebraic base change is finite over $R^{\mathrm{Ch}}$ and
therefore already complete. Equation~\eqref{v4-arith:ic:bcglue-core:eq:diagonal-pullback}
also specifies its topology.
\end{definition}

\begin{proposition}[the action on the formal pullback]
\label{v4-arith:ic:bcglue-core:def:bernstein-product-action}
\label{v4-arith:ic:bcglue-core:prop:formal-pullback-action}
Let $Z_p,Z_q$ be unital $W$-algebras acting
$R_p^{\mathrm{Ch}}$-linearly and $R_q^{\mathrm{Ch}}$-linearly on
$M_p,M_q$, respectively. Each acting endomorphism is continuous
for the maximal-ideal topology. Then $M_{p,q}^{\mathrm{res}}$ has
a canonical action of $Z_p\otimes_W Z_q$, given by
\[
 (z_p\otimes z_q)(a\otimes(m_p\otimes m_q))
   =a\otimes(z_pm_p\otimes z_qm_q).
\]
This action commutes with the $R^{\mathrm{Ch}}$-action.
\end{proposition}

\begin{proof}
Finite presentations of $M_p,M_q$ show that $M_{p,q}$ is finite
over $B_{p,q}$. Base change gives a finite module over
$R^{\mathrm{Ch}}$. Local linearity makes the displayed action
compatible with the balancing relations and with every quotient
in the inverse limits. It therefore defines an action on the
completed tensors.

The two local actions affect separate factors,
so they commute. Multiplication on the first factor commutes with
both. No topology on the abstract algebras $Z_p,Z_q$ is needed.
\end{proof}

The associated coherent sheaf on $\operatorname{Spf}R^{\mathrm{Ch}}$
is the ordinary pullback of the completed external product.
The derived construction replaces the external product and pullback
by their derived completed versions. Its boundedness requires
additional Tor conditions. The ordinary pullback is exact only on
a domain where the required higher Tor groups vanish.

\begin{proposition}[transport of ordinary central actions]
\label{v4-arith:ic:bcglue-core:thm:local-bernstein}
For a $W$-linear category $\mathcal C$, let
$Z(\mathcal C)=\operatorname{End}(\mathrm{id}_{\mathcal C})$
be its ordinary center. A $W$-linear equivalence
$F:\mathcal C\to\mathcal D$ induces an isomorphism
\begin{equation}
\label{v4-arith:ic:bcglue-core:eq:local-bernstein-equals-spectral}
 Z(\mathcal C)\xrightarrow{\sim}Z(\mathcal D).
\end{equation}
For a fully faithful functor, this statement applies to its essential
image. It does not identify the center of a larger target category.
\end{proposition}

\begin{proof}
Choose a quasi-inverse $G$ and the unit and counit isomorphisms.
A natural endomorphism $z$ of $\mathrm{id}_{\mathcal C}$ defines
an endomorphism of $\mathrm{id}_{\mathcal D}$ by conjugating
$FzG$ with the counit. Naturality gives a central action, and
composition and addition are preserved. This defines the map in
Equation~\eqref{v4-arith:ic:bcglue-core:eq:local-bernstein-equals-spectral}.
The construction with $G$
is its inverse. Restricting a fully faithful functor to its essential
image makes it an equivalence.
\end{proof}

\begin{remark}[local coefficients and Bernstein blocks]
\label{v4-arith:ic:bcglue-core:rem:local-domain}
For smooth representations of $GL_2(\mathbb Q_v)$ in
characteristic zero with coefficient prime different from $v$,
the Bernstein decomposition has the form
\begin{equation}
\label{v4-arith:ic:bcglue-core:eq:bernstein-decomp}
 \mathrm{Rep}(GL_2(\mathbb Q_v))
   =\prod_{\mathfrak s}\mathrm{Rep}_{\mathfrak s}(GL_2(\mathbb Q_v)).
\end{equation}
For the blocks in Equation~\eqref{v4-arith:ic:bcglue-core:eq:bernstein-decomp},
the ordinary scalar comparison and its precise domain are proved in
Proposition~\ref{v4-arith:w10-e:prop:excursion-bernstein}.
Its image in the spectral categorical center is the scalar subalgebra
of Proposition~\ref{v4-arith:w10-e:prop:scalar-centre-obstruction}.
Using a Bernstein center as $Z_v$ in the formal pullback requires
an action over the chosen coefficient ring and on the actual module
$M_v$. A categorical comparison additionally needs a functor with
specified essential image. A local $(\varphi,\Gamma)$-module at
residue prime $v$ is not identified here with a deformation over
$W(k)$ of coefficient characteristic $p_0\ne v$.
\end{remark}

\section{Equivariant isomorphisms on formal charts}
\label{v4-arith:ic:bcglue-core:core-theorem}

The module $M_{p,q}^{\mathrm{res}}$ is constructed from local data.
Let $N$ be an independently specified global module.
The pairwise comparison asks for an isomorphism from the former
to the latter which intertwines their central actions.
Equal determinant laws alone do not solve this problem.

\begin{example}[extension data invisible to determinant laws]
\label{v4-arith:ic:bcglue-core:ex:extension-data}
Let $R=k[[t]]$ and $M=N=R^2$. Give them actions of $R[z]$ by
\[
 z|_M=0,\qquad z|_N=J=\begin{pmatrix}0&t\\0&0\end{pmatrix}.
\]
For every polynomial $f$, one has
$f(J)=f(0)1+f'(0)J$, so both actions have characteristic
polynomial $(X-f(0))^2$ for $f(z)$.
The determinant polynomial laws agree after every coefficient
extension.

The actions also coincide modulo $t$.
An invertible intertwiner $U$ would satisfy $JU=0$, contradicting
$J\ne0$. These are reducible algebra representations.
They exhibit the loss of extension data for coherent objects.
They do not contradict reconstruction of a representation from an
absolutely irreducible residual determinant.
\end{example}

\begin{proposition}[the scheme of equivariant isomorphisms]
\label{v4-arith:ic:bcglue-core:thm:pairwise-reduces-to-chenevier}
Let $(R,\mathfrak m,k)$ be a complete Noetherian local ring.
Let $M,N$ be free $R$-modules of rank $n$ with actions of an
$R$-algebra generated by $a_1,\ldots,a_s$.
Choose bases and let $P_i,Q_i\in\operatorname{Mat}_n(R)$ denote
the action matrices on $M,N$.
The functor of equivariant isomorphisms is represented by
\begin{equation}
\label{v4-arith:ic:bcglue-core:eq:pairwise-chenevier-equivalent}
 \operatorname{Isom}(M,N)
 =\operatorname{Spec}
 \frac{R[u_{ab},d]}{
 (d\det(U)-1,\ (Q_iU-UP_i)_{ab}:i,a,b)}.
\end{equation}
Here $U=(u_{ab})$, with $1\le a,b\le n$ and $1\le i\le s$.
\end{proposition}

\begin{proof}
Over any commutative $R$-algebra $A$, a matrix $U$ is invertible
exactly when its determinant is a unit. The variable $d$ specifies
the unique inverse of that determinant. Intertwining the generators
is equivalent to intertwining the whole algebra. Thus maps from
the coordinate ring in
Equation~\eqref{v4-arith:ic:bcglue-core:eq:pairwise-chenevier-equivalent}
to $A$ are exactly the equivariant
isomorphisms $M\otimes_R A\to N\otimes_R A$, naturally in $A$.
\end{proof}

\begin{proposition}[successive lifting obstructions]
\label{v4-arith:ic:bcglue-core:cor:part-a-domain-reduction}
In the preceding situation, put
$I_r=\mathfrak m^r/\mathfrak m^{r+1}$ for $r\ge1$.
Let $U_r$ be an invertible intertwiner modulo $\mathfrak m^r$.
Choose a lift $\widetilde U$ modulo $\mathfrak m^{r+1}$ and set
\[
 E_i=Q_i\widetilde U-\widetilde U P_i
       \in\operatorname{Mat}_n(I_r).
\]
Write bars for reduction modulo $\mathfrak m$, and define
\[
 L_r:\operatorname{Mat}_n(I_r)\longrightarrow
       \operatorname{Mat}_n(I_r)^s,
 \qquad V\longmapsto(\overline Q_iV-V\overline P_i)_i.
\]

The class $[(E_i)_i]\in\operatorname{coker}L_r$ is independent
of the lift. It vanishes precisely when this $U_r$ extends to
an invertible intertwiner modulo $\mathfrak m^{r+1}$.
A compatible sequence of such lifts determines an equivariant
isomorphism over $R$.
\end{proposition}

\begin{proof}
Changing $\widetilde U$ by $V\in\operatorname{Mat}_n(I_r)$
changes the defects by $L_r(V)$, since $\mathfrak m I_r=0$.
Thus the class is independent of the lift and vanishes exactly
when a correction solves every intertwining equation.
The determinant of a corrected lift remains invertible modulo the
maximal ideal, hence is a unit. Completeness gives the matrix limit
of a compatible sequence. The intertwining equations pass to the
limit, and its determinant is a unit because its residue is.
Conversely an equivariant isomorphism over $R$ gives such a
compatible sequence by reduction.
\end{proof}

The obstruction depends on the specified partial lift $U_r$.
Vanishing for a compatible sequence, rather than separate choices
at unrelated orders, constructs the limit.
For general coherent complexes the comparison uses a derived
equivariant Hom complex. The finite free criterion does not assert
that the intended arithmetic sheaves satisfy its hypotheses.

\section{Global actions and geometric comparison}
\label{v4-arith:ic:bcglue-core:global}

\begin{proposition}[the action generated by local centers]
\label{v4-arith:ic:bcglue-core:thm:colimit-to-global-reduces-to-AG}
Let $E$ be a commutative ring and $Z_v$ commutative unital
$E$-algebras. For finite sets $S$ of places put
$A_S=\bigotimes_{v\in S,E}Z_v$, with transition maps inserting
units. Set
\begin{equation}
\label{v4-arith:ic:bcglue-core:eq:global-bernstein-adelic}
 A=\varinjlim_{S\text{ finite}} A_S.
\end{equation}
Pairwise commuting unital actions
$\alpha_v:Z_v\to\operatorname{End}_E(M)$ on an $E$-module
induce a unique action
\[
 \alpha:A\to\operatorname{End}_E(M),\qquad
 \alpha(\otimes_{v\in S}z_v)=\prod_{v\in S}\alpha_v(z_v).
\]
Its image is $A/\ker\alpha$. Natural actions on an $E$-linear
category instead give a map to its ordinary center.
\end{proposition}

\begin{proof}
The product of the commuting actions is multiplicative and
$E$-multilinear on each finite product. The tensor-product property
gives an action of $A_S$. Inserting a unit changes no endomorphism,
so these maps are compatible. The colimit in
Equation~\eqref{v4-arith:ic:bcglue-core:eq:global-bernstein-adelic}
therefore gives $\alpha$.
The quotient assertion is the first isomorphism theorem.
Products of natural endomorphisms are natural, which proves the
categorical assertion.
\end{proof}

\begin{proposition}[faithfulness and generation for a global scalar center]
\label{v4-arith:ic:bcglue-core:thm:global-bernstein-chenevier}
Suppose the preceding actions define
$\alpha:A\to Z(\mathcal C)$ on an $E$-linear category.
Then
\begin{equation}
\label{v4-arith:ic:bcglue-core:eq:global-bernstein-chenevier}
 A/\ker\alpha\xrightarrow{\sim}\operatorname{im}\alpha
       \subseteq Z(\mathcal C).
\end{equation}
This image is the entire ordinary center exactly when every
central endomorphism is generated by the finite local actions.
The map from $A$ itself is an isomorphism exactly when it is
also faithful. No identification with $R^{\mathrm{Ch}}$ is
implicit in this statement.
\end{proposition}

\begin{proof}
The quotient isomorphism in
Equation~\eqref{v4-arith:ic:bcglue-core:eq:global-bernstein-chenevier}
follows from the preceding proposition.
Surjectivity means precisely that every central endomorphism
belongs to this image. Injectivity means $\ker\alpha=0$.
Together these are exactly the isomorphism conditions.
\end{proof}

\begin{remark}[first-order fixed-determinant deformations]
\label{v4-arith:ic:bcglue-core:rem:chenevier-tangent}
Let $V=k^2$ with action $\bar\rho$, and let
$\operatorname{ad}\bar\rho=\operatorname{End}_k(V)$ with conjugation
action. A lift over $k[\epsilon]/(\epsilon^2)$ has the form
$(1+\epsilon c(g))\bar\rho(g)$.
Multiplicativity is the continuous cocycle equation for $c$.
The fixed-determinant condition is $\operatorname{tr}c(g)=0$.

Strict change of basis changes $c$ by the coboundary of an element
of $\operatorname{ad}\bar\rho$. Thus the tangent space before
additional local conditions is
\[
 \frac{Z^1_{\mathrm{cts}}(G_{\mathbb Q,S},
                  \operatorname{ad}^0\bar\rho)}
 {d\,\operatorname{ad}\bar\rho}.
\]
The denominator consists of traceless cocycles.
When $2$ is invertible, scalar subtraction identifies it with
$d\,\operatorname{ad}^0\bar\rho$, giving the usual
$H^1_{\mathrm{cts}}(G_{\mathbb Q,S},\operatorname{ad}^0\bar\rho)$.

At characteristic two this simplification needs justification.
Additional deformation conditions restrict this tangent problem.
A tangent-space comparison does not identify the complete rings.
\end{remark}


\section{Frobenius coefficients and global reconstruction}
\label{v4-arith:ic:frobenius}

Unramified characteristic polynomials provide scalar functions on a
Galois deformation space. The full collection determines the complete
local ring with its deformation topology. A finite generating subset
exists, but its choice depends on the deformation problem.

Let $K$ be a number field, let $k$ be a finite field of characteristic
$p_0$, and put $W=W(k)$. Fix a finite set $S$ containing the infinite
places and all places above $p_0$. Write $G=G_{K,S}$ for the Galois
group of the maximal extension of $K$ unramified outside $S$.
Fix a continuous absolutely irreducible representation
$\bar\rho:G\to\mathrm{GL}_2(k)$ and a continuous character
$\delta:G\to W^\times$ lifting $\det\bar\rho$.

For a finite local $W$-algebra $A$ with residue field $k$, a deformation
is a continuous lift $\rho:G\to\mathrm{GL}_2(A)$ with determinant
$\delta_A$. Two lifts are strictly equivalent if they are conjugate
by an element of $1+M_2(\mathfrak m_A)$.
Let $R=R^{\delta}_{\bar\rho,S}$ be the complete Noetherian local
$W$-algebra representing these deformation classes. Equivalently,
$R$ represents continuous rank-two determinant laws lifting
$\det\bar\rho$ with determinant character $\delta$.

The representability and comparison hypotheses are those of
Proposition~\ref{v4-arith:ic:w5-h:prop:chenevier-moduli-via-det-trace}.
In particular, Mazur's finiteness condition gives Noetherianity.
Let $\rho^u$ be a universal lift, let $T=\operatorname{tr}\rho^u$,
and let $\mathfrak m$ be the maximal ideal of $R$.
All topologies on $R$ below are $\mathfrak m$-adic.

For a finite place $v\notin S$, use arithmetic Frobenius
$\mathrm{Frob}_v$, acting on residue fields by the $\#k(v)$-power map.
Its conjugacy class in $G$ is well defined. Put
\[
 t_v=T(\mathrm{Frob}_v),\qquad
 d_v=\delta(\mathrm{Frob}_v),\qquad
 P_v(X)=X^2-t_vX+d_v.
\]
These conventions involve no choice of a square root of $\#k(v)$.

The distinction between a representation and its determinant law is
already visible over $\mathbb F_2[\epsilon]/(\epsilon^2)$. For the cyclic group
$C_2=\langle s\rangle$, the representations
\[
 \rho_0(s)=I_2,\qquad \rho_1(s)=(1+\epsilon)I_2
\]
have the same trace and determinant on every element. Their determinant
laws agree, since
\[
 \det\rho_1(a+bs)=(a+b(1+\epsilon))^2=(a+b)^2
 =\det\rho_0(a+bs)
\]
after every commutative coefficient extension. They are not conjugate,
since both images of $s$ are scalar and they differ. Thus equality of
determinant laws need not determine a representation when the residual
representation is reducible.

A framing retains a chosen basis, so framed deformations are lifts
without quotienting by strict conjugacy. Over
$A=\mathbb F_2[\epsilon]/(\epsilon^2)$, consider the constant lift
$\sigma_0:\mathrm{GL}_2(\mathbb F_2)\to\mathrm{GL}_2(A)$ of the
natural representation. Its reduction is absolutely irreducible:
the matrices
\[
 U=\begin{pmatrix}1&1\\0&1\end{pmatrix},\qquad
 V=\begin{pmatrix}1&0\\1&1\end{pmatrix}
\]
have different unique invariant lines over $\overline{\mathbb F}_2$.
Put $P=I_2+\epsilon E_{12}$ and $\sigma_1=P\sigma_0P^{-1}$.

Since $\epsilon^2=0$ and the characteristic is $2$,
\[
 PUP^{-1}=U,\qquad
 PVP^{-1}=V+\epsilon(E_{12}V-VE_{12})=V+\epsilon I_2.
\]
Both lifts have determinant $1$. Conjugation preserves all traces and
determinant laws, but the displayed images of $V$ differ.
These are the same unframed deformation and distinct framed deformations.

\begin{lemma}[trace determination with absolutely irreducible reduction]
\label{v4-arith:ic:w10-e:trace-conjugacy}
Let $A$ be a commutative local ring with residue field $k$.
Suppose $\rho_1,\rho_2:G\to\mathrm{GL}_2(A)$ reduce to the same
absolutely irreducible representation over $k$.
If $\operatorname{tr}\rho_1(g)=\operatorname{tr}\rho_2(g)$ for every
$g\in G$, then $\rho_1$ and $\rho_2$ are strictly equivalent.
The assertion includes residue characteristic $2$.
\end{lemma}

\begin{proof}
The $k$-linear span of the residual group matrices is $M_2(k)$.
Here is a proof in dimension two. Extend scalars to an algebraic
closure $\bar k$ and let $E$ be the algebra spanned by the group image.
Absolute irreducibility says that $E$ acts irreducibly.

It contains a
nonscalar matrix, and subtracting an eigenvalue produces a nonzero
rank-one matrix $u\varphi$. Irreducibility gives $Eu=\bar k^2$ and
$\varphi E=(\bar k^2)^*$ over the algebraic closure. Indeed, the annihilator
of $\varphi E$ is invariant and cannot be the whole space.
Thus the matrices $a(u\varphi)b$ span all endomorphisms.
Descent of dimension gives the asserted span over $k$.

Choose $g_1,\ldots,g_4\in G$ whose residual matrices form a basis.
For $a=1,2$, the matrices $A_{a,i}=\rho_a(g_i)$ form an $A$-basis of
$M_2(A)$, since their coordinate determinant is a unit.
The trace pairing is perfect over every commutative ring:
\[
 \operatorname{tr}(E_{ij}E_{rs})=\delta_{jr}\delta_{is}.
\]
Hence the common Gram matrix
$Q_{ij}=\operatorname{tr}\rho_a(g_i g_j)$ is invertible.

The coordinates $c_i(g)$ of $\rho_a(g)$ in this basis satisfy
\[
 \sum_i c_i(g)Q_{ij}=\operatorname{tr}\rho_a(gg_j).
\]
They are therefore independent of $a$.
The $A$-linear map $\Phi(A_{1,i})=A_{2,i}$ sends
$\rho_1(g)$ to $\rho_2(g)$ for every $g$.
Since products of basis matrices are group matrices, $\Phi$ is a
unital algebra isomorphism. Its reduction is the identity.

Put $F_{ij}=\Phi(E_{ij})$, $v=F_{11}e_1$, and $w=F_{21}v$.
The matrix $P=(v,w)$ is congruent to $I_2$ modulo $\mathfrak m_A$,
so it is invertible. The matrix-unit relations give
$F_{ij}P=PE_{ij}$. Consequently $\Phi(X)=PXP^{-1}$, which proves
strict equivalence. No division by $2$ occurs.
\end{proof}

For example, over $\mathbb F_2$ the matrices
\[
 I=\begin{pmatrix}1&0\\0&1\end{pmatrix},\quad
 U=\begin{pmatrix}1&1\\0&1\end{pmatrix},\quad
 V=\begin{pmatrix}1&0\\1&1\end{pmatrix},\quad UV
\]
form a basis. Their trace Gram matrix is
\[
 \begin{pmatrix}
 0&0&0&1\\0&0&1&0\\0&1&0&0\\1&0&0&1
 \end{pmatrix},
\]
whose determinant is $1$, although $\operatorname{tr}(I)=0$.

\begin{lemma}[unramified coefficients determine continuous determinants]
\label{v4-arith:ic:w10-e:frobenius-law}
Consider two continuous rank-two determinant laws over a complete
separated local coefficient ring. They agree if and only if their
characteristic polynomials agree at every $v\notin S$.
For deformations of $\bar\rho$ with fixed determinant $\delta$,
equality of Frobenius traces suffices. Their representation classes
then agree.
\end{lemma}

\begin{proof}
Chebotarev's theorem implies that the union of the conjugacy classes
of $\mathrm{Frob}_v$, for $v\notin S$, is dense in $G$.
Indeed, every conjugacy class in every finite quotient of $G$ occurs
at infinitely many such places. The inverse image of any nonempty
open subset of a finite quotient therefore meets this union.
Trace and determinant are continuous conjugacy-invariant functions.
Equality on Frobenius classes consequently implies equality on $G$.

The passage from group coefficients to polynomial laws is integral.
For a rank-two determinant law $D$ with trace $T_D$ and determinant
character $d$, its value after any commutative coefficient extension
$A\to C$ is
\[
 D_C\left(\sum_{i=1}^r a_i g_i\right)
 =\sum_i a_i^2d(g_i)
  +\sum_{i<j}a_i a_jd(g_j)T_D(g_i g_j^{-1}).
\]
To derive the mixed coefficient, multiply $Ug_i+Vg_j$ on the right
by $g_j^{-1}$ and use multiplicativity. The characteristic polynomial
of $g_i g_j^{-1}$ gives
\[
 D(Ug_i+Vg_j)
 =d(g_i)U^2+d(g_j)T_D(g_i g_j^{-1})UV+d(g_j)V^2.
\]

A homogeneous quadratic polynomial has only diagonal and mixed
coefficients, so this determines the displayed formula for every
finite sum. Thus equality of $(T_D,d)$ implies equality of the laws.
The previous lemma proves the assertion about representation classes.
\end{proof}

\begin{theorem}[global scalar reconstruction from Frobenius coefficients]
\label{v4-arith:ic:w10-e:thm:bc-bernstein-galois}
Define the polynomial parameter algebra
\[
 H^S=W[\mathsf T_v,\mathsf D_v,\mathsf D_v^{-1}:v\notin S]
\]
and its fixed-determinant quotient
\[
 H^{S,\delta}=H^S/(\mathsf D_v-d_v:v\notin S)
 \simeq W[\mathsf T_v:v\notin S].
\]
The map
\[
 \eta:H^{S,\delta}\longrightarrow R,\qquad
 \mathsf T_v\longmapsto t_v
\]
has $\mathfrak m$-adically dense image.

Choose $a_v\in W$ lifting $\operatorname{tr}\bar\rho(\mathrm{Frob}_v)$,
put $\mathfrak n=(p_0,\mathsf T_v-a_v:v\notin S)$, and define
\[
 \widehat H^{S,\delta}
 =\varprojlim_n H^{S,\delta}/\mathfrak n^n.
\]
The induced continuous homomorphism
$\widehat\eta:\widehat H^{S,\delta}\to R$ is surjective.
It gives the topological quotient
\begin{equation}
\label{v4-arith:ic:w10-e:eq:bc-global}
 R\simeq\widehat H^{S,\delta}/\ker\widehat\eta.
\end{equation}
Moreover, finitely many places $v_1,\ldots,v_d\notin S$ give a
continuous surjection
\[
 W[[X_1,\ldots,X_d]]\longrightarrow R,\qquad
 X_i\longmapsto t_{v_i}-a_{v_i}.
\]
The same density conclusion holds after any quotient imposing a
closed deformation condition.
\end{theorem}

\begin{proof}
Let $B$ be the closure of $W[t_v:v\notin S]$ in $R$.
It is a closed subring. Chebotarev and continuity show that
$T(g)\in B$ for every $g\in G$.
The ring $B$ is profinite, with finite quotients
$B_n=B/(B\cap\mathfrak m^n)$.

It is local with residue field $k$.
Indeed, if $b\in B$ has nonzero residue, choose a unit $u\in W$
with the same residue. Writing $b=u(1-x)$ with
$x\in B\cap\mathfrak m$, the convergent series
$u^{-1}\sum_{j\geq0}x^j$ belongs to $B$ and inverts $b$.

The quadratic formula in the previous proof defines the universal
determinant law over $B$, because all its coefficients belong to $B$.
Its multiplicativity follows as a polynomial identity: the identities
hold after the injection $B[\mathbf X]\hookrightarrow R[\mathbf X]$,
and hence already hold over $B$. This also proves the identities after
arbitrary coefficient extensions, without a flatness assumption.
Its coefficient functions are continuous for the subspace topology
on $B$, and its determinant character is $\delta$.

The representing property gives compatible maps $f_n:R\to B_n$.
Their inverse limit is a continuous map $f:R\to B$.
If $i:B\hookrightarrow R$ is the inclusion, then $i\circ f$
induces the universal determinant over each $R/\mathfrak m^n$.
Uniqueness in the representing property gives
$i\circ f=\operatorname{id}_R$. Therefore $B=R$.
This argument uses finite local quotients of $B$ and does not assume
that $B$ is Noetherian before proving the equality.

Density implies that the classes of $t_v-a_v$ span
$\mathfrak m/(\mathfrak m^2+p_0R)$ over $k$.
In detail, approximate any element of $\mathfrak m$ modulo
$\mathfrak m^2+p_0R$ by a polynomial in the $t_v$. Substituting
$t_v=a_v+(t_v-a_v)$ leaves only its linear part in this quotient.
Noetherianity makes this vector space finite dimensional, so choose
finitely many places whose classes form a basis.

Nakayama's lemma gives
$\mathfrak m=(p_0,t_{v_1}-a_{v_1},\ldots,t_{v_d}-a_{v_d})$.
Completeness then gives the displayed power-series surjection:
approximate an element successively modulo $\mathfrak m^n$, lifting
its degree-$n$ correction to $(p_0,X_1,\ldots,X_d)^n$.
The resulting corrections converge in the power-series ring.

The polynomial map $\eta$ sends $\mathfrak n^n$ into
$\mathfrak m^n$, so it extends to $\widehat\eta$.
The preceding finite power-series map factors through
$\widehat H^{S,\delta}$, proving surjectivity.
If $F_n$ is the kernel of the projection
$\widehat H^{S,\delta}\to H^{S,\delta}/\mathfrak n^n$, then
$\widehat\eta(F_n)=\mathfrak m^n$.
The inclusion from left to right follows by construction.

The reverse
inclusion follows from the finite power-series surjection on the
$n$th powers of the maximal ideals. Thus the quotient topology is
exactly the $\mathfrak m$-adic topology.
Density passes to every quotient of $R$.
\end{proof}

\begin{remark}[topology and global relations]
\label{v4-arith:ic:w10-e:global-relations}
If $A_0=\eta(H^{S,\delta})$, its completion for the induced filtration
$A_0\cap\mathfrak m^n$ is canonically $R$:
\[
 \varprojlim_n A_0/(A_0\cap\mathfrak m^n)\simeq R.
\]
The kernel of $\widehat\eta$ records the global deformation relations.
The theorem does not identify it with the closure of the kernel of
$\eta$. Such an identification requires a separate comparison of
completion and quotient. The finite set of generating places depends
on the deformation problem. An arbitrarily prescribed pair of places
need not generate its ring.
\end{remark}

\begin{proposition}[obstruction to a product of deformation components]
\label{v4-arith:ic:w10-e:local-product-obstruction}
The fixed-residual ring $R$ has no decomposition as a product of two
nonzero unital rings. More generally it cannot be a unital restricted
product containing the characteristic idempotents of two nonzero
factors.
\end{proposition}

\begin{proof}
If $e^2=e$ in a local ring, one of $e$ and $1-e$ is a unit.
The equation $e(1-e)=0$ then gives $e=1$ or $e=0$.
A product of two nonzero rings contains the nontrivial idempotent
$(1,0)$. The same obstruction applies whenever a restricted product
contains these characteristic idempotents.
\end{proof}

The algebra $H^{S,\delta}$ consists of formal characteristic-polynomial
parameters. An automorphic action requires a separate module or
category and a homomorphism from these parameters to its endomorphisms.
An identification of its completed action image with $R$ additionally
requires equality of the two kernels. The scalar reconstruction gives
no such equality. It also supplies no local type decomposition,
flatness assertion, prismatic filtration, untilt descent, or
identification with a higher center.


\begin{example}[a split generic fiber of a local ring]
\label{v4-arith:ic:w6-g:ex:generic-fibre-split}
Put $R=\mathbb Z_2[[x]]/(x(x-2))$.
It is local with residue field $\mathbb F_2$.
The relation $x^m=2^{m-1}x$ for $m\ge1$ reduces every series
uniquely to $a+bx$ with $a,b\in\mathbb Z_2$.
The reductions converge $2$-adically.

Uniqueness follows by evaluating at $x=0$ and $x=2$:
$a=0$ and $2b=0$ imply $a=b=0$.
After inverting $2$, these evaluations give
\begin{equation}
\label{v4-arith:ic:w6-g:eq:bc-global-bernstein-galois}
 R[1/2]\xrightarrow{\sim}\mathbb Q_2\times\mathbb Q_2,
 \qquad a+bx\longmapsto(a,a+2b).
\end{equation}
The inverse in Equation~\eqref{v4-arith:ic:w6-g:eq:bc-global-bernstein-galois}
sends $(u,v)$ to $u+(v-u)x/2$.
Integrally, the image consists of pairs in $\mathbb Z_2^2$ with
$v-u\in2\mathbb Z_2$, so $(1,0)$ is not in the image.
The generic-fiber splitting therefore does not split $R$.
\end{example}

This distinction is necessary before relating a Bernstein component
to a formal deformation space. Such a comparison must specify
whether it concerns the integral ring, its generic fiber, a quotient,
or an action image. A correspondence of closed points establishes
none of those ring identifications by itself.


\begin{remark}[coefficient variants]
\label{v4-arith:ic:w10-e:coefficient-variants}
The same proof applies over a complete discrete valuation ring
$\mathcal O$ finite over $W$. Fix its finite residue field,
residual representation and determinant character.
Replace $p_0$ by a uniformizer of $\mathcal O$, and assume the
corresponding complete Noetherian representing ring exists.
For a deformation problem over $k$ itself, omit the uniformizer
from the residual ideal. These are statements about the indicated
coefficient problems. A comparison between their representing rings
requires its own base-change map and proof.

The density argument uses only the representing determinant functor.
Absolute irreducibility is used to identify that functor with
representation deformations and to obtain strict conjugacy from traces.
\end{remark}

The polynomial-law identities are Chenevier, arXiv:0809.0415v2,
Example~1.8 and Lemma~1.9, pp.~10--11.
Definition--Proposition~2.18 and Theorem~2.22, pp.~32--34,
give the Cayley--Hamilton and absolutely irreducible comparisons.
Proposition~3.3 through Proposition~3.7, pp.~42--43,
supply the deformation-ring assertions. The number-field density theorem is stated in
Stevenhagen--Lenstra, \emph{Chebotarev and his density theorem},
Appendix, p.~35. The finite set of generating places in the theorem
is not a prescribed pair: the two-prime nonuniqueness example of
Proposition~\ref{v4-arith:ic:bcglue-core:prop:two-prime-nonuniqueness}
retains its different quantifiers.

\subsection{The remaining global automorphic comparison}

\begin{proposition}[a specified action quotient and its comparison]
\label{v4-arith:ic:w10-e:thm:bcglue-glue}
Use the coefficient ring and deformation problem of
Theorem~\ref{v4-arith:ic:w10-e:thm:bc-bernstein-galois}.
Suppose an automorphic module or category has a complete separated
commutative $W$-algebra $A_{\mathrm{act}}$ of scalar actions.
Assume there is a continuous surjection
\[
 \theta:\widehat H^{S,\delta}\longrightarrow A_{\mathrm{act}}.
\]
Give $A_{\mathrm{act}}$ the quotient topology for this map.

If
$\ker\theta=\ker\widehat\eta$, there is a unique topological
$W$-algebra isomorphism $A_{\mathrm{act}}\simeq R$ intertwining the
maps from $\widehat H^{S,\delta}$.
To identify this algebra with the full ordinary center of the global
category, additionally require that its central action be faithful and
exhaust that center. An identification of higher centers additionally
requires a specified equivalence of their derived algebras, not only
this degree-zero isomorphism.
\end{proposition}
\begin{proof}
Both algebras are topological quotients of the same algebra by the
same ideal. The identity of the source descends to the claimed
isomorphism, with inverse induced in the same way. It is a homeomorphism
because both topologies are quotient topologies. Surjectivity of the
source map proves uniqueness.

A faithful central action identifies an
algebra with its image. Equality with the full center requires that
image to be the whole center. The argument has made no assertion about
higher cohomology or operations.
\end{proof}

The Frobenius theorem constructs $\widehat\eta$ and its scalar quotient.
It does not construct $\theta$ or compare these kernels. The unframed
Galois deformation problem is fixed at one residual representation.
An automorphic spherical locus, local Bernstein types and their
coefficient changes must be specified independently. Almost-everywhere
sphericity does not provide the missing ring maps.

\section{Coherent gluing and its hypotheses}
\label{v4-arith:ic:descent}

A restriction diagram reconstructs objects only when it retains
their overlap and extension maps. Closed and open restrictions
already show the issue on the affine line. A formal neighborhood
provides a positive coherent gluing theorem.

\subsection{Closed and open restriction on the affine line}

Let $k$ be a field and put $A=k[t]$, $C=k[t,t^{-1}]$.
Let $X=\operatorname{Spec}A$, $Z=V(t)$, and $U=D(t)$.
Write $\operatorname{mod}_{\mathrm{fg}}(R)$ for the category of
finitely generated left modules over a ring $R$.
The intersection $Z\times_X U$ is empty, since its coordinate ring
is $(A/(t))[t^{-1}]=0$.
Consequently a fiber product of categories over the category of this
intersection is a product.

\begin{proposition}
\label{v4-arith:ic:w5-h:closed-open-failure}
The restriction functor
\[
 \operatorname{mod}_{\mathrm{fg}}(A)\longrightarrow
 \operatorname{mod}_{\mathrm{fg}}(k)\times
 \operatorname{mod}_{\mathrm{fg}}(C),\qquad
 M\longmapsto(M/tM,M[t^{-1}]),
\]
is not full. The corresponding derived restriction functor on
perfect complexes is not full either.
\end{proposition}
\begin{proof}
The object $A$ maps to $(k,C)$. Its endomorphism homomorphism is
\[
 A\longrightarrow k\times C,\qquad f\longmapsto(f(0),f).
\]
The endomorphism $(1,0)$ does not belong to its image.
The same calculation applies to degree-zero derived endomorphisms.
\end{proof}

\subsubsection{The extension retained by localization}

Let $D(A)$ be the stable $k$-linear $\infty$-category obtained by
inverting quasi-isomorphisms of unbounded cochain complexes of
$A$-modules. Cohomological shifts satisfy $H^n(M[1])=H^{n+1}(M)$.
Let $D_t(A)$ be its full subcategory of complexes whose cohomology
modules are elementwise killed by powers of $t$.
Extension of scalars $L=C\otimes_A^{\mathbf L}-$ is left adjoint
to restriction of scalars $j_*:D(C)\to D(A)$.

Because $C$ is a localization, it is flat and
$C\otimes_A C=C$. Thus $Lj_*=\operatorname{id}$ and $j_*$ is fully
faithful. Define
\[
 \Gamma_t M=\operatorname{fib}(M\longrightarrow j_*LM).
\]
Localization is exact, so $L\Gamma_tM=0$.

Exactness of localization on modules identifies $\ker L$ with
$D_t(A)$. If $K\in D_t(A)$, adjunction gives
$\operatorname{Map}_{D(A)}(K,j_*N)=0$ for every $N\in D(C)$.
Applying mapping spaces to the defining fiber sequence therefore
shows that $\Gamma_t$ is right adjoint to $D_t(A)\hookrightarrow D(A)$.

Define a stable $\infty$-category $\mathcal G$ whose objects are
\[
 (K,N,\delta),\qquad K\in D_t(A),\quad N\in D(C),\quad
 \delta:j_*N\longrightarrow K[1].
\]
Precisely, $\mathcal G$ is the pullback of
$\operatorname{Fun}(\Delta^1,D(A))\to D(A)\times D(A)$ along
$(N,K)\mapsto(j_*N,K[1])$.
Morphisms include the homotopies in the resulting commutative
squares, rather than merely an equality in a homotopy category.

\begin{proposition}
\label{v4-arith:ic:w5-h:localization-gluing}
There is an equivalence of stable $\infty$-categories
\[
 D(A)\simeq\mathcal G.
\]
It sends $M$ to $(\Gamma_tM,LM,\delta_M)$, where $\delta_M$ is the
connecting morphism of the localization triangle.
Its inverse sends $(K,N,\delta)$ to $\operatorname{fib}(\delta)$.
\end{proposition}
\begin{proof}
A triple defines the fiber sequence
\[
 K\longrightarrow \operatorname{fib}(\delta)
 \longrightarrow j_*N\xrightarrow{\delta}K[1].
\]
Applying $L$ identifies $L\operatorname{fib}(\delta)$ with $N$.
Under this identification, the middle map is the unit of $L\dashv j_*$:
by adjunction it corresponds to the identity of $N$.
Its fiber is therefore $\Gamma_t\operatorname{fib}(\delta)=K$.

These identifications recover the connecting morphism $\delta$.
Conversely, the fiber of the connecting map in a localization
triangle is its middle term. Fiber constructions are functorial
in a stable $\infty$-category, so these identifications give mutually
inverse functors, including their mapping spaces.
\end{proof}

For $M=A$, the complex $\Gamma_t A$ is represented by
$[A\to C]$ in degrees $0,1$.
Hence $\Gamma_t A\simeq(C/A)[-1]$, and the connecting map is
\[
 C\longrightarrow C/A.
\]
Its fiber is $A$ in degree zero. This nonzero quotient map is the
extension datum that reconstructs $A$.

The two objects $C$ and $\Gamma_t A$ are not perfect over $A$:
their nonzero cohomology modules are not finitely generated.
Thus this localization construction takes place in $D(A)$.
It does not present perfect complexes as a product of perfect
complexes on the two strata.
Nor can the supported category be replaced by $D(k)$ under closed
pushforward. Indeed, the resolution
$0\to A\xrightarrow{t}A\to k\to0$ gives
$\operatorname{Ext}^1_A(k,k)=k$, whereas
$\operatorname{Ext}^1_k(k,k)=0$.

\subsection{Coherent gluing across a formal neighborhood}

To retain coherent modules, replace the closed point by its formal
neighborhood. Put
\[
 B=k[[t]],\qquad D=k((t)),\qquad C=k[t,t^{-1}].
\]
The maps $A\to B$, $A\to C$, $B\to D$, and $C\to D$ are the
inclusions. The overlap $\operatorname{Spec}D$ is the punctured
formal disk.

It is not the empty intersection $Z\cap U$.
The ring $B$ is the coordinate ring of the $t$-adic formal disk.
Here its coherent category means finitely generated $B$-modules,
with their $t$-adic topology. Tensor products displayed below are
ordinary algebraic tensor products unless explicitly derived.

Let $\mathcal P$ be the category of triples
\[
 (P,Q,\eta),\qquad P\in\operatorname{mod}_{\mathrm{fg}}(B),\quad
 Q\in\operatorname{mod}_{\mathrm{fg}}(C),\quad
 \eta:D\otimes_B P\xrightarrow{\sim}D\otimes_C Q.
\]
Morphisms are compatible pairs of module homomorphisms.

\begin{theorem}
\label{v4-arith:ic:w5-h:formal-patching}
The functor
\[
 \operatorname{mod}_{\mathrm{fg}}(A)\longrightarrow\mathcal P,
 \qquad M\longmapsto(B\otimes_A M,C\otimes_A M,\mathrm{can})
\]
is an equivalence. Its inverse sends a triple to
\[
 \{(p,q)\in P\oplus Q:\eta(1\otimes p)=1\otimes q\}.
\]
\end{theorem}
\begin{proof}
Every Laurent series is the sum of a power series and a Laurent
polynomial. Their intersection is $k[t]$. Consequently
\[
 0\longrightarrow A\longrightarrow B\oplus C
 \xrightarrow{(b,c)\mapsto b-c}D\longrightarrow0
\]
is exact. The modules $B,C,D$ are flat over $A$.

For $C,D$ this follows from localization and torsion-freeness over
the principal ideal domain $A$. A torsion-free module over a principal ideal domain is the
filtered union of its finite free submodules. It is therefore flat,
which also applies to $B$.
Since $D$ is flat, tensoring the displayed sequence with any
$A$-module $M$ remains exact.

In particular $M$ is the indicated
fiber product of its two scalar extensions.
A compatible pair of homomorphisms on these extensions sends this
fiber product into the corresponding fiber product for the target.
The resulting $A$-linear homomorphism recovers the pair, since the
elements $1\otimes m$ generate both scalar extensions.
This proves full faithfulness.

For essential surjectivity, use the structure theorem for finitely
generated modules over the principal ideal domains $B$ and $C$.
A torsion summand $B/(t^m)$ descends from $A/(t^m)$, whose scalar
extension to $C$ vanishes. A torsion summand of $Q$ has the form
$C/(f)$, where $f\in k[t]$ and $f(0)\ne0$.
It descends from $A/(f)$, whose extension to $B$ vanishes because
$f$ is a unit in $B$. Both kinds of torsion disappear over $D$.
After splitting off these summands, $P$ and $Q$ are free of the
same rank $r$.

Choose $Q=C^r$ and view $P$ as a $B$-lattice
$\Lambda\subset D^r$ through $\eta$.
There is an integer $n\ge0$ such that
\[
 t^nB^r\subset\Lambda\subset t^{-n}B^r.
\]
Set $M=\Lambda\cap C^r$. Intersecting the bounds with $C^r$ gives
\[
 t^nA^r\subset M\subset t^{-n}A^r.
\]

Thus $M$ is a finite free $A$-module of rank $r$, and
$C\otimes_A M=C^r$.
For $\lambda\in\Lambda$, truncate its Laurent expansion below
degree $n$. The truncation $q$ is a Laurent polynomial, and
$\lambda-q\in t^nB^r\subset\Lambda$. Hence $q\in M$ and
$\Lambda=M+t^nB^r$.

The map $B\otimes_A M\to\Lambda$ is therefore
surjective. It is injective: an $A$-basis of $M$ becomes a basis
of $C^r$, so its determinant remains nonzero over the field $D$.
This proves $B\otimes_A M=\Lambda$ with the specified overlap map.
Adding the torsion summands proves essential surjectivity.
Full faithfulness identifies the stated fiber-product construction
with a quasi-inverse on all triples.
\end{proof}

For a $k$-linear category $\mathcal C$, write
$Z(\mathcal C)=\operatorname{End}(\operatorname{id}_{\mathcal C})$
for its ordinary center, with multiplication given by composition.

\begin{corollary}
\label{v4-arith:ic:w5-h:ordinary-patching-centre}
Under the preceding equivalence,
\[
 Z(\operatorname{mod}_{\mathrm{fg}}(A))
 \simeq B\times_D C
 \simeq Z(\operatorname{mod}_{\mathrm{fg}}(B))
       \times_{Z(\operatorname{mod}_{\mathrm{fg}}(D))}
       Z(\operatorname{mod}_{\mathrm{fg}}(C)).
\]
The comparison maps on centers are the indicated inclusions of rings.
\end{corollary}
\begin{proof}
For every commutative ring $R$, a natural endomorphism $z$ of
$\operatorname{mod}_{\mathrm{fg}}(R)$ is multiplication by
$a=z_R(1)\in R$. Indeed, for each $m\in M$, apply naturality to
$R\to M$, $1\mapsto m$. This yields $z_M(m)=am$.
Thus the center is $R$, compatibly with scalar extension.
The exact ring sequence above gives $A=B\times_D C$.
\end{proof}


\subsection{Effective equivariant descent}

Let \(X\) be a specified global parameter space in a fixed
coefficient theory, and let \(U_\bullet\to X\) be an augmented
simplicial diagram. Fix the coherent or ind-coherent support
subcategories being used. Suppose restriction is an equivalence
\[
 \mathcal C(X)\xrightarrow{\sim}
       \operatorname{Tot}\mathcal C(U_\bullet).
\]
An object of the totalization includes the objects on the charts,
their isomorphisms on overlaps and every higher compatibility.

Applying the inverse equivalence constructs its global object.
Full faithfulness identifies the mapping spaces and gives uniqueness
as a descent object. Coherence follows from this equivalence.
It does not follow from an arbitrary filtered colimit of coherent
objects in an ind-category.

For a derived pairwise construction there are two separate
Tor conditions if ordinary tensors are to compute it.
One condition makes the external product of the two local objects
ordinary. The other makes its pullback to the global base ordinary.
Perfectness of the local objects and flatness of the pullback
do not imply the first condition.
For \(W=k[[t]]\), both local modules \(k\) are perfect over \(W\),
but
\[
 k\otimes_W^{\mathbb L}k\simeq k\oplus k[1].
\]
An identity pullback is flat and retains this negative-degree term.

Suppose, in addition, a compatible family of algebra actions
\(A_S\to Z^0(\mathcal C)\) is given for every finite set \(S\)
of places. The ordinary center here consists of natural
endomorphisms of the identity on the homotopy category.
Their colimit \(A=\operatorname*{colim}_S A_S\) acts, because
every element and relation comes from a finite stage.
An equivalence which intertwines each action intertwines the
assembled action by that same property.

For derived actions, the corresponding input consists of maps
of the specified derived algebras and coherent comparisons.
The identity
\[
 \operatorname{Map}(\operatorname*{colim}_S A_S,Z)
 \simeq\operatorname*{lim}_S\operatorname{Map}(A_S,Z)
\]
assembles these comparisons, including their homotopies.
Equal degree-zero actions alone do not supply them.
All finite sets enter this indexing system. The poset of pairs
alone is not filtered.

\begin{question}\label{v4-arith:ic:arithmetic-descent}
For the continuous Selmer deformation problem over one coefficient
base, construct the geometric charts, overlap correspondences and
effective equivariant descent above. Compare the local spectral
objects with their arithmetic coefficient objects, preserving the
full local parameters, the specified support and the global actions.
If a prismatic refinement is required, construct its filtered
objects and cross-prime maps on this same recipient.
\end{question}

The ordinary formal pullback modules and their central actions
have already been constructed. This question asks for their
geometric identification and descent. An actual Nygaard filtration
on a local prism does not, by itself, provide cross-prime maps.


\section{Centers, categorical traces and character classes}
\label{v4-arith:ic:traces}

Fix a commutative coefficient ring \(k\). Let \(\mathcal C\)
be a stable presentable \(k\)-linear monoidal category whose
tensor preserves colimits in each variable. Its categorical center
and trace are
\[
 Z_{\mathrm{cat}}(\mathcal C)
 =\operatorname{Fun}^{\mathrm L}_{\mathcal C\otimes\mathcal C^{\mathrm{op}}}
                     (\mathcal C,\mathcal C),
 \qquad
 \operatorname{Tr}_{\mathrm{cat}}(\mathcal C)
 =\mathcal C\otimes_{\mathcal C\otimes\mathcal C^{\mathrm{op}}}
                  \mathcal C.
\]
Here tensors are in \(k\)-linear presentable categories, and
the functors preserve colimits.

Both expressions are categories.
The center is braided and the trace has
a circle action.
For an associative \(k\)-algebra \(A\), the scalar Hochschild
complexes instead have the forms
\[
 C^\bullet(A,A)
   =\operatorname{RHom}_{A\otimes^{\mathbb L}A^{\mathrm{op}}}(A,A),
 \qquad
 C_\bullet(A)
   =A\otimes_{A\otimes^{\mathbb L}A^{\mathrm{op}}}^{\mathbb L}A.
\]
Their degree-zero center and trace are still different objects.
An identity trace of a chosen perfect module is a class in
\(HH_0\), not the whole Hochschild complex.

For a perfect derived stack \(X\), Ben-Zvi--Francis--Nadler,
Theorem 1.7, prove
\[
 Z_{\mathrm{cat}}(\mathrm{QCoh}(X))
 \simeq\mathrm{QCoh}(\mathcal LX)
 \simeq\operatorname{Tr}_{\mathrm{cat}}(\mathrm{QCoh}(X)).
\]
Here perfect means that \(X\) has affine diagonal and
\(\mathrm{QCoh}(X)=\operatorname{Ind}\mathrm{Perf}(X)\).
For \(X=\operatorname{Spec}A\), scalar Hochschild chains of
\(\mathrm{Perf}(A)\) identify with functions on the derived
self-intersection
\(X\times_{X\times X}^{\mathbb L}X\).
This affine algebra calculation and the categorical theorem
occur at different levels.

A commutative algebra can still have nontrivial Hochschild
information: \(HH_1(k[t])=k[t]\,dt\).
Nor does an \(E_1\)-algebra automatically have \(E_2\) operations.
These observations rule out deductions of arithmetic trace
identities from spherical commutativity alone.

\subsection{The finite Hecke character theorem}

Let \(G\) be a complex reductive group, \(B\) a Borel subgroup,
and \(N\) its unipotent radical.
Put
\[
 \mathcal H_G=D(B\backslash G/B).
\]
Its unipotently monodromic version
\(\widetilde{\mathcal H}_G\) is the full colimit-generated
subcategory of \(D(N\backslash G/N)\) generated by pullbacks
from \(\mathcal H_G\).
Both use convolution.

The horocycle correspondence
\[
 G/\operatorname{Ad}G
 \xleftarrow{p}G/\operatorname{Ad}B
 \xrightarrow{\delta}B\backslash G/B
\]
defines the Harish--Chandra transform \(p_*\delta^!\).
Let \(\mathrm{Ch}_G\) be the full subcategory of
\(D(G/\operatorname{Ad}G)\) generated under colimits by its image.
Ben-Zvi--Nadler, Theorem 1.8 of
\emph{The character theory of a complex group}, prove
\[
 Z_{\mathrm{cat}}(\mathcal H_G)
 \simeq\mathrm{Ch}_G
 \simeq\operatorname{Tr}_{\mathrm{cat}}(\mathcal H_G),
\]
and the same for \(\widetilde{\mathcal H}_G\).
These Hecke categories are semi-rigid and have the specified
pivotal and Calabi--Yau structures.
Semi-rigid means compactly generated with dualizable compact
generators. It does not require a compact unit.

A dualizable Hecke module has a character object in
\(\mathrm{Ch}_G\).
This is a positive categorical character construction.
It is not an identification of \(HH(D(G/\operatorname{Ad}G))\)
with functions on an unspecified flat locus.
An arithmetic identity trace requires its own source category,
chosen object and comparison to the intended arithmetic recipient.

\subsection{Continuity and arithmetic trace comparisons}

There is an elementary continuity model on a small differential
graded presentation. In an orthogonal union of small \(K\)-flat
dg categories, every cyclic composable word stays in one block.
The Hochschild complex is therefore the direct sum of the block
complexes, compatibly with all faces, degeneracies and the
Connes operator. Perfect completion retains finite block support.
This proves the corresponding circle-equivariant trace formula
on that carrier.

An infinite coordinatewise sum of nonzero unital algebras has
no unit, and a categorical product is not that algebraic sum.
A multiplicative trace comparison consequently needs its
actual monoidal enhancement.
For a filtered dg diagram, a Hochschild continuity statement
likewise requires its compact or Morita presentation and the
specified maps on cyclic bars.
Neither a cofiltered limit of geometric spaces nor its functions
is exchanged with an arithmetic completed trace by notation.

\begin{question}\label{v4-arith:ic:trace-comparison}
Construct the arithmetic representation or Hecke category,
its chosen object and identity trace, and an operation-preserving
map to the specified geometric trace recipient.
Compare the coefficient theory, completion, singular support and
circle action. If a scalar evaluation is prescribed, prove that
the map intertwines it on its stated domain.
\end{question}


\section{Restricted tensors and local actions}
\label{v4-arith:ic:restricted-tensors}

A finite collection of local categories has a tensor product over
a common coefficient category. Passing to all places requires maps
which insert the omitted factors. These maps must preserve the
operations being assembled.

Write \(\mathcal V=\mathrm{Pr}^{\mathrm L}_{\mathrm{st}}\) for the
symmetric monoidal \(\infty\)-category of stable presentable categories
and colimit-preserving functors. Its tensor represents functors
preserving colimits separately in each input. Its unit is
\(\mathrm{Sp}\), with sphere object \(\mathbb S\).
The same construction can be made in a specified linear version,
whose tensor unit is the corresponding module category.

Fix a small coherent one-colored \(\infty\)-operad \(\mathcal O\),
for example \(E_n\) or \(E_\infty\). Let \(V\) be a set of places,
and let \(E\subset V\) be finite. For each \(v\), give
\(\mathcal C_v\) an \(\mathcal O\)-algebra structure in \(\mathcal V\).
For \(v\notin E\), supply a morphism
\[
 \eta_v:\mathrm{Sp}\longrightarrow\mathcal C_v
       \quad\text{in }\mathrm{Alg}_{\mathcal O}(\mathcal V),
 \qquad u_v=\eta_v(\mathbb S).
\]
The map includes every operation and its coherences. In the unital
\(E_n\) cases it sends the unit to the unit, so
\(u_v\simeq\mathbb 1_{\mathcal C_v}\).

Let \(J_E\) be the filtered poset of finite subsets \(S\subset V\)
containing \(E\). Define
\[
 \mathcal C_S=\bigotimes_{v\in S}\mathcal C_v,\qquad
 \iota_{S,S'}=
 \operatorname{id}\otimes\bigotimes_{v\in S'\setminus S}\eta_v.
\]
For a chain \(S\subset S'\subset S''\), both composites insert
the same maps at the same places. Tensor coherence supplies their
identification and all higher compatibilities. Thus these are a
diagram of \(\mathcal O\)-algebras.

\begin{theorem}\label{v4-arith:ic:restricted-colimit}
The underlying restricted tensor
\[
 \mathcal C_V'=\operatorname*{colim}_{S\in J_E}\mathcal C_S
       \quad\text{in }\mathcal V
\]
is the underlying object of the colimit in
\(\mathrm{Alg}_{\mathcal O}(\mathcal V)\).
Every finite-stage map is an \(\mathcal O\)-algebra map.
Fixing the full diagram and its underlying colimit cocone, the
space of lifts to this structured colimit cocone is contractible.

A canonical single-place map \(\mathcal C_v\to\mathcal C_V'\)
is supplied directly when \(E\subset\{v\}\).
In general, such a map additionally requires pointings
\(\mathrm{Sp}\to\mathcal C_w\) for \(w\in E\setminus\{v\}\).
The theorem makes no full-faithfulness assertion about these maps.
\end{theorem}

\begin{proof}
Tensor product in \(\mathcal V\) preserves colimits separately.
For a filtered diagram and each finite arity, its diagonal indexing
functor is cofinal. Hence tensor powers commute with its colimit.
The forgetful functor from \(\mathcal O\)-algebras therefore creates
filtered colimits, by the operadic colimit theorem
of \emph{Higher Algebra}, \S3.2.3.
Apply it to the displayed diagram. The structure maps of its colimit
are morphisms of \(\mathcal O\)-algebras.

An initial object in a category of cocones has a contractible space
of choices. A comparison between two structured colimit cocones
has the underlying equivalence determined by their fixed colimit
cocones. The forgetful functor detects equivalences.
Thus the space asserted in the theorem concerns the entire
coherent cocone, and is contractible.

If \(E\subset\{v\}\), the singleton is an object of \(J_E\).
Otherwise the extra maps give
\[
 \mathcal C_v\longrightarrow
 \bigotimes_{w\in E\cup\{v\}}\mathcal C_w
 \longrightarrow\mathcal C_V'.
\]
The first arrow inserts the additionally supplied pointings.
Its compatibility with larger finite sets follows from tensor
coherence. It depends on these choices.
\end{proof}

The need for a unital map is already visible at one place.
In \(\mathcal C=D(k)\times D(k)\), with componentwise tensor,
the object \(u=(k,0)\) is a compact central idempotent algebra
under the unit \((k,k)\). The insertion \(D(k)\to\mathcal C\)
with value \(u\) does not preserve that unit.
One can instead use the monoidal subcategory with unit \(u\),
or a separately specified lax or nonunital theory.
These change the recipient of the construction.

Nor is a unit insertion generally fully faithful.
Extension of scalars \(D(k)\to D(k[t])\) sends the unit to the
unit, but its map on unit endomorphisms is \(k\to k[t]\).
Unitality and full faithfulness are different conditions.

\subsection{Factorization, coefficient change and iteration}

For a partition \(V=V_1\sqcup V_2\), finite subsets separate as pairs
of finite subsets. Cofinality and separate colimit preservation give
\[
 \mathcal C_V'\simeq
 \mathcal C_{V_1}'\otimes\mathcal C_{V_2}',
\]
using \(E\cap V_i\) in each factor.
The same proof handles every finite partition, with its associativity
and permutation coherences: all maps are induced by the corresponding
finite tensor diagrams.

Suppose colimit-preserving \(\mathcal O\)-algebra maps
\(F_v:\mathcal C_v\to\mathcal D_v\) come with coherent
identifications \(F_v\eta_v\simeq\eta_v'\).
Their finite tensor products commute with every insertion.
Taking their colimit gives the induced map of restricted tensors.
An enlargement of the exceptional finite set restricts to a cofinal
subdiagram and leaves the colimit unchanged.

If \(B:\mathcal V\to\mathcal W\) is a symmetric monoidal,
colimit-preserving functor between the chosen coefficient categories,
then
\[
 B(\mathcal C_V')\simeq
 \operatorname*{colim}_{S\in J_E}
     \bigotimes_{v\in S}B(\mathcal C_v).
\]
The comparison retains \(B(\eta_v)\), including their nullary
compatibilities. This is the coefficient-change map.
The same argument commutes two filtered systems of finite subsets.
It gives iterated restricted tensors whenever the stated pointings
define both diagrams and the paired indexing systems are cofinal.

Dunn additivity identifies an \(E_{m+n}\)-structure with a compatible
\(E_m\)-structure in \(E_n\)-algebras.
When this full structure is supplied at each finite stage and on
the insertion maps, the preceding colimit transports it.
An \(E_m\)-structure by itself supplies no \(E_{m+n}\)-structure.
A matrix algebra in degree zero is a simple test: an \(E_2\)-algebra
has commutative degree-zero homology, whereas
\(\operatorname{Mat}_2(k)\) is noncommutative.

\subsection{Module categories and compact objects}

For \(E_\infty\)-ring spectra \(R_v\) over a common \(E_\infty\)-ring
\(A\), extension of scalars gives the pointings of
\(\mathrm{Mod}_A\)-linear module categories.
The universal property of the module-category tensor gives
\[
 \bigotimes_{v\in S,\mathrm{Mod}_A}\mathrm{Mod}_{R_v}
 \simeq
 \mathrm{Mod}_{\bigotimes_{v\in S,A}R_v}.
\]
Here the tensor on rings is the derived spectral tensor over \(A\).
The finite comparisons are natural under unit insertion.
Module categories preserve this filtered algebra colimit, so
\[
 \bigotimes_{v\in V}'\mathrm{Mod}_{R_v}
 \simeq
 \mathrm{Mod}_{\,\operatorname*{colim}_{S}
                    \bigotimes_{v\in S,A}R_v}.
\]
This formula does not replace a tensor of rings by their product,
or an arithmetic completion by an uncompleted tensor.

For compact generation, suppose every \(\mathcal C_S\) is compactly
generated and each insertion preserves compact objects.
Write \(\mathcal C_S^c\) for its small idempotent-complete stable
category of compacts. The colimit is then
\[
 \mathcal C_V'\simeq
 \operatorname{Ind}\!
 \left(\operatorname*{colim}_{S\in J_E}
                    \mathcal C_S^c\right),
\]
where the inner colimit is taken in small idempotent-complete stable
categories. Indeed, Ind-completion identifies that category with
compactly generated presentable stable categories and
compact-preserving colimit functors.

The universal property on both
sides gives the displayed equivalence.
A selected spherical generator only produces the sector it generates.
Taking \(\operatorname{Ind}(\mathcal C^c)\) does not remove a
continuous spectrum from an already compactly generated category.

\subsection{Actions at separate places}

Let \(T_v:\mathcal C_v\to\mathcal C_v\) be an exact
colimit-preserving endofunctor at one place.
The subposet of finite sets containing \(v\) is cofinal.
On that diagram \(T_v\) acts on the \(v\)-factor and the identity
acts on the other factors. It commutes with every insertion,
which occurs away from \(v\).
Thus it defines an exact colimit-preserving endofunctor on
\(\mathcal C_V'\).

Finitely many such actions extend on the cofinal diagram containing
all their places. The tensor interchange gives their coherent
commutation. No monoidal structure on the individual endofunctors
is needed. A construction on a diagram which omits an acting place,
or a specified simultaneous infinite action, additionally needs
compatibility with the relevant pointings.
A Hecke correspondence is therefore usable through its actual
endofunctor. It need not be a unital monoidal endofunctor.

\subsection{Finite labeled disks}

The finite-set tensor diagram gives a discrete factorization
construction: disjoint sets of labels tensor independently.
A geometric factorization algebra additionally requires a space
and categories on its opens or configurations.
It also requires maps at their collisions.

There is also a distinction within a finite disk construction.
Putting a fixed algebra \(A_S=\bigotimes_{v\in S}A_v\) on each of
\(r\) disjoint disks gives \(A_S^{\otimes r}\).
Putting \(A_v\) on one disk carrying label \(v\) gives one factor
per label. Their products need not agree.

For example, for \(A_S=k[x]\), two copies give \(k[x_1,x_2]\),
whereas a single labeled factor gives \(k[x]\).
A varying-label factorization diagram must retain this assignment
and the maps for every labeled embedding.
The finite-set colimit theorem alone supplies no arithmetic
collision comparison.


\section{Tensor sectors and extension data}

An orthogonal decomposition into local sectors gives a direct-sum
decomposition of their tensor product. Even with two sectors
at each of two places, tensor distributivity gives four terms:
\[
 (\mathcal A_0\oplus\mathcal A_1)\otimes
 (\mathcal B_0\oplus\mathcal B_1)
 \simeq
 (\mathcal A_0\otimes\mathcal B_0)\oplus
 (\mathcal A_0\otimes\mathcal B_1)\oplus
 (\mathcal A_1\otimes\mathcal B_0)\oplus
 (\mathcal A_1\otimes\mathcal B_1).
\]
The middle two terms record different sector choices at the two places.
A relation between those choices requires additional structure.

\subsection{Orthogonal decompositions}

Let $\mathcal V$ be a symmetric monoidal $\infty$-category with small
colimits. Assume its tensor product preserves colimits separately in each
variable. Write $\mathbf1$ for its tensor unit and $\coprod$ for its
categorical coproduct. The following argument applies to the
$\infty$-category of presentable stable $\infty$-categories and
colimit-preserving functors.

Its tensor represents functors preserving colimits separately
in each variable. See Lurie, \emph{Higher Algebra} (2017),
Remarks~4.8.1.18 and~4.8.1.24, and Proposition~4.8.2.18.
For finite families of stable categories, this coproduct is their
orthogonal direct sum, denoted $\oplus$.

Indeed, in a finite product $\prod_i\mathcal C_i$, every object is the
finite direct sum of its components. A colimit-preserving functor out of
this product is therefore determined by its restrictions to the factors.
Conversely, functors $F_i:\mathcal C_i\to\mathcal D$ extend by
$(x_i)_i\mapsto\bigoplus_iF_i(x_i)$. This proves the coproduct universal
property. Distinct factors have zero mapping spectra.

\begin{proposition}[Finite sector decomposition]
\label{v4-arith:ic:eis:finite-sectors}
Let $S$ be a finite set. For each $v\in S$, fix a finite set $I_v$ and a
specified equivalence
$\mathcal C_v\simeq\coprod_{r\in I_v}\mathcal C_{v,r}$ in $\mathcal V$.
There is an equivalence, induced by the given sector inclusions,
\[
 \bigotimes_{v\in S}\mathcal C_v
 \simeq
 \coprod_{\alpha\in\prod_{v\in S}I_v}
       \bigotimes_{v\in S}\mathcal C_{v,\alpha(v)}.
\]
For $S=\varnothing$, both sides are $\mathbf1$.
\end{proposition}

\begin{proof}
For one factor the assertion is the specified decomposition. Tensoring
that decomposition with another factor preserves its coproduct. Applying
the same fact to the new factor gives the Cartesian product of the two
index sets. Induction proves the formula. The maps are induced by the
sector inclusions, so the equivalence respects maps that preserve the
specified decompositions.
\end{proof}

Write $D(k)$ for the stable $\infty$-category of unbounded cochain
complexes of $k$-vector spaces, with quasi-isomorphisms inverted.
For example, in complex-linear presentable stable categories, put
$\mathcal A_i=\mathcal B_j=D(\mathbb C)$ and tensor over $D(\mathbb C)$.
This category is the tensor unit, so the result is $D(\mathbb C)^{\oplus4}$.
Keeping only the two equal-index factors gives $D(\mathbb C)^{\oplus2}$.
These categories are not equivalent.

Their compact objects have
Grothendieck groups $\mathbb Z^4$ and $\mathbb Z^2$, respectively.
Here the Grothendieck group imposes the relation $[y]=[x]+[z]$ for each
exact triangle $x\to y\to z$. In $\operatorname{Perf}(\mathbb C)$,
splitting a bounded complex into its cohomology identifies this group
with $\mathbb Z$, by Euler characteristic. Finite products give the
stated ranks. An equivalence preserves compact objects and exact
triangles, hence would preserve these groups.

The restricted tensor construction also depends on the sector of each
insertion. Let $V$ be a small set. For each $v\in V$, fix a finite nonempty set
$I_v$, an element $0_v\in I_v$, and a specified decomposition of
$\mathcal C_v$ as above. Fix a morphism
$e_v:\mathbf1\to\mathcal C_{v,0_v}$ in $\mathcal V$.
For finite $S\subset V$, set
\[
 \mathcal C_S=\bigotimes_{v\in S}\mathcal C_v.
\]

For $S\subset T$, tensor with the morphisms $e_v$, followed by sector
inclusion, for $v\in T\setminus S$. The monoidal coherence maps give a
diagram on the poset of finite subsets of $V$. Define its restricted
product by
\[
 \mathcal C^{\mathrm{res}}=\mathop{\operatorname{colim}}_{S\subset V,
                   \,S\text{ finite}}\mathcal C_S.
\]
This is an algebraic colimit in $\mathcal V$. No analytic or topological
completion is implicit.

Let $I^{\mathrm{fin}}$ consist of functions $\alpha$ with
$\alpha(v)\in I_v$ whose support
$\operatorname{supp}(\alpha)=\{v:\alpha(v)\ne0_v\}$ is finite.
For such $\alpha$, define
\[
 \mathcal C^{\mathrm{res}}_\alpha
 =\mathop{\operatorname{colim}}_{\substack{S\subset V\text{ finite}\\
                       S\supset\operatorname{supp}(\alpha)}}
       \bigotimes_{v\in S}\mathcal C_{v,\alpha(v)},
\]
with the same distinguished insertions.

\begin{proposition}[Restricted sector decomposition]
\label{v4-arith:ic:eis:restricted-sectors}
Under these hypotheses, the sector inclusions induce an equivalence
\[
 \mathcal C^{\mathrm{res}}
 \simeq\coprod_{\alpha\in I^{\mathrm{fin}}}
                   \mathcal C^{\mathrm{res}}_\alpha.
\]
\end{proposition}

\begin{proof}
For each $\alpha$, extend its displayed diagram to all finite $S$ by
assigning the initial object when $S$ does not contain its support.
Call the resulting diagram $E_\alpha$. At a fixed finite $S$, only
functions supported in $S$ contribute noninitial terms. The finite sector
decomposition therefore gives
$\mathcal C_S\simeq\coprod_\alpha E_\alpha(S)$.

The insertion lies in sector $0_v$, so the transition maps preserve this
identification of diagrams. Colimits commute with coproducts. Thus
\[
 \operatorname{colim}_S\mathcal C_S
 \simeq\coprod_\alpha\operatorname{colim}_S E_\alpha(S).
\]
The finite sets containing a fixed support form a cofinal subposet:
$S\mapsto S\cup\operatorname{supp}(\alpha)$ supplies a finite upper
bound compatible with inclusions. Each inner colimit is consequently
$\mathcal C^{\mathrm{res}}_\alpha$.
\end{proof}

A general insertion need not lie in one sector. If it has several
components, its transition maps need not preserve the finite-support
indexing above. The finite sector formula still holds. The restricted
colimit must then retain the actual transition maps between those
finite decompositions.

\subsection{A nonsplit extension}

Orthogonal direct sums exclude extensions between distinct factors.
The quiver $1\to2$ gives a concrete extension that must be retained when
orthogonality is absent. Let $k$ be a field, with no restriction on its
characteristic. A representation is a linear map $u:A\to B$ of
$k$-vector spaces. A morphism is a pair $(f,g)$ satisfying $gu=u'f$.
Kernels and cokernels are computed at the two vertices.

Consider
\[
 S_1=(k\to0),\qquad S_2=(0\to k),\qquad P=(k\xrightarrow{1}k).
\]
There is an exact sequence of representations
\[
 0\longrightarrow S_2\longrightarrow P\longrightarrow S_1
   \longrightarrow0.
\]
It does not split: a morphism $S_1\to P$ has source component
$f:k\to k$ satisfying $f=0$. It cannot be a section.

Moreover, $\operatorname{Hom}(P,(A\to B))\cong A$ and
$\operatorname{Hom}(S_2,(A\to B))\cong B$.
Both functors are exact, so $P$ and $S_2$ are projective.
The sequence is therefore a projective resolution of $S_1$.
Applying $\operatorname{Hom}(-,S_2)$ gives
\[
 \operatorname{Hom}(P,S_2)=0,\qquad
 \operatorname{Hom}(S_2,S_2)=k,\qquad
 \operatorname{Ext}^1(S_1,S_2)=k.
\]

The displayed sequence represents the image of $1\in k$ under the
connecting isomorphism. Projectivity of $S_2$ also gives
$\operatorname{Ext}^n(S_2,S_1)=0$ for every $n\ge0$.
Thus one direction is orthogonal, while the opposite direction carries
an extension.

Equivalently, these are left modules over
\[
 T=\left\{\begin{pmatrix}a&0\\m&b\end{pmatrix}:a,b,m\in k\right\},
 \qquad
 \begin{pmatrix}a&0\\m&b\end{pmatrix}
 \begin{pmatrix}a'&0\\m'&b'\end{pmatrix}
 =\begin{pmatrix}aa'&0\\ma'+bm'&bb'\end{pmatrix}.
\]
The two diagonal idempotents recover $A$ and $B$. The lower-left matrix
unit acts by $u:A\to B$. Associativity gives exactly the representation
relations. Replacing this algebra by $k\times k$ removes its lower-left
entry and the extension above.

\subsection{Stable gluing}

To retain extensions in a stable category, the connecting morphism must
be part of each object. Let $\mathcal A$ and $\mathcal B$ be stable
$\infty$-categories and let $F:\mathcal A\to\mathcal B$ be exact.
All mapping objects below are spectra, written
$\operatorname{map}_{\mathcal A}(-,-)$ and similarly for other categories.
The shift $[1]$ denotes suspension. For cochain complexes it has
$K[1]^n=K^{n+1}$ and differential $-d_K$.

Define $\mathcal G(F)$ to be the pullback
\[
 \mathcal G(F)
 = (\mathcal A\times\mathcal B)
   \mathop{\times}_{\mathcal B\times\mathcal B}
   \operatorname{Fun}(\Delta^1,\mathcal B),
\]
where the first map is $(a,b)\mapsto(Fa,b)$ and the second records the
source and target of an arrow. Thus an object is $(a,b,u)$ with
$u:Fa\to b$. A morphism consists of $f:a\to a'$, $g:b\to b'$, and a
specified homotopy from $u'Ff$ to $gu$. The pullback retains all higher
coherences.

\begin{proposition}[Mapping spectra and extensions]
\label{v4-arith:ic:eis:extension-category}
The category $\mathcal G(F)$ is stable, with finite limits and colimits
computed in the two components. For $x=(a,b,u)$ and $x'=(a',b',u')$,
\begin{equation}\label{v4-arith:ic:eis:sector-gluing-maps}
 \operatorname{map}_{\mathcal G(F)}(x,x')
 \simeq\operatorname{fib}\left(
 \operatorname{map}_{\mathcal A}(a,a')\oplus
 \operatorname{map}_{\mathcal B}(b,b')
 \xrightarrow{\ d\ }
 \operatorname{map}_{\mathcal B}(Fa,b')\right).
\end{equation}
Here \(d(f,g)=u'Ff-gu\).

The exact fully faithful functors
\[
 i(b)=(0,b,0),\qquad j(a)=(a,0,0)
\]
satisfy
\[
 \operatorname{map}(i(b),j(a))=0,\qquad
 \operatorname{map}(j(a),i(b))
 \simeq\Omega\operatorname{map}_{\mathcal B}(Fa,b)
 \simeq\operatorname{map}_{\mathcal B}(Fa[1],b).
\]
Every object has a natural exact triangle
\[
 i(b)\longrightarrow(a,b,u)\longrightarrow j(a)
       \xrightarrow{\delta_u}i(b)[1].
\]
Under the induced equivalence
$\operatorname{map}(j(a),i(b)[1])\simeq
\operatorname{map}_{\mathcal B}(Fa,b)$, the connecting morphism
$\delta_u$ corresponds to $u$. The choice of sign in this last
equivalence is fixed by this exact triangle.
\end{proposition}

\begin{proof}
Exactness of $F$ allows componentwise finite limits and colimits to
supply the required arrows. The zero components give a zero object.
Pushout squares are pullback squares componentwise, proving stability.
A map between two arrows is a homotopy-commutative square.

Its mapping
space is the homotopy equalizer of the two compositions $u'Ff$ and $gu$.
The same statement after every suspension gives the homotopy equalizer
of mapping spectra. In spectra this is the fiber of their difference,
which proves~\eqref{v4-arith:ic:eis:sector-gluing-maps}.

For two objects in the image of $i$, the formula reduces to
$\operatorname{map}_{\mathcal B}$. For two in the image of $j$, it
reduces to $\operatorname{map}_{\mathcal A}$.
The two cross formulas follow by substituting the zero components.
In particular, the fiber of $0\to E$ is $\Omega E$, not $E$.

The cofiber of $i(b)\to(a,b,u)$ is $(a,0,0)$ componentwise.
Its connecting morphism includes the homotopy encoded by $u$.
The suspension equivalence identifies that homotopy with $u$ and fixes
the stated sign convention.
\end{proof}

The triangle splits precisely when $\delta_u$ is nullhomotopic,
equivalently when $u$ is nullhomotopic. Indeed, a nullhomotopy identifies
its fiber with $i(b)\oplus j(a)$. Conversely, a splitting makes the
connecting morphism zero in the homotopy category. When $F=0$,
$\mathcal G(F)\simeq\mathcal A\times\mathcal B$.

Taking $\mathcal A=\mathcal B=D(k)$ and $F=\mathrm{id}$ gives the
arrow category of complexes. Here $h\mathcal G(F)$ denotes the homotopy category.
For $a=b=k$ in degree zero, the second
cross formula gives
\[
 \operatorname{Hom}_{h\mathcal G(F)}(j(k),i(k)[1])\cong k.
\]
The object $(k,k,1)$ gives its nonzero connecting class. This recovers
the extension mechanism in the quiver calculation, including its degree.

\subsection{Reconstruction from a semiorthogonal filtration}

The same description applies to a category whose extension structure is
already specified. Let $\mathcal C$ be stable and let
$i:\mathcal B\to\mathcal C$ and $j:\mathcal A\to\mathcal C$ be exact
fully faithful functors. Assume
\[
 \operatorname{map}_{\mathcal C}(i(b),j(a))=0
 \quad\text{for all }a\in\mathcal A,\ b\in\mathcal B.
\]
Assume every $x\in\mathcal C$ admits an exact triangle
\[
 i(b)\longrightarrow x\longrightarrow j(a)\longrightarrow i(b)[1].
\]
These hypotheses state the orientation of the semiorthogonal filtration.
They are stronger than generation by the two subcategories alone.

Let $\mathcal E$ be the full pullback of the arrow category of
$\mathcal C$ whose arrows have source $j(a)$ and target $i(b)[1]$.
Its objects are triples $(a,b,\delta)$ with
$\delta:j(a)\to i(b)[1]$. Then
\[
 \operatorname{map}_{\mathcal E}((a,b,\delta),(a',b',\delta'))
 \simeq\operatorname{fib}\left(
 \operatorname{map}_{\mathcal A}(a,a')\oplus
 \operatorname{map}_{\mathcal B}(b,b')
 \longrightarrow\operatorname{map}_{\mathcal C}(j(a),i(b')[1])
 \right),
\]
where the map sends $(f,g)$ to $\delta'j(f)-i(g)[1]\delta$.

\begin{proposition}[Semiorthogonal reconstruction]
\label{v4-arith:ic:eis:semiorthogonal-reconstruction}
Under these hypotheses, the functor
\[
 \mathcal E\longrightarrow\mathcal C,\qquad
 (a,b,\delta)\longmapsto\operatorname{fib}(\delta)
\]
is an equivalence. Suppose, in addition, that $i$ has an exact left
adjoint $L:\mathcal C\to\mathcal B$. Put $F=(Lj)[-1]$.
Then $\mathcal C\simeq\mathcal G(F)$, with the given embeddings $i,j$.
\end{proposition}

\begin{proof}
A connecting arrow has a fiber triangle
$i(b)\to\operatorname{fib}(\delta)\to j(a)$.
The filtration hypothesis gives essential surjectivity.
To prove full faithfulness, consider two such triangles with middle
objects $x,x'$. Applying $\operatorname{map}_{\mathcal C}(-,j(a'))$
to the first triangle gives an equivalence
\[
 \operatorname{map}_{\mathcal C}(j(a),j(a'))
       \xrightarrow{\sim}
 \operatorname{map}_{\mathcal C}(x,j(a')),
\]
because the mapping spectrum from $i(b)$ to $j(a')$ is zero.
Consequently a map $h:x\to x'$ determines, with contractible choices,
a map $j(f):j(a)\to j(a')$ and a homotopy between
$j(f)(x\to j(a))$ and $x\xrightarrow{h}x'\to j(a')$.

Taking fibers gives a map $i(g):i(b)\to i(b')$.
Taking cofibers supplies the compatible square of connecting arrows.
Conversely, a square of connecting arrows gives $h$ by taking fibers.
The fiber and cofiber universal properties make these constructions
inverse on mapping spaces, including all their higher homotopies.
Thus the functor is fully faithful.

For the additional assertion, adjunction and exactness give
\[
 \operatorname{map}_{\mathcal C}(j(a),i(b)[1])
 \simeq\operatorname{map}_{\mathcal B}(Lj(a),b[1])
 \simeq\operatorname{map}_{\mathcal B}(F(a),b).
\]
This natural equivalence identifies objects, squares, and their higher
coherences in $\mathcal E$ with those in $\mathcal G(F)$.
\end{proof}

In particular, a recollement can supply such a filtration without
supplying an orthogonal direct sum. Use the convention
$i_*:\mathcal B\to\mathcal C$ for the closed embedding and
$j_!:\mathcal A\to\mathcal C$ for the open extension functor.
The recollement adjunctions include $i^*\dashv i_*\dashv i^!$ and
$j_!\dashv j^*\dashv j_*$. Its defining vanishing $j^*i_*=0$ gives
$\operatorname{map}(j_!a,i_*b)=0$, and its exact triangle is
\[
 j_!j^*x\longrightarrow x\longrightarrow i_*i^*x
               \longrightarrow j_!j^*x[1].
\]

The first assertion of the proposition applies with the roles of the
closed and open categories interchanged. It retains the connecting
arrow $i_*i^*x\to j_!j^*x[1]$. It does not require a left adjoint to
$j_!$ or assert one exists.

The semiorthogonal description becomes an orthogonal direct sum if the
other cross mapping spectrum also vanishes. Then every connecting
spectrum is zero, and $\mathcal E\simeq\mathcal A\times\mathcal B$.
Without this vanishing, tensor assembly must retain the gluing functors
or connecting spectra and their composition data. The two sector names
alone do not determine the category.

For automorphic categories, applying these propositions requires actual
local decompositions or semiorthogonal filtrations and specified
insertion functors. An identification with a global automorphic category
also requires a comparison functor preserving the Hecke actions,
parameter constraints, and extension data. Neither tensor distributivity
nor semiorthogonal reconstruction constructs that comparison functor.


\section{Adelic factorization and automorphic realization}
\label{v4-arith:ic:adelic}

The restricted tensor construction has an algebraic representation
theorem as a model. For a restricted product of local groups, choose
nonzero reference vectors in the almost-everywhere spherical lines
of given local representations. The algebraic restricted tensor
is the colimit of finite tensors under insertion of these vectors.

Flath's factorization theorem says that every admissible
irreducible adelic representation has this form.
For a number field, the archimedean factors in its algebraic
form are \((\mathfrak g,K)\)-modules.
Cuspidality is not a hypothesis of this factorization theorem.
Changing the nonzero vectors in the almost-everywhere spherical
lines changes the concrete insertions but not the isomorphism
class of the restricted tensor.
It does not choose canonical vectors.

The unitary Hilbert restricted tensor is a different completion.
Its reference vectors are unit vectors and its inner product is
defined first on the algebraic finite tensors.
An infinite tensor without these reference choices has not been
specified by listing its local factors.

Neither construction places every independent tuple of local
representations in automorphic functions.
For \(GL_1\), the idele character must be trivial on the diagonal
number field's multiplicative group.
For \(GL_n\), automorphic realization likewise imposes its global
relations, growth and spectral conditions.
The local tuples and their tensor are positive constructions
before that additional comparison.

\subsection{The real place}

Let \(\mathfrak g_{\mathbb C}\) be the complexified Lie algebra
of the real group and let \(K\) be a maximal compact subgroup.
A Harish--Chandra module has compatible
\(\mathfrak g_{\mathbb C}\)- and \(K\)-actions. Admissibility asks
for finite multiplicity of each irreducible \(K\)-type.
It does not ask for only finitely many types.

The category of admissible modules is not closed under arbitrary
direct sums. An infinite sum of trivial representations has an
infinite-dimensional trivial \(K\)-isotype.
Consequently it is not the asserted presentable category merely
by being admissible.
One algebraic recipient is the Ind-completion of a chosen small
bounded derived category of finite-length Harish--Chandra modules.
Its analytic realization and comparison with Hilbert representations
require their own functors.

A spherical principal series is one object in a family, not
the tensor unit of this category.
On a cofinal finite-place diagram which always includes infinity,
the smallest stage is the specified archimedean category.
In a linear tensor theory the genuinely empty tensor is the
coefficient module category.

The real Weil group is the nonsplit extension generated by
\(\mathbb C^\times\) and \(j\), with
\[
 j^2=-1,\qquad jzj^{-1}=\bar z.
\]
For a lift \(zj\) of the nontrivial Galois element,
\((zj)^2=-|z|^2\ne1\), proving that this extension does not split.
Using \(\mathbb C^\times\rtimes C_2\) would change its
representation theory.

Strong approximation for \(SL_2\), with its specified omitted
places and noncompactness hypotheses, can simplify a group quotient.
For \(GL_2\) the determinant still gives the idele-class quotient.
Neither group-theoretic statement identifies all independent
local representation data with an automorphic category.

\subsection{Comparisons of enhanced categories}

Let \((\mathcal C_S)\) and \((\mathcal D_S)\) be supplied diagrams
of stable presentable categories on finite sets of places.
Suppose exact colimit-preserving equivalences
\(F_S:\mathcal C_S\simeq\mathcal D_S\) intertwine every insertion,
with their higher coherences.
Taking the colimit gives
\[
 \operatorname*{colim}_S\mathcal C_S
      \simeq\operatorname*{colim}_S\mathcal D_S.
\]
Indeed, their inverse diagrams give an inverse colimit functor,
and the coherent units and counits pass to these colimits.
This is a comparison of the supplied diagrams.
Identifying either colimit with an independently specified
geometric or automorphic category is additional data.

For a global functor
\(F:\mathcal C_{\mathrm{arith}}\to\mathcal D_{\mathrm{aut}}\),
Theorem~\ref{v4-arith:ic:w5-e:local:generator-criterion}
gives a sufficient test. Compare the full mapping spectra of
compact generators and preserve their compactness.
Prove that their images generate the intended target.
Natural compatibility with finite-place Hecke, ramified,
Whittaker and archimedean operations must also be supplied.
A match of central characters does not provide it.

The cuspidal, residual and continuous summands in a Hilbert
spectral decomposition have a precise analytic meaning.
Hilbert orthogonality does not imply vanishing of cross mapping
spectra in an enhanced category.
The nonsplit principal-series and adelic extensions above
exhibit the data a categorical gluing must retain.
Such a gluing uses its actual connecting functor, all mixed
finite-place sectors and coherent transitions.

\begin{question}\label{v4-arith:ic:global-comparison}
Construct an enhanced global automorphic category and a geometric
recipient for the fixed Selmer problem, together with the
comparison functor. Identify the local and global generators,
their full derived endomorphisms and the essential image.
Prove compatibility with finite and archimedean actions,
Whittaker realization, coefficient completion and the specified
trace or pairing.
\end{question}

The local full embeddings, determinant-ring reconstruction,
formal patching and restricted tensor theorems supply parts of
this construction. Their hypotheses do not state a necessity
theorem for every possible arithmetic realization.
A different geometric recipient or a restricted natural domain
remains possible when its comparison maps are constructed.


\section{Formal operators and arithmetic recognition}
\label{v4-arith:ic:formal-operators}

Formal spectral variables make sense over a common coefficient ring.
Their existence is separate from identifying the resulting algebra
with arithmetic operators.
In particular, \(\operatorname{Spec}\mathbb Z\) is not discrete:
its generic point is dense.
The lack of a chosen comparison between different local fields
does not prevent the formal expression \(u-v\) over \(\mathbb Z\).

\subsection{An integral RTT construction}

Let \(R\) be a commutative ring and set
\[
 T_{ij}(u)=\delta_{ij}+\sum_{r\ge1}t_{ij}^{(r)}u^{-r}.
\]
Define an \(R\)-algebra by the coefficient relations
\[
 (u-v)[T_{ij}(u),T_{kl}(v)]
  =T_{kj}(u)T_{il}(v)-T_{kj}(v)T_{il}(u).
\]
The denominator has been cleared, so these are integral relations
among the formal coefficients.
Equivalently, they are RTT with
\(\mathcal R(u-v)=1-P/(u-v)\), where \(P\) exchanges the two
auxiliary factors, in the corresponding rational formal notation.

\begin{proposition}\label{v4-arith:ic:rtt}
These relations admit the coproduct and counit
\[
 \Delta T_{ij}(u)=\sum_aT_{ia}(u)\otimes T_{aj}(u),
 \qquad \epsilon(T_{ij}(u))=\delta_{ij}.
\]
The formal matrix inverse defines the antipode.
There is an evaluation homomorphism to
\(U_R(\mathfrak{gl}_n)\) given by
\(T_{ij}(u)\mapsto\delta_{ij}+E_{ij}u^{-1}\).
\end{proposition}

\begin{proof}
Suppose \(T\) and \(S\) satisfy RTT and their coefficients
commute with one another. Move the \(R\)-matrix through
\(T_1T_2\), and then through \(S_1S_2\).
This proves RTT for their matrix product. The coefficients in the two tensor factors commute,
so the intermediate factors can be reordered as required.

This proves that the coproduct preserves the relations.
Associativity of matrix multiplication gives coassociativity.
The identity matrix gives the counit.

The constant term of \(T\) is the identity.
Its inverse is therefore defined coefficientwise, with each
coefficient a finite sum of products of the generators.
Inverting RTT gives the opposite-algebra relation required of
an antipode, and matrix inversion gives both antipode identities.
Finally, substitute the displayed evaluation. The relation reduces
to
\[
 [E_{ij},E_{kl}]
   =\delta_{jk}E_{il}-\delta_{il}E_{kj},
\]
which holds in the enveloping algebra over \(R\).
\end{proof}

This is a formal algebraic construction. It neither identifies
primes with spectral coordinates nor produces an arithmetic
coproduct on an independently specified Hecke category.
A proposed shuffle obstruction must likewise specify its kernel
and the identity which fails. A name for a zeta-valued kernel
is not an obstruction proof.

\subsection{Bethe algebras in specified representations}

A Bethe algebra is a commutative subalgebra of the Yangian.
Its image on a chosen representation is a quotient of that
subalgebra. It is not thereby a quotient of the entire Yangian.

The discrete Wronski theorem gives a precise positive model.
Fix distinct nonzero \(q_1,\ldots,q_N\in\mathbb C\) and
nonnegative integers \(\lambda_i\) with \(\sum_i\lambda_i=n\).
Let
\[
 f_i(u)=q_i^u
       (u^{\lambda_i}+p_{i,1}u^{\lambda_i-1}
                      +\cdots+p_{i,\lambda_i}).
\]
The \(p_{i,j}\) are coordinates on an affine \(n\)-space.

For \(\operatorname{Wr}(g_1,\ldots,g_N)
=\det(g_i(u-j+1))_{i,j}\), write
\[
 \operatorname{Wr}(f_1(u-1),\ldots,f_N(u-1))
 =\prod_iq_i^{u-1}
   \prod_{i<j}(q_j^{-1}-q_i^{-1})
   \left(u^n+\sum_{s=1}^n(-1)^s\Sigma_su^{n-s}\right).
\]
The map given by the \(\Sigma_s\) is the discrete Wronski map.
Let \(\mathcal O_{\lambda,\mathbf a}^q\) be the coordinate ring
of its scheme-theoretic fiber at
\(\mathbf a=(a_1,\ldots,a_n)\).

Choose \(b_1,\ldots,b_n\in\mathbb C\) with
\(b_i\ne b_j+1\) for \(i>j\), and put
\(a_s=e_s(b_1,\ldots,b_n)\).
For the tensor of vector evaluation modules
\[
 V(\mathbf b)=\mathbb C^N(b_1)\otimes\cdots
                              \otimes\mathbb C^N(b_n),
\]
Mukhin--Tarasov--Varchenko, Theorem 5.2 and Corollary 5.4,
identify \(V(\mathbf b)_\lambda\) with the regular module of
\(\mathcal O_{\lambda,\mathbf a}^q\).
This identifies its Bethe action for the diagonal twist \(q\).

The resulting Bethe image is maximal commutative and Frobenius.
Their Theorem 6.3 and Corollary 6.4 give the polynomial
Wronski version for the untwisted singular weight space.
These specified fibers and evaluation modules are the recipients
of the theorem.

In contrast, a local spherical Hecke algebra for \(GL_n\) is a
Laurent symmetric Satake algebra with its chosen normalization.
Its determinant generator is invertible.
Already for \(GL_1\), valuation translation acts on a bilateral
lattice, whereas a unilateral creation operator has no inverse
at the vacuum.
Arbitrary characters of a tensor of local Satake algebras specify
local parameter tuples. Their global automorphic and Galois
realizations still require the maps of
Question~\ref{v4-arith:ic:global-comparison}.
An oper or critical-center comparison must also specify its
actual Lie algebra, representation and spectral curve.

\section{Scalar invariants from chiral geometry}
\label{v4-arith:ic:scalar-examples}

A numerical invariant can be computed before its proposed geometric
realization is constructed. The examples below retain the scalar
calculation and identify the additional object it would have to describe.

\begin{example}[A rational function and its square root]
\label{v4-arith:ic:rational-shadow}
Define
\[
 W(c)=\frac{c^2(5c+22)}{4(45c+218)}.
\]
It has a simple pole at \(c=-218/45\), a double zero at \(c=0\),
and a simple zero at \(c=-22/5\). These statements follow from
the numerator and denominator, which have no common root.
For \(c>0\), let \(\Omega(c)\) be its positive square root.

Then
\[
 \Omega(c)=\frac c6-\frac1{27}+O(c^{-1})
       \qquad(c\longrightarrow+\infty).
\]
Indeed,
\[
 W(c)=\frac{c^2}{36}
       \left(1-\frac4{9c}+O(c^{-2})\right),
\]
and the positive square-root expansion gives the result.
This computation specifies a rational function and a real branch.
Identifying either with a Borel invariant requires its defining
transform and an equality on the common domain.
\end{example}

There is a separate differential identity. Let \(Q(t)\) and
\(H(t)\) be twice differentiable functions on a domain where \(t,Q,H\)
are nonzero, and suppose \(H^2=t^4Q\). Then
\[
 U=\frac{H'}H-\frac2t=\frac{Q'}{2Q},
 \qquad U'+2U^2=\frac{Q''}{2Q}.
\]
The first equality follows by logarithmic differentiation.

Differentiating \(Q'/(2Q)\) proves the second, including the
cancellation of \((Q')^2/(2Q^2)\).
The identity is valid for any such \(Q,H\).
It therefore does not identify a particular chiral operation by itself.

A weighting also changes the scalar coefficients. The two functions
\[
 S_6(c)=\frac{4(240c+1031)}{c^3(5c+22)^2},
 \qquad
 \widehat S_6(c)=\frac{80(45c+193)}{3c^3(5c+22)^2}
\]
satisfy
\[
 \widehat S_6(c)-S_6(c)
   =\frac{4(180c+767)}{3c^3(5c+22)^2}.
\]
They cannot be interchanged without a comparison of their
weighting conventions. If a formal series
\(H(t)=\sum_{r\ge2}S_rt^r\) satisfies
\([t^6]H^2=0\), its coefficients instead obey the literal equation
\[
 2S_2S_4+S_3^2=0.
\]
The equation \(2S_2S_6+2S_3S_5+S_4^2=0\) is the
coefficient equation in degree eight. Which degree vanishes
must be determined from the specified right-hand side of \(H^2\).

\begin{example}[The Bernoulli numerator \(691\)]
\label{v4-arith:ic:bernoulli}
Define \(B_n\in\mathbb Q\) by
\[
 \frac{x}{e^x-1}=\sum_{n\ge0}B_n\frac{x^n}{n!}.
\]
Multiplication by \(e^x-1\) gives \(B_0=1\) and
\[
 \sum_{j=0}^{m}\binom{m+1}{j}B_j=0
       \qquad(m\ge1).
\]
Successive substitution yields \(B_1=-1/2\), vanishing of the
odd values from \(B_3\) through \(B_{11}\), and
\[
 \begin{array}{c|rrrrrr}
 n&2&4&6&8&10&12\\ \hline
 B_n&1/6&-1/30&1/42&-1/30&5/66&-691/2730 .
 \end{array}
\]
This is an exact rational calculation.

It specifies neither a
two-dimensional Galois representation nor its determinant law.
A proposed representation attached to a shadow coefficient must
first be constructed, with its coefficient ring and continuous action.
Its characteristic polynomials can then be compared by the
determinant methods of this chapter.
\end{example}

An automorphic product has a similarly limited numerical test:
a seed coefficient \(c(0)\) gives the proposed weight \(c(0)/2\).
An identification with a chiral central invariant also requires
the product, its transformation law and the comparison map.
A form of half-integral or quarter-integral weight requires its
actual multiplier system or covering group.
The numerical weight alone chooses none of these structures.

\section{Physical realizations and scalar tests}
\label{v4-arith:ic:physical}

Gauge-theoretic constructions give further possible recipients
for arithmetic data.
The Kapustin--Witten construction starts with a smooth geometric
curve and a twisted four-dimensional gauge theory.
The Costello--Witten--Yamazaki construction uses a topological
surface and a complex curve with a chosen meromorphic one-form,
together with a partial connection and its boundary conditions.
Replacing that curve by an arithmetic compactification requires
new geometric and physical data.
The absence of those data does not prove nonexistence of every
arithmetic bulk theory.

The protected chiral construction of
Beem--Lemos--Liendo--Peelaers--Rastelli--van Rees assigns a chiral
algebra to the protected sector of a four-dimensional
\(\mathcal N=2\) superconformal theory.
Its normalization has
\(c_{2d}=-12c_{4d}\) and \(k_{2d}=-k_{4d}/2\).
A precise finite-current obstruction requires a finite-dimensional
weight-one space. Under that hypothesis, there cannot be an
infinite linearly independent family of weight-one currents,
by dimension alone.
This argument does not establish that hypothesis for every
possible superconformal theory.

The oscillator commutator and its logarithmic energy are independent
normalizations. For example,
\[
 [a_{p,m},a_{q,n}]
   =m\delta_{pq}\delta_{m+n,0},
 \qquad
 E=\sum_p(\log p)N_p
\]
retain level one at each prime while weighting the occupation
energy by \(\log p\).
They do not force Chern--Simons levels \(1/\log p\).
Nor does finite rank of a root lattice bound the number of states
or generators in all of its infinite-dimensional representations.

A stronger concrete automorphic model comes from the odd
Maaß eigenfunctions on the Dirichlet half-modular domain.
The Eisenstein continuum of the full modular quotient does not
satisfy that odd boundary condition.
For a fixed Hecke eigenstate, assume its normalized parameters
are \(e^{i\theta_p},e^{-i\theta_p}\).
Then on two charged oscillators at each prime,
\[
 \operatorname{Tr}\!
 \left(e^{i\theta_p(N_{p,1}-N_{p,2})}
       p^{-s(N_{p,1}+N_{p,2})}\right)
 =
 \frac1{(1-e^{i\theta_p}p^{-s})
          (1-e^{-i\theta_p}p^{-s})}.
\]

The local identity is an ordinary geometric-series calculation.
Under these unit-modulus parameters, the product is absolutely
convergent for \(\operatorname{Re}s>1\).
Real phases retain the Ramanujan hypothesis. General Satake
parameters require the corresponding convergence bound.

In the Hartnoll--Yang model the dilatation wavefunction contains
the completed \(L\)-function as a factor.
Its zeros are wavefunction zeros.
They are not, on that account, eigenvalues or a density of
states of the closed dilatation operator.
A fixed-prime spectral average, a fixed-state prime-distribution
statement and a large-prime continuum limit have separate
hypotheses. In particular, the formal multiplicity
\((p-1)/2\) is not an integer at \(p=2\).

\subsection{The evaluation map remains part of the comparison}

For a finite set of primes, an exterior generator of energy
\(\log p\) contributes \(1-p^{-s}\) to a supertrace.
The natural dual Ext generator instead has energy \(-\log p\)
and contributes \(1-p^s\).
Assigning it positive energy is an additional grading.
The restricted exterior limits have their different
absolute-convergence half-planes,
\(\operatorname{Re}s>1\) and \(\operatorname{Re}s<-1\).
They do not identify unrestricted Ext over an infinite polynomial
algebra with that restricted direct limit.

Likewise, a local Whittaker pairing has its full measure and
adjoint-\(L\) normalization.
A common nonzero complex rescaling preserves an equality of
sesquilinear pairings, but does not establish positivity.
A regularized global product is not an ordinary convergent
Euler product at one.
The arithmetic comparison must supply the actual maps and
their evaluation domains.

Satake estimates also have a precise scope.
For a unitary normalized \(GL_2\) parameter with extremal exponent
\(\theta_p\), the extremal symmetric-\(m\) exponent is \(m\theta_p\).
The general rank-four and rank-five estimates give
\(5/34\) through the symmetric cube and \(3/26\) through
the symmetric fourth. Its refined estimate gives \(7/64\).
Conditional cuspidal automorphy of the Maaß symmetric ninth
would give \(11/202\) by the general rank-ten bound.
Automorphy for an unbounded sequence of these symmetric powers
would force \(\theta_p=0\).

These are implications under the indicated automorphy hypotheses.
Newton--Thorne's symmetric-power theorem concerns non-CM regular
algebraic forms. A holomorphic companion does not multiply
a Maaß exponent by its own symmetric-power degree.
The identity
\[
 \Lambda^2\operatorname{Sym}^3(V)
 \simeq
 \bigl(\operatorname{Sym}^4(V)\otimes\det V\bigr)
       \oplus(\det V)^3
\]
retains both summands.

It does not supply missing cuspidality or temperedness.
In Ramakrishnan's tensor-product criterion, nondihedral
twist-related pairs are also noncuspidal.
A bound on local Satake exponents still requires an actual
comparison estimate before it gives a summable adelic error.

The arithmetic realization problem therefore retains the object,
its operators and their full comparison, followed by the scalar
evaluation. The finite coefficient, local representation and
formal gluing constructions above remain positive at those
levels. The automorphic trace, physical bulk and analytic
positivity comparisons require their separately stated maps.

\section*{References for the comparison theorems}

For determinants, see Chenevier,
\emph{The \(p\)-adic analytic space of pseudocharacters
of a profinite group and pseudorepresentations over arbitrary rings},
arXiv:0809.0415.
Kisin's \emph{The structure of potentially semi-stable deformation rings}
supplies the characteristic-zero lifting conditions.

The coefficient-prime categorical comparison is
Dotto--Emerton--Gee, \emph{A categorical \(p\)-adic Langlands
correspondence for \(GL_2(\mathbb Q_p)\)}, arXiv:2603.26887.
Its torsion and determinant conventions are specified above.
Pa\v{s}k\={u}nas, \emph{The image of Colmez's Montreal functor},
arXiv:1005.2008, supplies the locally finite block and Banach results.
Emerton--Gee, \emph{Moduli stacks of étale \((\varphi,\Gamma)\)-modules
and the existence of crystalline lifts}, arXiv:1908.07185,
supplies the formal-stack and cycle statements.

The affine-flag equivalences are Bezrukavnikov,
\emph{On two geometric realizations of an affine Hecke algebra},
arXiv:1209.0403.
The derived-loop construction is Ben-Zvi--Chen--Helm--Nadler,
\emph{Coherent Springer theory and categorical Deligne--Langlands},
arXiv:2010.02321.
Fargues--Scholze, \emph{Geometrization of the local Langlands correspondence},
arXiv:2102.13459, supplies the spectral action.
These results have the distinct coefficient theories stated in their respective comparisons.

The critical factorization theorem is Arinkin--Beraldo--Chen--F{\ae}rgeman--Gaitsgory--Lin--Raskin--Rozenblyum,
\emph{Proof of the geometric Langlands conjecture II:
Kac--Moody localization and the FLE}, arXiv:2405.03648.
The good-level comparison is Campbell--Dhillon--Raskin,
\emph{Fundamental local equivalences in quantum geometric Langlands},
arXiv:1907.03204.
The renormalized small equivalence is Gaitsgory--Lysenko,
\emph{Metaplectic Whittaker category and quantum groups:
the small FLE}, arXiv:1903.02279.

The perfect-stack center theorem is Ben-Zvi--Francis--Nadler,
\emph{Integral transforms and Drinfeld centers in derived algebraic geometry},
arXiv:0805.0157.
The finite Hecke character theorem is Ben-Zvi--Nadler,
\emph{The character theory of a complex group}, arXiv:0904.1247.
Singular support is developed by Arinkin--Gaitsgory,
\emph{Singular support of coherent sheaves and the geometric
Langlands conjecture}, arXiv:1201.6343.

The local irreducibility criterion is Bernstein--Zelevinsky,
\emph{Induced representations of reductive \(p\)-adic groups I}.
Flath's \emph{Decomposition of representations into tensor products}
gives the admissible irreducible factorization theorem.
The local cohomology input is Milne,
\emph{Arithmetic Duality Theorems}, second edition, Chapter I.

The categorical tensor and descent results use Lurie,
\emph{Higher Algebra}, with the locators given in their proofs.

The specified Bethe-algebra theorem is Mukhin--Tarasov--Varchenko,
\emph{Spaces of quasi-exponentials and representations of the Yangian
\(Y(\mathfrak{gl}_N)\)}, arXiv:1303.1578.
The charged oscillator model is Hartnoll--Yang,
\emph{The conformal primon gas at the end of time}, arXiv:2502.02661.
The protected chiral construction is Beem--Lemos--Liendo--Peelaers--Rastelli--van Rees,
\emph{Infinite chiral symmetry in four dimensions}, arXiv:1312.5344.

